% Leçon 25 — Vocabulaire et compréhension de texte : gestion du budget — Chinois Tle A — Cours signé OnBuch+
% Programme officiel MINESEC. Compiler avec : tectonic ZH25-vocabulaire-gestion-du-budget.tex
\newcommand{\DOCMATIERE}{Chinois}
\newcommand{\DOCNIVEAU}{Tle A}
\newcommand{\DOCMODULE}{Module 4 : Activités économiques et monde du travail}
\newcommand{\DOCLECON}{ZH25}
\newcommand{\DOCTITRE}{Leçon 25 — Vocabulaire et compréhension de texte : gestion du budget}
% OnBuch+ — préambule « cours signés » ENRICHI (compilé avec tectonic / XeLaTeX)
% Système visuel « Pop » d'OnBuch+ : fond crème (jamais blanc), bordures encre
% épaisses, ombres « dures » décalées, titres Archivo Black en capitales.
% Version MATHÉMATIQUES : graphes (pgfplots/tikz), boîtes pédagogiques
% (définition, propriété, méthode, attention, le savais-tu, exercice résolu,
% exercice, TP, bilan), page de garde et signature OnBuch+ ancrée en bas.
%
% Macros attendues dans main.tex AVANT \input{preamble} :
%   \DOCMATIERE  \DOCNIVEAU  \DOCMODULE  \DOCLECON (numéro)  \DOCTITRE
\documentclass[11pt]{article}

\usepackage{fontspec}
\usepackage[french]{babel} % français = langue principale ; le chinois via \zh{..} / textebox (xeCJK)
\usepackage{xeCJK}
\usepackage{pifont} % \ding{51} \ding{55} utilisés par les modèles
\usepackage[a4paper,margin=2cm,top=3.4cm,bottom=2.8cm,headheight=1.4cm,footskip=1.3cm]{geometry}
\usepackage{amsmath,amssymb,amsfonts}
\usepackage{mathtools} % \xmapsto, \coloneqq... souvent utilisés par les modèles
\usepackage{cancel} % simplifications barrées (bilans de forces)
% \abs{z} (module, valeur absolue) et \norm{u} (norme), utilisés par les modèles
\providecommand{\abs}{}\let\abs\relax\DeclarePairedDelimiter{\abs}{\lvert}{\rvert}
\providecommand{\norm}{}\let\norm\relax\DeclarePairedDelimiter{\norm}{\lVert}{\rVert}
\usepackage[table]{xcolor} % [table] : fournit aussi \rowcolors (alternance de lignes)
\usepackage{pagecolor}
\usepackage{tikz}
\usetikzlibrary{arrows.meta,positioning,calc,shapes.geometric,shapes.misc,shapes.arrows,decorations.pathmorphing,decorations.pathreplacing,decorations.markings,patterns,shadows,mindmap,trees,fit,backgrounds,matrix,intersections,angles,quotes}
\usepackage{pgfplots}
\usepackage[version=4]{mhchem} % équations nucléaires \ce{^{235}_{92}U}
\usepackage[european,siunitx]{circuitikz} % schémas électriques
\pgfplotsset{compat=1.18}
\usepgfplotslibrary{fillbetween,groupplots} % aires entre courbes (calcul intégral)
\usepackage{enumitem}
\usepackage{fancyhdr}
% [most] charge toutes les bibliothèques tcolorbox (listings, minted, raster,
% documentation...) et gonfle inutilement la mémoire TeX sur un document long
% (~300 boîtes) : on ne charge que ce qui est réellement utilisé.
\usepackage[skins,breakable]{tcolorbox}
\usepackage{graphicx}
\usepackage{colortbl}
\usepackage{booktabs}
\usepackage{tabularx}
\usepackage{diagbox} % case d'en-tête diagonale des tableaux à double entrée
% Colonnes extensibles centrée / gauche / droite (C, L, R), souvent utilisées par les modèles
\newcolumntype{C}{>{\centering\arraybackslash}X}
\newcolumntype{L}{>{\raggedright\arraybackslash}X}
\newcolumntype{R}{>{\raggedleft\arraybackslash}X}
\usepackage{array}
\usepackage{multicol}
\usepackage{multirow}
\usepackage{siunitx}
% parse-numbers=false : accepte aussi \SI{2,5 \times 10^{-3}}{...}
\sisetup{locale=FR,output-decimal-marker={,},group-minimum-digits=4,parse-numbers=false}
\DeclareSIUnit\ohm{\ensuremath{\symup{\Omega}}} % U+2126 absent de la police
\DeclareSIUnit\division{div} % réglages d'oscilloscope (ms/div, V/div)
\DeclareSIUnit\tr{tr} % tours (vitesse de rotation, ex. \SI{1500}{\tr\per\minute})
\DeclareSIUnit\molar{M} % concentration molaire (ex. \SI{3}{\molar}, \SI{1}{\micro\molar})
\DeclareSIUnit\an{an} % année (le modèle confond parfois avec \year, un compteur TeX) — ex. \SI{5}{\kilo\meter\per\an}
\DeclareSIUnit\FCFA{FCFA} % franc CFA (ex. \SI{500}{\FCFA\per\kilo\gram})
\DeclareSIUnit\degreeBrix{\textdegree Brix} % teneur en sucre d'un sirop/jus (réfractométrie)
\DeclareSIUnit\month{mois} % le modèle utilise parfois \month, un compteur TeX, comme unité de durée
\usepackage{qrcode}
% \degree : commande fréquemment utilisée par les modèles pour un angle ou une
% température en dehors de \SI{}{} (non fournie par les paquets chargés).
\providecommand{\degree}{^{\circ}}
% \qed (fin de démonstration), utilisé par les modèles sans amsthm
\providecommand{\qed}{\hfill$\square$}
% \barre : double barre des valeurs interdites dans un tableau de variations (tkz-tab)
\providecommand{\barre}{\ensuremath{\|}}
% environnement proof (amsthm non chargé) utilisé par les modèles
\ifdefined\proof\else\newenvironment{proof}[1][Démonstration]{\par\noindent\textit{#1.}\ }{\qed\par}\fi
\usepackage{lastpage}
\usepackage{hyperref}
\hypersetup{hidelinks}

% ============================================================
% Palette « Pop » OnBuch+ — mêmes hex que la classe P (plus_ui.dart)
% ============================================================
\definecolor{poporange}{HTML}{F59321}
\definecolor{popdark}{HTML}{E07A0C}
\definecolor{popcream}{HTML}{FFF6E8}
\definecolor{popink}{HTML}{1C1714}
\definecolor{poppurple}{HTML}{7A5AE0}
\definecolor{popgreen}{HTML}{2E9E6A}
\definecolor{poppink}{HTML}{E0457B}
\definecolor{popblue}{HTML}{2F7DE1}
\definecolor{popgold}{HTML}{C98A12}
\definecolor{popmuted}{HTML}{8A7F76}
\definecolor{popline}{HTML}{EADFD2}
\definecolor{popgray}{HTML}{8A7F76}
\colorlet{popgrey}{popgray}
% Alias tolérants (le modèle nomme parfois les couleurs différemment)
\colorlet{popwhite}{white}
\colorlet{popblack}{popink}
\colorlet{popred}{poppink}
\colorlet{popyellow}{popgold}
% Teintes claires (fonds de tableaux/encadrés) — le modèle nomme parfois
% les couleurs avec un suffixe L (ex. popblueL) sans les avoir définies.
\colorlet{poporangeL}{poporange!20}
\colorlet{popdarkL}{popdark!20}
\colorlet{poppurpleL}{poppurple!20}
\colorlet{popgreenL}{popgreen!20}
\colorlet{poppinkL}{poppink!20}
\colorlet{popblueL}{popblue!20}
\colorlet{popgoldL}{popgold!20}
\colorlet{popredL}{popred!20}
\colorlet{popgrayL}{popgray!20}
% Teintes claires pour fonds de tableaux / zones
\colorlet{popblueL}{popblue!10!white}
\colorlet{poporangeL}{poporange!14!white}
\colorlet{popgreenL}{popgreen!12!white}
\colorlet{poppurpleL}{poppurple!12!white}
\colorlet{poppinkL}{poppink!12!white}
\colorlet{popgoldL}{popgold!14!white}

\pagecolor{popcream}

% ============================================================
% Typographie
% ============================================================
\defaultfontfeatures{Ligatures=TeX}
\setmainfont{PlusJakartaSans-Medium.ttf}[Path=../../../fonts/,BoldFont=PlusJakartaSans-Bold.ttf]
% Symboles, API et emojis absents de la police principale : repli sur DejaVu Sans (caractères actifs)
\newfontface\symfont{DejaVuSans.ttf}[Path=../../../fonts/]
\makeatletter
\newcommand\symactive[1]{\catcode#1=13 \begingroup\lccode`\~=#1 \lowercase{\endgroup\def~}{{\symfont\char#1}}}
\newcommand\symblank[1]{\catcode#1=13 \begingroup\lccode`\~=#1 \lowercase{\endgroup\def~}{}}
\newcommand\symhyph[1]{\catcode#1=13 \begingroup\lccode`\~=#1 \lowercase{\endgroup\def~}{\mbox{-}}}
\newcount\symc
\newcommand\symrange[2]{\symc=#1 \@whilenum\symc<#2 \do{\expandafter\symactive\expandafter{\the\symc}\advance\symc by 1 }}
\newcommand\symblankrange[2]{\symc=#1 \@whilenum\symc<#2 \do{\expandafter\symblank\expandafter{\the\symc}\advance\symc by 1 }}
\makeatother
\symrange{"0250}{"02FF}\symactive{"2713}\symactive{"2714}\symactive{"2717}\symactive{"2718}\symactive{"2610}\symactive{"2611}\symactive{"2612}
\symrange{"2600}{"27BF}
\catcode"2705=13 \begingroup\lccode`\~="2705 \lowercase{\endgroup\def~}{{\symfont\char"2713}}\symblank{"26BD}\symblank{"26A1}
\symrange{"2460}{"257F}\symactive{"223C}\symactive{"2248}\symactive{"2260}
\symactive{"01F8}\symactive{"01F9}\symactive{"1E3E}\symactive{"1E3F}
\symhyph{"2011}\symhyph{"2E1A}\symblankrange{"1F300}{"1FAFF}\symblank{"FE0F}
% Caractères chinois : WenQuanYi Zen Hei (sous-ensemble GB2312 + pinyin), gras/italique simulés
\setCJKmainfont{WenQuanYiZenHei-subset.ttf}[Path=../../../fonts/,AutoFakeBold=3,AutoFakeSlant=0.2]
\xeCJKsetup{CJKspace=false}
% Pinyin : voyelles à caron (ǎ ǐ ǒ ǔ ǖ ǘ ǚ ǜ) absentes de la police principale -> caractères actifs
\newfontface\pyfont{WenQuanYiZenHei-subset.ttf}[Path=../../../fonts/]
\newcommand\pyactive[1]{\catcode#1=13 \begingroup\lccode`\~=#1 \lowercase{\endgroup\def~}{{\pyfont\char#1}}}
\pyactive{"01CD}\pyactive{"01CE}\pyactive{"01CF}\pyactive{"01D0}\pyactive{"01D1}\pyactive{"01D2}\pyactive{"01D3}\pyactive{"01D4}\pyactive{"01D5}\pyactive{"01D6}\pyactive{"01D7}\pyactive{"01D8}\pyactive{"01D9}\pyactive{"01DA}\pyactive{"01DB}\pyactive{"01DC}
% Maths en sans-serif (Fira Math) avec chiffres et lettres droites de Plus
% Jakarta, pour que les formules se fondent dans le texte.
\usepackage[math-style=ISO,bold-style=ISO]{unicode-math}
\setmathfont{FiraMath-Regular.otf}[Path=../../../fonts/]
\setmathfont{PlusJakartaSans-Medium.ttf}[Path=../../../fonts/,range={up/num,up/{Latin,latin},\mathup}]
\usepackage{icomma} % virgule décimale sans espace en mode math
% Caractères mathématiques Unicode absents des polices de texte (Plus Jakarta,
% Archivo Black) : rendus par leur équivalent mathématique, en texte comme en
% formule — sinon ils s'affichent en carré vide (ex. « Arithmétique dans ℤ »).
\usepackage{newunicodechar}
\newunicodechar{ℝ}{\ensuremath{\mathbb{R}}}
\newunicodechar{ℤ}{\ensuremath{\mathbb{Z}}}
\newunicodechar{ℂ}{\ensuremath{\mathbb{C}}}
\newunicodechar{ℕ}{\ensuremath{\mathbb{N}}}
\newunicodechar{ℚ}{\ensuremath{\mathbb{Q}}}
\newunicodechar{𝔻}{\ensuremath{\mathbb{D}}}
\newunicodechar{α}{\ensuremath{\alpha}}
\newunicodechar{β}{\ensuremath{\beta}}
\newunicodechar{γ}{\ensuremath{\gamma}}
\newunicodechar{δ}{\ensuremath{\delta}}
\newunicodechar{ε}{\ensuremath{\varepsilon}}
\newunicodechar{θ}{\ensuremath{\theta}}
\newunicodechar{λ}{\ensuremath{\lambda}}
\newunicodechar{μ}{\ensuremath{\mu}}
\newunicodechar{σ}{\ensuremath{\sigma}}
\newunicodechar{τ}{\ensuremath{\tau}}
\newunicodechar{φ}{\ensuremath{\varphi}}
\newunicodechar{ψ}{\ensuremath{\psi}}
\newunicodechar{ω}{\ensuremath{\omega}}
\newunicodechar{Σ}{\ensuremath{\Sigma}}
% majuscules produites par \MakeUppercase dans les titres (α-aminés -> Α)
\newunicodechar{Α}{\ensuremath{\alpha}}
\newunicodechar{Β}{\ensuremath{\beta}}
\newunicodechar{Γ}{\ensuremath{\Gamma}}
\newunicodechar{Δ}{\ensuremath{\Delta}}
\newunicodechar{Ε}{\ensuremath{\varepsilon}}
\newunicodechar{Θ}{\ensuremath{\Theta}}
\newunicodechar{Λ}{\ensuremath{\Lambda}}
\newunicodechar{Μ}{\ensuremath{\mu}}
\newunicodechar{Τ}{\ensuremath{\tau}}
\newunicodechar{Φ}{\ensuremath{\Phi}}
\newunicodechar{Ψ}{\ensuremath{\Psi}}
\newunicodechar{Ω}{\ensuremath{\Omega}}
\newunicodechar{↦}{\ensuremath{\mapsto}}
\newunicodechar{∈}{\ensuremath{\in}}
\newunicodechar{∉}{\ensuremath{\notin}}
\newunicodechar{∧}{\ensuremath{\wedge}}
\newunicodechar{⇔}{\ensuremath{\Leftrightarrow}}
\newunicodechar{⟺}{\ensuremath{\Longleftrightarrow}}
\newunicodechar{⇒}{\ensuremath{\Rightarrow}}
\newunicodechar{⟹}{\ensuremath{\Longrightarrow}}
\newunicodechar{∘}{\ensuremath{\circ}}
\newunicodechar{∪}{\ensuremath{\cup}}
\newunicodechar{∩}{\ensuremath{\cap}}
\newunicodechar{≡}{\ensuremath{\equiv}}
\newunicodechar{⊂}{\ensuremath{\subset}}
\newunicodechar{⊥}{\ensuremath{\perp}}
\newunicodechar{∃}{\ensuremath{\exists}}
\newunicodechar{∀}{\ensuremath{\forall}}
\newunicodechar{∛}{\ensuremath{\sqrt[3]{\,}}}
\newunicodechar{∜}{\ensuremath{\sqrt[4]{\,}}}
\newunicodechar{⃗}{\ensuremath{{}^{\to}}}
\newunicodechar{ⁿ}{\ifmmode^{n}\else\textsuperscript{\ensuremath{n}}\fi}
\newunicodechar{ˣ}{\ifmmode^{x}\else\textsuperscript{\ensuremath{x}}\fi}
\newunicodechar{ᵇ}{\ifmmode^{b}\else\textsuperscript{\ensuremath{b}}\fi}
\newunicodechar{ʳ}{\ifmmode^{r}\else\textsuperscript{\ensuremath{r}}\fi}
\newunicodechar{ᵘ}{\ifmmode^{u}\else\textsuperscript{\ensuremath{u}}\fi}
\newunicodechar{ʸ}{\ifmmode^{y}\else\textsuperscript{\ensuremath{y}}\fi}
\newunicodechar{ᵃ}{\ifmmode^{a}\else\textsuperscript{\ensuremath{a}}\fi}
\newunicodechar{ᵉ}{\ifmmode^{e}\else\textsuperscript{\ensuremath{e}}\fi}
\newunicodechar{ᵏ}{\ifmmode^{k}\else\textsuperscript{\ensuremath{k}}\fi}
\newunicodechar{ᵖ}{\ifmmode^{p}\else\textsuperscript{\ensuremath{p}}\fi}
\newunicodechar{ᴺ}{\ifmmode^{N}\else\textsuperscript{\ensuremath{N}}\fi}
\newunicodechar{ᶜ}{\ifmmode^{c}\else\textsuperscript{\ensuremath{c}}\fi}
\newunicodechar{ʰ}{\ifmmode^{h}\else\textsuperscript{\ensuremath{h}}\fi}
\newunicodechar{ᵠ}{\ifmmode^{q}\else\textsuperscript{\ensuremath{q}}\fi}
\newunicodechar{ₙ}{\ifmmode_{n}\else\textsubscript{\ensuremath{n}}\fi}
\newunicodechar{ₐ}{\ifmmode_{a}\else\textsubscript{\ensuremath{a}}\fi}
\newunicodechar{ᵢ}{\ifmmode_{i}\else\textsubscript{\ensuremath{i}}\fi}
\newunicodechar{ᵣ}{\ifmmode_{r}\else\textsubscript{\ensuremath{r}}\fi}
\newunicodechar{ₖ}{\ifmmode_{k}\else\textsubscript{\ensuremath{k}}\fi}
\newunicodechar{ᵦ}{\ifmmode_{\beta}\else\textsubscript{\ensuremath{\beta}}\fi}
\newunicodechar{ₓ}{\ifmmode_{x}\else\textsubscript{\ensuremath{x}}\fi}
\newfontfamily\popdisplay{ArchivoBlack-Regular.ttf}[Path=../../../fonts/]
\newfontfamily\poplogo{SpaceGrotesk-Bold.ttf}[Path=../../../fonts/]
\newfontfamily\popsemi{PlusJakartaSans-SemiBold.ttf}[Path=../../../fonts/]
\newfontfamily\popxbold{PlusJakartaSans-ExtraBold.ttf}[Path=../../../fonts/]
\newfontfamily\popxboldls{PlusJakartaSans-ExtraBold.ttf}[Path=../../../fonts/,LetterSpace=3.0]

\setlist[enumerate]{leftmargin=1.7em,itemsep=5pt,topsep=4pt}
\setlist[itemize]{leftmargin=1.7em,itemsep=4pt,topsep=4pt,label={\color{poporange}\textbullet}}
\setlength{\parindent}{0pt}
\setlength{\parskip}{5pt}
\renewcommand{\arraystretch}{1.25}

% pgfplots : style Pop par défaut
\pgfplotsset{
  every axis/.append style={axis line style={popink,line width=1pt},tick style={popink},
    grid style={popline,line width=0.5pt},label style={font=\small},tick label style={font=\footnotesize}},
  popcurve/.style={poporange,line width=1.8pt,no markers,smooth},
}
% Même style disponible dans un \draw TikZ simple (hors pgfplots)
\tikzset{popcurve/.style={poporange,line width=1.8pt}}
\tikzset{popgrid/.style={popline,very thin}}
\colorlet{popcurve}{poporange} % le nom du style est parfois utilisé comme couleur

% ============================================================
% Pill / squiggle
% ============================================================
\newcommand{\poppill}[3]{%
  \tikz[baseline=-0.55ex]\node[draw=popink,line width=1.2pt,rounded corners=2.6pt,%
    fill=#2,inner xsep=6.5pt,inner ysep=3pt]{%
    \popxboldls\fontsize{7}{7}\selectfont\color{#3}\MakeUppercase{#1}};%
}
\newcommand{\popsquiggle}[1]{%
  \tikz[baseline=-0.4ex]{
    \draw[#1,line width=2.6pt,line cap=round]
      plot [smooth] coordinates {(0,0.09) (0.34,0.01) (0.68,0.21) (1.02,0.01) (1.36,0.10)};
  }%
}
% Mot-clé surligné (marqueur orange sous le texte)
\newcommand{\cle}[1]{{\popxbold\color{popdark}#1}}

% ============================================================
% En-tête / pied
% ============================================================
\pagestyle{fancy}
\fancyhf{}
\renewcommand{\headrulewidth}{0pt}
\renewcommand{\footrulewidth}{0pt}
\lhead{\raisebox{-0.15ex}{\poplogo\fontsize{13}{13}\selectfont\color{popink}OnBuch\color{poporange}+}}
\rhead{\poppill{\DOCMATIERE\ \textbullet\ \DOCNIVEAU\ \textbullet\ Leçon \DOCLECON}{white}{popink}}
\lfoot{\popxbold\fontsize{7.6}{7.6}\selectfont\color{popmuted}\MakeUppercase{Cours signé OnBuch+}\ \textbullet\ {\popsemi\DOCTITRE}}
\rfoot{\tikz[baseline=-0.6ex]\node[rounded corners=8.5pt,fill=popink,minimum height=17pt,minimum width=17pt,inner xsep=5pt,inner sep=0]{\popxbold\fontsize{8}{8}\selectfont\color{popcream}\thepage\,/\,\pageref*{LastPage}};}

% ============================================================
% Titres
% ============================================================
\newcommand{\chaptitre}[2]{%
  \begin{tcolorbox}[enhanced,colback=poporange,colframe=popink,boxrule=2.6pt,arc=11pt,
      drop shadow=popink,left=15pt,right=15pt,top=12pt,bottom=11pt,before skip=3pt,after skip=18pt]
    \poppill{#2}{popink}{popcream}\par\vspace{9pt}
    {\raggedright\hyphenpenalty=10000\exhyphenpenalty=10000\popdisplay\fontsize{22}{24}\selectfont\color{popcream}\MakeUppercase{#1}\par}
  \end{tcolorbox}
}
% Section numérotée : pastille orange avec le numéro + titre Archivo
\newcounter{popsec}
\newcommand{\coursec}[1]{%
  \stepcounter{popsec}\setcounter{popsub}{0}%
  \par\vspace{16pt}\noindent
  \begin{minipage}{\linewidth}
  \tikz[baseline=-0.7ex]\node[draw=popink,line width=1.6pt,fill=poporange,rounded corners=4pt,minimum size=22pt,inner sep=0]{\popdisplay\fontsize{12}{12}\selectfont\color{popcream}\thepopsec};%
  \hspace{8pt}{\popdisplay\fontsize{15}{17}\selectfont\color{popink}\MakeUppercase{#1}}\par
  \vspace{4pt}\noindent{\color{poporange}\rule{36pt}{3.6pt}}
  \end{minipage}\par\nopagebreak\vspace{8pt}
}
\newcounter{popsub}
\newcommand{\courssub}[1]{%
  \stepcounter{popsub}%
  \par\vspace{9pt}\noindent{\popxbold\fontsize{11.5}{13}\selectfont\color{popink}{\color{poporange}\thepopsec.\thepopsub}\hspace{6pt}#1}\par\nopagebreak\vspace{4pt}
}

% ============================================================
% Boîtes pédagogiques — même recette : fond blanc, bordure épaisse colorée,
% ombre dure encre, badge plein. Titre optionnel [#1].
% ============================================================
\tcbset{popbase/.style={breakable,enhanced,colback=white,boxrule=2.3pt,arc=9pt,
  shadow={3pt}{-3pt}{0pt}{popink},left=14pt,right=14pt,top=11pt,bottom=11pt,before skip=11pt,after skip=15pt,
  fontupper=\popsemi\fontsize{10.6}{14.5}\selectfont\color{popink},
  fontlower=\popsemi\fontsize{10.6}{14.5}\selectfont\color{popink}}}
% Coche dessinée (le glyphe ✓ n'existe pas dans les polices du document)
\newcommand{\popcheck}[1]{\tikz[baseline=-0.5ex]\draw[#1,line width=1.6pt,line cap=round,line join=round](-0.13,0.0)--(-0.04,-0.09)--(0.13,0.1);}
\providecommand{\coche}{\popcheck{popgreen}}
\providecommand{\resultat}[1]{\par\smallskip\noindent\fcolorbox{popgreen}{popgreenL}{\parbox{\dimexpr\linewidth-2\fboxsep-2\fboxrule\relax}{\textbf{Résultat.} #1}}\par\smallskip}
\newcommand{\popboxhead}[3]{\poppill{#1}{#2}{white}\ifx\relax#3\relax\else\hspace{6pt}{\popxbold\fontsize{10.6}{12}\selectfont\color{popink}#3}\fi\par\vspace{7pt}}

\newtcolorbox{aretenir}[1][]{popbase,colframe=popgreen,before upper={\popboxhead{À retenir}{popgreen}{#1}}}
% Texte en chinois (document de lecture, dialogue, extrait) : cadre
% d'étude. Un titre optionnel (source, auteur) s'affiche dans l'en-tête.
\newtcolorbox{textebox}[1][]{popbase,colframe=popblue,before upper={\popboxhead{文本 · Texte}{popblue}{#1}},after upper={}}
% Marques ✓ / ✗ (forme correcte / fautive) souvent utilisées par les modèles
\providecommand{\cross}{\textcolor{poppink}{\ensuremath{\times}}}
\providecommand{\xmark}{\cross}\providecommand{\crossmark}{\cross}\providecommand{\cmark}{\ensuremath{\checkmark}}
% Ligne de réponse / justification à compléter (inventée par les modèles)
\providecommand{\justification}[1]{\par\noindent\textit{Justification~:} \dotfill\par}
% \textipa (package tipa non chargé) : la police ne couvre pas l'API -> simple passage
\providecommand{\textipa}[1]{#1}
% Soulignement (ulem non chargé) : \uline, \ul
\providecommand{\uline}[1]{\underline{#1}}\providecommand{\ul}[1]{\underline{#1}}
% \glossaire{\item ...} inventée par les modèles : liste de glossaire
\providecommand{\glossaire}[1]{\begin{itemize}#1\end{itemize}}
% \alert (beamer) inventée par les modèles : texte mis en évidence
\providecommand{\alert}[1]{\textbf{\textcolor{poppink}{#1}}}
% variantes ASCII d'ordinaux écrites en macro
\providecommand{\iere}{\textsuperscript{re}}\providecommand{\ieme}{\textsuperscript{e}}\providecommand{\ere}{\textsuperscript{re}}\providecommand{\eme}{\textsuperscript{e}}\providecommand{\er}{\textsuperscript{er}}\providecommand{\ie}{\textsuperscript{e}}
% Ordinaux français écrits en macro par les modèles : 1\ière, 3\ième
\newcommand{\ière}{\textsuperscript{re}}\newcommand{\ième}{\textsuperscript{e}}
% \exemple (inventée par les modèles) : étiquette « Exemple : »
\providecommand{\exemple}{\textbf{Exemple~:}\ }
% \square utilisable aussi hors mode math (cases à cocher)
\AtBeginDocument{\let\squareMath\square\renewcommand{\square}{\ensuremath{\squareMath}}}
% Mot ou phrase en chinois à l'intérieur d'un paragraphe français
\newcommand{\zh}[1]{\ifmmode\text{#1}\else{#1}\fi}
\let\eng\zh \let\esp\zh \let\all\zh % alias tolérés
% flèches d'intonation inventées par les modèles
\providecommand{\textsearrow}{}\renewcommand{\textsearrow}{\ensuremath{\searrow}}\providecommand{\textnearrow}{}\renewcommand{\textnearrow}{\ensuremath{\nearrow}}
% \fr{..} : mot français (inventé par les modèles)
\providecommand{\fr}[1]{\textit{#1}}
% zhbox : encadré de texte chinois inventé par les modèles
\newenvironment{zhbox}{\begin{quote}}{\end{quote}}
% \courssubsub : titre de niveau 3 inventé par les modèles
\providecommand{\courssubsub}[1]{\par\medskip\noindent\textbf{#1}\par\smallskip}
% \vs : « versus » inventé par les modèles
\providecommand{\vs}{\textit{vs}\ }
% \blank : case à compléter inventée par les modèles
\providecommand{\blank}[1][]{\underline{\hspace{1.6em}}}
\newtcolorbox{exemplebox}[1][]{popbase,colframe=poporange,before upper={\popboxhead{Exemple}{poporange}{#1}}}
\newtcolorbox{definition}[1][]{popbase,colframe=popblue,before upper={\popboxhead{Définition}{popblue}{#1}}}
\newtcolorbox{propriete}[1][]{popbase,colframe=poppurple,before upper={\popboxhead{Propriété}{poppurple}{#1}}}
\newtcolorbox{methode}[1][]{popbase,colframe=popgold,before upper={\popboxhead{Méthode}{popgold}{#1}}}
\newtcolorbox{attention}[1][]{popbase,colframe=poppink,before upper={\popboxhead{Attention}{poppink}{#1}}}
\newtcolorbox{experience}[1][]{popbase,colframe=popblue,colback=popblueL,before upper={\popboxhead{Expérience}{popblue}{#1}}}
% « Le savais-tu ? » : carte violette pleine (contexte, culture, Cameroun)
\newtcolorbox{savaistu}[1][]{popbase,colback=poppurple,colframe=popink,
  fontupper=\popsemi\fontsize{10.4}{14.2}\selectfont\color{white},
  before upper={\poppill{Le savais-tu\,?}{white}{poppurple}\ifx\relax#1\relax\else\hspace{6pt}{\popxbold\color{white}#1}\fi\par\vspace{7pt}}}
% Exercice résolu : énoncé en haut, \tcblower puis la solution sur fond crème
\newcounter{popexo}\newcounter{popexr}
\newtcolorbox{exoresolu}[1][]{popbase,colframe=poporange,colbacklower=popcream,
  segmentation style={popink,line width=1.2pt,dashed},
  before upper={\stepcounter{popexr}\popboxhead{Exercice résolu \thepopexr}{poporange}{#1}},
  before lower={\poppill{Solution}{popink}{popcream}\par\vspace{6pt}}}
% Exercice (non résolu en place) — la correction est dans la partie Corrigés
\newtcolorbox{exercice}[1][]{popbase,colframe=popink,
  before upper={\stepcounter{popexo}\popboxhead{Exercice \thepopexo}{popink}{#1}}}
% Corrigé
\newtcolorbox{corrige}[1][]{popbase,colframe=popgreen,colback=popgreenL,
  before upper={\popboxhead{Corrigé}{popgreen}{#1}}}
% Objectifs / prérequis / bilan
\newtcolorbox{objectifs}{popbase,colframe=popink,colback=poporangeL,
  before upper={\popboxhead{Objectifs de la leçon}{popink}{}}}
\newtcolorbox{prerequis}{popbase,colframe=popink,colback=popblueL,
  before upper={\popboxhead{Prérequis}{popblue}{}}}
\newtcolorbox{bilan}{popbase,colframe=popink,colback=white,boxrule=2.8pt,
  before upper={\popboxhead{Fiche bilan}{poporange}{L'essentiel à connaître pour l'examen}}}
% Figure encadrée avec légende
\newtcolorbox{popfigure}[1][]{enhanced,breakable=false,colback=white,colframe=popline,boxrule=1.4pt,arc=8pt,
  left=10pt,right=10pt,top=10pt,bottom=8pt,before skip=10pt,after skip=12pt,halign=center,
  lower separated=false,halign lower=center,
  fontlower=\popsemi\fontsize{9}{11.5}\selectfont\color{popmuted},
  #1}
\newcommand{\legende}[1]{\par\vspace{6pt}{\popsemi\fontsize{9}{11.5}\selectfont\color{popmuted}#1}}

% ============================================================
% Signature OnBuch+
% ============================================================
\newcommand{\poplogoinline}[1]{{\poplogo\fontsize{#1}{#1}\selectfont\color{popink}OnBuch\color{poporange}+}}
\newcommand{\signatureonbuch}{%
  \begin{tcolorbox}[enhanced,colback=popink,colframe=popink,boxrule=2.2pt,arc=10pt,
      drop shadow=poporange,left=14pt,right=14pt,top=10pt,bottom=10pt,before skip=0pt,after skip=0pt]
    \tikz[baseline=-0.7ex]\node[circle,fill=popgreen,draw=popcream,line width=1.3pt,minimum size=22pt,inner sep=0]{\popcheck{white}};%
    \hspace{9pt}\begin{minipage}[c]{0.62\linewidth}
      {\popxbold\fontsize{12}{13}\selectfont\color{popcream}Signé\ \textbullet\ {\poplogo OnBuch\color{poporange}+}}\par\vspace{2pt}
      {\raggedright\popsemi\fontsize{8}{10.5}\selectfont\color{popline}\MakeUppercase{Équipe pédagogique OnBuch+}\\\MakeUppercase{Conforme au programme officiel MINESEC}\par}
    \end{minipage}\hfill
    {\poplogo\fontsize{22}{22}\selectfont\color{popcream}OnBuch\color{poporange}+}
  \end{tcolorbox}%
}

% ============================================================
% Page de garde
% #1 titre · #2 matière · #3 niveau · #4 module · #5 n° leçon · #6 durée ·
% #7 description / objectifs (texte ou itemize)
% ============================================================
\newcommand{\pagegarde}[7]{%
  \thispagestyle{empty}
  \newgeometry{a4paper,margin=1.7cm,top=1.6cm,bottom=1.4cm}%
  \begin{tikzpicture}[remember picture,overlay]
    \fill[poporange!18] ([xshift=-3.2cm,yshift=-3.6cm]current page.north east) circle (4.6cm);
    \fill[poppurple!14] ([xshift=2.4cm,yshift=7.6cm]current page.south west) circle (3.8cm);
  \end{tikzpicture}%
  {\poplogo\fontsize{30}{30}\selectfont\color{popink}OnBuch\color{poporange}+}\hfill
  \raisebox{0.6ex}{\poppill{Cours signé \textbullet\ Premium}{popink}{popcream}}\par
  \vspace{1pt}\popsquiggle{poppurple}\par
  \vspace{8pt}
  \begin{tcolorbox}[enhanced,colback=poporange,colframe=popink,boxrule=2.8pt,arc=13pt,
      drop shadow=popink,left=18pt,right=18pt,top=12pt,bottom=13pt,after skip=10pt]
    \poppill{#2 \textbullet\ #3}{popink}{popcream}\hspace{5pt}\poppill{#4}{white}{popink}\par\vspace{8pt}
    {\popdisplay\fontsize{14}{14}\selectfont\color{popink}LEÇON #5}\par\vspace{4pt}
    {\raggedright\hyphenpenalty=10000\exhyphenpenalty=10000\popdisplay\fontsize{25}{27}\selectfont\color{popcream}\MakeUppercase{#1}\par}
    \vspace{8pt}
    {\popxbold\fontsize{9.5}{9.5}\selectfont\color{popink}\MakeUppercase{Durée conseillée : #6}}
  \end{tcolorbox}
  \begin{tcolorbox}[enhanced,colback=white,colframe=popink,boxrule=2.2pt,arc=10pt,
      drop shadow=popink,left=16pt,right=16pt,top=10pt,bottom=10pt,after skip=8pt]
    \poppill{Dans cette leçon}{poppurple}{white}\par\vspace{5pt}
    \popsemi\fontsize{9.3}{12.6}\selectfont\color{popink}#7
  \end{tcolorbox}
  \noindent\begin{minipage}[t]{0.487\linewidth}
    \begin{tcolorbox}[enhanced,colback=popgreen,colframe=popink,boxrule=2.2pt,arc=10pt,
        drop shadow=popink,left=11pt,right=11pt,top=10pt,bottom=10pt,height=3.1cm,valign=center]
      \tcbox[colback=white,colframe=popink,boxrule=1.3pt,arc=5pt,boxsep=0pt,nobeforeafter,
        left=3pt,right=3pt,top=3pt,bottom=3pt,valign=center]{\qrcode[height=1.9cm]{https://play.google.com/store/apps/details?id=com.luvvix.onbuch&hl=fr}}%
      \hspace{8pt}\begin{minipage}[c]{0.5\linewidth}
        {\popdisplay\fontsize{11}{12.5}\selectfont\color{white}\MakeUppercase{Télécharge OnBuch}}\par\vspace{4pt}
        {\raggedright\popsemi\fontsize{8.4}{11}\selectfont\color{white}Gratuit \textbullet\ Google Play\\Tuteur IA, annales, concours, résultats\par}
      \end{minipage}
    \end{tcolorbox}
  \end{minipage}\hfill
  \begin{minipage}[t]{0.487\linewidth}
    \begin{tcolorbox}[enhanced,colback=poppurple,colframe=popink,boxrule=2.2pt,arc=10pt,
        drop shadow=popink,left=11pt,right=11pt,top=10pt,bottom=10pt,height=3.1cm,valign=center]
      \tikz[baseline=-0.7ex]\node[circle,fill=white,draw=popink,line width=1.3pt,minimum size=1.6cm,inner sep=0]{\popdisplay\fontsize{24}{24}\selectfont\color{poppurple}+};%
      \hspace{8pt}\begin{minipage}[c]{0.6\linewidth}
        {\popdisplay\fontsize{11}{12.5}\selectfont\color{white}\MakeUppercase{Passe à OnBuch+}}\par\vspace{4pt}
        {\raggedright\popsemi\fontsize{8.4}{11}\selectfont\color{white}Tous les cours signés, Tuteur IA illimité, fiches PDF — onglet OnBuch+ de l'app\par}
      \end{minipage}
    \end{tcolorbox}
  \end{minipage}
  \vspace{8pt}
  \popcontenu
  \vfill
  \signatureonbuch
  \restoregeometry
  \newpage
}

% Rangée « Ce cours contient » (page de garde)
\newcommand{\poptile}[3]{% #1 couleur #2 gros texte #3 légende
  \begin{tcolorbox}[enhanced,colback=#1,colframe=popink,boxrule=2pt,arc=9pt,shadow={2.5pt}{-2.5pt}{0pt}{popink},
      left=6pt,right=6pt,top=9pt,bottom=9pt,halign=center,valign=center,height=1.75cm,width=\linewidth,nobeforeafter]
    {\popdisplay\fontsize{12.5}{13}\selectfont\color{white}\MakeUppercase{#2}}\par\vspace{3pt}
    {\centering\popsemi\fontsize{7.8}{9.5}\selectfont\color{white}#3\par}
  \end{tcolorbox}}
\newcommand{\popcontenu}{%
  \par\noindent{\popxboldls\fontsize{7.5}{7.5}\selectfont\color{popmuted}CE COURS SIGNÉ CONTIENT}\par\vspace{7pt}
  \noindent\begin{minipage}[t]{0.235\linewidth}\poptile{popblue}{Cours}{complet, illustré, pas à pas}\end{minipage}\hfill
  \begin{minipage}[t]{0.235\linewidth}\poptile{poporange}{Méthodes}{et exercices résolus}\end{minipage}\hfill
  \begin{minipage}[t]{0.235\linewidth}\poptile{poppink}{Exercices}{type examen + corrigés}\end{minipage}\hfill
  \begin{minipage}[t]{0.235\linewidth}\poptile{popgreen}{Fiche bilan}{l'essentiel à retenir}\end{minipage}\par}

% Sommaire
\newtcolorbox{sommaire}{enhanced,colback=white,colframe=popink,boxrule=2.2pt,arc=10pt,
  drop shadow=popink,left=18pt,right=18pt,top=13pt,bottom=13pt,before skip=6pt,after skip=20pt,
  before upper={\poppill{Sommaire}{poppurple}{white}\par\vspace{9pt}\popsemi\fontsize{10.6}{14.5}\selectfont\color{popink}}}

% Signature de fin de cours — poussée tout en bas de la dernière page
\newcommand{\findecours}{%
  \par\vfill\nopagebreak
  \begin{center}\popsquiggle{poporange}\end{center}\vspace{4pt}
  \signatureonbuch
}
\AtBeginDocument{\def\llbracket{\mathopen{[\mkern-3.5mu[}}\def\rrbracket{\mathclose{]\mkern-3.5mu]}}} % intervalles d'entiers (glyphe ⟦ absent de la police)
% Glyphes absents de Fira Math : redéfinis à partir de glyphes disponibles
\AtBeginDocument{%
  \def\setminus{\mathbin{\backslash}}\let\smallsetminus\setminus
  \def\vdots{\mathord{\vbox{\baselineskip3.2pt\lineskiplimit0pt\kern4pt\hbox{.}\hbox{.}\hbox{.}}}}%
  \def\ddots{\mathinner{\mkern1mu\raise6.5pt\hbox{.}\mkern2mu\raise3.6pt\hbox{.}\mkern2mu\raise0.7pt\hbox{.}\mkern1mu}}%
  \def\triangle{\mathord{\scalebox{0.9}{$\symup{\Delta}$}}}%
  \def\nabla{\mathord{\raisebox{\depth}{\scalebox{1}[-1]{$\symup{\Delta}$}}}}%
  \def\checkmark{\popcheck{popdark}}%
  \def\star{\mathbin{\ast}}\def\bigstar{\mathord{\ast}}%
  \def\diamond{\mathbin{\scalebox{0.6}{$\Diamond$}}}%
  \def\bigcup{\mathop{\mathchoice{\raisebox{-0.3em}{\scalebox{1.7}{$\cup$}}}{\scalebox{1.25}{$\cup$}}{\cup}{\cup}}\displaylimits}%
  \def\bigcap{\mathop{\mathchoice{\raisebox{-0.3em}{\scalebox{1.7}{$\cap$}}}{\scalebox{1.25}{$\cap$}}{\cap}{\cap}}\displaylimits}%
  \def\lesseqqgtr{\mathrel{\lessgtr}}\def\bigoplus{\mathop{\oplus}\displaylimits}%
}
\providecommand{\ssi}{si et seulement si} % abréviation fréquente
\AtBeginDocument{\providecommand{\wideparen}{\overparen}} % arc de cercle

%% FIGFIX-BEGIN v3 — correctifs globaux des figures (docs/onbuch-generation/tools/figfix)
\usepackage{adjustbox}\usepackage{etoolbox}
% toute figure TikZ/pgfplots plus large que la ligne est réduite à la largeur disponible
\BeforeBeginEnvironment{tikzpicture}{\begin{adjustbox}{max width=\linewidth}}
\AfterEndEnvironment{tikzpicture}{\end{adjustbox}}
% légendes pgfplots lisibles par-dessus les courbes
\pgfplotsset{every axis/.append style={legend style={fill=white,fill opacity=0.92,text opacity=1,draw=black!15,inner sep=3pt}}}
% étiquettes d'arêtes (syntaxe edge["..."]) avec fond blanc
\tikzset{every edge quotes/.append style={fill=white,inner sep=1.2pt}}
% bulles de cartes mentales : texte à la ligne dans la bulle, bulle un peu plus grande
\tikzset{every concept/.append style={text width=2.0cm,align=center,minimum size=2.3cm,inner sep=2pt}}
%% FIGFIX-END
\begin{document}

\pagegarde{Leçon 25 — Vocabulaire et compréhension de texte : gestion du budget}{Chinois}{Tle A}{Module 4 : Activités économiques et monde du travail}{ZH25}{2 h}{Cette leçon de vocabulaire et compréhension de texte porte sur la gestion du budget : à travers un texte original sur les finances d'une famille chinoise, les élèves acquièrent le lexique essentiel de l'argent (revenus, dépenses, épargne) et s'entraînent à lire, comprendre et réutiliser ce vocabulaire dans des situations concrètes.
\vspace{4pt}
\begin{multicols}{2}\small
\begin{itemize}
\item À la fin de cette leçon, je sais lire à voix haute un court texte chinois sur le budget avec une prononciation correcte (pinyin et tons).
\item À la fin de cette leçon, je sais reconnaître et traduire au moins 15 mots du thème de la gestion du budget.
\item À la fin de cette leçon, je sais identifier les clés (radicaux) et tracer l'ordre des traits des caractères 钱, 费, 省, 算.
\item À la fin de cette leçon, je sais répondre à des questions de compréhension portant sur un texte authentique simple.
\item À la fin de cette leçon, je sais employer les collocations essentielles : 花钱, 省钱, 存钱, 做预算.
\item À la fin de cette leçon, je sais comparer des prix et exprimer une dépense avec 花 + argent + 买 + objet.
\end{itemize}\end{multicols}}

\begin{sommaire}
\begin{multicols}{2}
\begin{enumerate}
\item Mise en route et découverte du thème
\item Texte support : lecture et écoute
\item Glossaire du thème
\item Clés et ordre des traits
\item Compréhension du texte et collocations
\item Production guidée et évaluation
\item Activité d'intégration
\item Méthodes et exercices résolus
\item Exercices
\item Corrigés des exercices
\item Fiche bilan
\end{enumerate}
\end{multicols}
\end{sommaire}

\begin{objectifs}
\begin{itemize}
\item À la fin de cette leçon, je sais lire à voix haute un court texte chinois sur le budget avec une prononciation correcte (pinyin et tons).
\item À la fin de cette leçon, je sais reconnaître et traduire au moins 15 mots du thème de la gestion du budget.
\item À la fin de cette leçon, je sais identifier les clés (radicaux) et tracer l'ordre des traits des caractères 钱, 费, 省, 算.
\item À la fin de cette leçon, je sais répondre à des questions de compréhension portant sur un texte authentique simple.
\item À la fin de cette leçon, je sais employer les collocations essentielles : 花钱, 省钱, 存钱, 做预算.
\item À la fin de cette leçon, je sais comparer des prix et exprimer une dépense avec 花 + argent + 买 + objet.
\item À la fin de cette leçon, je sais rédiger 4 à 6 phrases guidées présentant le budget de ma famille ou mon budget de poche.
\item À la fin de cette leçon, je sais distinguer les caractères proches 钱/线 et 费/佛.
\end{itemize}
\end{objectifs}

\begin{prerequis}
\begin{itemize}
\item Les nombres et les quantités en chinois (Première)
\item La monnaie et les prix : 块, 毛, 钱 (leçon sur les achats)
\item Les verbes courants 有, 是, 买, 用 et la négation 不/没
\item La structure de base sujet-verbe-complément et les compléments de quantité
\item Le vocabulaire de la famille et de la vie quotidienne (爸爸, 妈妈, 家)
\end{itemize}
\end{prerequis}

\newpage

\coursec{Mise en route et découverte du thème}

\courssub{Situation d'accroche : le salaire de fin de mois}

Imaginez la scène. Nous sommes à Yaoundé, quartier Mvog-Ada, un vendredi soir, fin du mois. La maman de Ngo Bell vient de recevoir son salaire. Avant même de rentrer à la maison, elle a déjà tout calculé dans sa tête : une partie pour les frais de scolarité des enfants, une partie pour la nourriture au marché, une enveloppe pour le transport en taxi et en moto, et enfin la cotisation de la tontine du quartier, qu'elle met soigneusement de côté. Rien n'est laissé au hasard : chaque franc a sa mission.

Savez-vous qu'en Chine, des millions de familles font exactement la même chose ? Elles aussi partagent leurs revenus entre les dépenses et l'épargne. Et elles ont un mot pour cela : \zh{预算} (\emph{yùsuàn}), le \cle{budget}. Toute la leçon d'aujourd'hui tourne autour de ce mot.

\begin{textebox}[Proverbe chinois pour ouvrir la réflexion]
钱不是万能的，但是没有钱是不行的。\\
\emph{Qián bú shì wànnéng de, dànshì méiyǒu qián shì bù xíng de.}\\
« L'argent ne fait pas tout, mais sans argent on ne peut rien faire. »
\end{textebox}

Ce proverbe moderne, très populaire en Chine, résume une sagesse que partagent les familles camerounaises et chinoises : l'argent n'est pas un but, mais un outil qu'il faut \emph{gérer}. C'est exactement le sens du mot \zh{预算}.

\begin{definition}[Le mot-clé de la leçon : 预算]
\zh{预算} (\emph{yùsuàn}) signifie « budget ». C'est un nom, mais il peut aussi fonctionner comme un verbe : « budgétiser, prévoir un budget ».
\begin{itemize}
\item \zh{预} (\emph{yù}) = « à l'avance, pré- » (comme dans \zh{预习} \emph{yùxí}, « préparer la leçon à l'avance ») ;
\item \zh{算} (\emph{suàn}) = « calculer » (comme dans \zh{算数} \emph{suànshù}, « l'arithmétique »).
\end{itemize}
Littéralement, \zh{预算} = « calculer à l'avance ». C'est exactement ce que fait la maman de Ngo Bell chaque fin de mois !
\end{definition}

\textbf{Question directrice de la leçon.} Écrivons au tableau la grande question qui guidera tout notre travail :

\begin{center}
\zh{你们家怎么做预算？}\\
\emph{Nǐmen jiā zěnme zuò yùsuàn ?}\\
« Comment votre famille fait-elle son budget ? »
\end{center}

Remarquez la structure : \zh{你们家} (\emph{nǐmen jiā}, « votre famille ») + \zh{怎么} (\emph{zěnme}, « comment ») + \zh{做预算} (\emph{zuò yùsuàn}, « faire un budget »). En chinois, on dit \zh{做预算} (« faire un budget ») et non « avoir un budget » : le verbe \zh{做} insiste sur l'action de planifier. Notez aussi que le mot interrogatif \zh{怎么} se place \emph{avant le verbe}, et non en tête de phrase comme le « comment » français.

\courssub{Brainstorming : et chez nous, comment gère-t-on l'argent ?}

Avant d'ouvrir le texte chinois, parlons de nous. Le professeur pose la question à la classe : \emph{comment votre famille gère-t-elle l'argent ?} Les élèves répondent d'abord en français, puis le professeur écrit au tableau les mots chinois correspondants. Voici les réponses typiques attendues :

\begin{itemize}
\item « Papa reçoit un salaire » → \zh{工资} (\emph{gōngzī}, le salaire) ;
\item « Maman vend des tomates au marché » → \zh{做生意} (\emph{zuò shēngyi}, faire du commerce) ;
\item « On paie les frais de scolarité » → \zh{学费} (\emph{xuéfèi}, les frais de scolarité) ;
\item « On met de l'argent dans la tontine » → \zh{存钱} (\emph{cún qián}, épargner de l'argent) ;
\item « On achète la nourriture » → \zh{买东西} (\emph{mǎi dōngxi}, acheter des choses).
\end{itemize}

Ce brainstorming a un double objectif : d'une part, il fait surgir du vécu des élèves le \emph{champ lexical} de l'argent ; d'autre part, il montre que les notions de revenus, de dépenses et d'épargne sont universelles — seules les langues changent.

\begin{popfigure}
\begin{tikzpicture}[mindmap, grow cyclic, every node/.style={concept, align=center},
root concept/.append style={concept color=popgold, font=\Large\bfseries},
level 1/.append style={level distance=3.6cm, sibling angle=90},
level 1 concept/.append style={concept color=popblueL, font=\small},
text=popdark]
\node[root concept, align=center]{钱\\ \emph{qián}\\ l'argent}
  child { node[align=center]{收入\\ \emph{shōurù}\\ revenus\\ 工资 salaire} }
  child { node[align=center]{支出\\ \emph{zhīchū}\\ dépenses\\ 学费 scolarité} }
  child { node[align=center]{储蓄\\ \emph{chǔxù}\\ épargne\\ 存钱 épargner} }
  child { node[align=center]{预算\\ \emph{yùsuàn}\\ budget\\ 做预算 budgétiser} };
\end{tikzpicture}
\legende{Carte mentale du champ lexical : autour de \zh{钱} (l'argent), quatre grandes familles de mots — les revenus, les dépenses, l'épargne et le budget.}
\end{popfigure}

Observez bien cette carte mentale : elle sera votre fil conducteur pendant toute la leçon. Chaque nouveau mot du texte support pourra être rangé dans l'une de ces quatre branches.

\textbf{Premiers mots du thème.} Voici le tableau des mots déjà rencontrés dans cette mise en route. Apprenez-les dès maintenant : ils reviendront dans le texte support.

\begin{center}
\begin{tabularx}{\linewidth}{|X|X|X|X|}
\hline
\rowcolor{popblueL} Caractères & Pinyin & Nature & Traduction \\
\hline
钱 & qián & nom & argent \\
\hline
收入 & shōurù & nom & revenus \\
\hline
支出 & zhīchū & nom & dépenses \\
\hline
储蓄 & chǔxù & nom & épargne \\
\hline
预算 & yùsuàn & nom / verbe & budget ; budgétiser \\
\hline
工资 & gōngzī & nom & salaire \\
\hline
学费 & xuéfèi & nom & frais de scolarité \\
\hline
存钱 & cún qián & verbe + compl. & épargner de l'argent \\
\hline
\end{tabularx}
\end{center}

\begin{attention}[Piège de prononciation : 钱 et ses homophones]
Le mot \zh{钱} (\emph{qián}, argent) porte le \textbf{deuxième ton} (ton montant). Ne le confondez pas avec \zh{前} (\emph{qián}, devant/avant), qui se prononce exactement de la même façon : ce sont des \emph{homophones}. En contexte, tout est clair : \zh{我没有钱。} (\emph{Wǒ méiyǒu qián.}, « Je n'ai pas d'argent. ») ne peut pas vouloir dire « je n'ai pas de devant ». Autre piège : dans \zh{收入} (\emph{shōurù}), la première syllabe est au premier ton (haut et plat) et la deuxième au quatrième ton (descendant) ; les francophones ont tendance à aplatir les deux syllabes. Entraînez-vous : \emph{shōu — rù}, en marquant bien la chute sur \emph{rù}.
\end{attention}

\courssub{Objectifs de la leçon}

À la fin de cette leçon de deux heures, vous serez capables de :

\begin{enumerate}
\item lire et comprendre un court texte chinois intitulé \zh{我们家的预算} (\emph{Wǒmen jiā de yùsuàn}, « Le budget de notre famille ») ;
\item reconnaître, prononcer avec les bons tons et employer une quinzaine de mots du thème de l'argent et du budget ;
\item répondre, en chinois, à des questions de compréhension sur le texte ;
\item comparer la gestion du budget en Chine et au Cameroun.
\end{enumerate}

Le texte support que nous lirons dans la section suivante raconte comment une famille chinoise organise son budget mensuel : le salaire des parents, les dépenses de nourriture et de scolarité, et l'argent mis de côté chaque mois. Vous verrez : il ressemble beaucoup à la vie de la famille de Ngo Bell !

\begin{savaistu}[Du téléphone à la tontine]
En Chine, beaucoup de familles n'utilisent presque plus d'argent liquide : elles paient et suivent leur budget mensuel avec des applications de téléphone comme Alipay (\zh{支付宝}, \emph{Zhīfùbǎo}) ou WeChat Pay. Chaque dépense est enregistrée automatiquement, et l'application peut afficher un résumé du mois : combien pour la nourriture, le transport, les études... Au Cameroun, les tontines se modernisent aussi : de nombreuses tontines fonctionnent désormais par Mobile Money, avec des cotisations envoyées par téléphone. Deux pays, deux technologies, une même idée : \emph{bien gérer son argent, c'est d'abord savoir où il va}.
\end{savaistu}

\courssub{Activation des prérequis : les nombres et les prix}

Pour comprendre un texte sur un budget, il faut être à l'aise avec les nombres et les prix. Révisons rapidement ce que vous avez appris en classe de Première.

\textbf{Les nombres.} Rappel : \zh{一} (1), \zh{十} (10), \zh{百} (100), \zh{千} (1000), \zh{万} (10 000). Attention, le chinois compte par \emph{dizaines de mille} : 50 000 se dit \zh{五万} (\emph{wǔ wàn}, « cinq dix-mille »), et non « cinquante mille ».

\textbf{Les prix : 块 et 毛.} La monnaie chinoise est le yuan. À l'oral, on dit :

\begin{propriete}[Dire un prix en chinois]
Structure : \textbf{nombre + 块 (\emph{kuài}) + (nombre + 毛 (\emph{máo}))}\\
\zh{块} est l'unité principale (le yuan, à l'oral) ; \zh{毛} est le dixième de yuan (comme nos « dix centimes »).\\
Exemples :
\begin{itemize}
\item \zh{五块} (\emph{wǔ kuài}) = 5 yuans ;
\item \zh{三块五} (\emph{sān kuài wǔ}) = 3,50 yuans (le \zh{毛} final est souvent omis) ;
\item \zh{这本书二十块。} (\emph{Zhè běn shū èrshí kuài.}) « Ce livre coûte vingt yuans. »
\end{itemize}
À l'écrit et dans les contextes formels, on utilise \zh{元} (\emph{yuán}) au lieu de \zh{块} : \zh{一百元} (\emph{yìbǎi yuán}, 100 yuans).
\end{propriete}

Petit entraînement oral : le professeur dicte des prix, les élèves les écrivent en chiffres : \zh{十五块} (15), \zh{八块五} (8,50), \zh{两百块} (200), \zh{三千块} (3000).

\begin{exoresolu}[Traduire des prix]
\textbf{Énoncé.} Traduisez en chinois : a) « Le repas coûte douze yuans. » b) « Le cahier coûte 3,50 yuans. » c) « Son salaire est de cinq mille yuans par mois. »
\tcblower
\textbf{Solution détaillée.}
\begin{itemize}
\item a) On utilise la structure \emph{sujet + prix}, où le prix fonctionne comme prédicat : \zh{这顿饭十二块。} (\emph{Zhè dùn fàn shí'èr kuài.}) « Ce repas (coûte) douze yuans. » Notez le classificateur \zh{顿} (\emph{dùn}) pour les repas. On peut aussi employer le verbe \zh{要} (« il faut, coûter ») : \zh{这顿饭要十二块。}
\item b) 3,50 = 3 \zh{块} 5 \zh{毛}, et on omet \zh{毛} à l'oral : \zh{这个本子三块五。} (\emph{Zhège běnzi sān kuài wǔ.}) « Ce cahier coûte 3,50 yuans. »
\item c) « Par mois » se dit \zh{每个月} (\emph{měi ge yuè}), placé avant le prédicat : \zh{他每个月的工资是五千块。} (\emph{Tā měi ge yuè de gōngzī shì wǔ qiān kuài.}) « Son salaire mensuel est de cinq mille yuans. »
\end{itemize}
\textbf{Piège à éviter :} ne dites pas \zh{*五千个块} : les prix n'ont pas besoin de classificateur, on dit directement \zh{五千块}. Le classificateur \zh{个} n'apparaît que devant les \emph{noms} (comme \zh{一个月}, « un mois »), jamais devant \zh{块}.
\end{exoresolu}

\courssub{Prédiction : deviner le texte avant de le lire}

Bonne nouvelle : pour comprendre un texte chinois, il ne faut pas attendre de connaître tous les mots. Les bons lecteurs \emph{anticipent}. Voici la méthode que nous utiliserons dans toutes les leçons de compréhension de texte.

\begin{methode}[Prédire le contenu d'un texte en trois étapes]
\begin{enumerate}
\item \textbf{Lire le titre.} Le titre annonce presque toujours le sujet. Ici : \zh{我们家的预算} (\emph{Wǒmen jiā de yùsuàn}) = « le budget de notre famille ». Je sais donc déjà que le texte parlera d'argent, dans une famille.
\item \textbf{Repérer les mots connus.} Dans le titre, je reconnais \zh{我们} (nous), \zh{家} (famille/maison) et le nouveau mot \zh{预算} (budget). La particule \zh{的} relie les deux : « le budget \emph{de} notre famille ».
\item \textbf{Formuler une hypothèse.} Avant de lire, je me dis : « Ce texte va probablement expliquer combien la famille gagne, ce qu'elle achète, et combien elle économise. » Je note deux ou trois mots que je m'attends à rencontrer : \zh{工资} (salaire), \zh{买} (acheter), \zh{存钱} (épargner).
\end{enumerate}
Pendant la lecture, je vérifierai si mes prédictions étaient justes : cela rend la lecture active et bien plus facile.
\end{methode}

\begin{experience}[À vous de prédire !]
Par groupes de trois, appliquez la méthode au titre \zh{我们家的预算} :
\begin{enumerate}
\item Notez trois mots chinois que vous pensez trouver dans le texte (aidez-vous de la carte mentale).
\item Formulez une phrase d'hypothèse en français : « Nous pensons que le texte va parler de... »
\item Chaque groupe présente sa prédiction à la classe. Le professeur les note au tableau : nous les confronterons au texte dans la section suivante.
\end{enumerate}
Durée : cinq minutes. L'objectif n'est pas d'avoir « juste », mais d'activer votre vocabulaire avant la lecture.
\end{experience}

\begin{aretenir}[Ce qu'il faut retenir de cette mise en route]
\begin{itemize}
\item Le mot-clé de la leçon est \zh{预算} (\emph{yùsuàn}), « budget », littéralement « calculer à l'avance » ; on « fait » un budget : \zh{做预算}.
\item Le champ lexical de l'argent s'organise en quatre branches : \zh{收入} (revenus), \zh{支出} (dépenses), \zh{储蓄} (épargne), \zh{预算} (budget).
\item Les prix se disent avec \zh{块} et \zh{毛} à l'oral (\zh{三块五} = 3,50 yuans), sans classificateur devant \zh{块}.
\item Avant de lire un texte, on prédit son contenu : titre, mots connus, hypothèse.
\end{itemize}
\end{aretenir}

Nous sommes maintenant prêts : le vocabulaire de base est activé, les prix sont révisés, et vos prédictions sont au tableau. Passons à la lecture du texte support \zh{我们家的预算}.

\coursec{Texte support : lecture et écoute}

Après la \emph{mise en route} où nous avons activé le lexique de l'argent et du budget à travers des images de la vie quotidienne à Yaoundé et à Douala (prix du marché, factures \textsc{ENEO}/\textsc{CAMWATER}, carnets d'épargne), nous entrons maintenant dans le cœur de la leçon : l'étude d'un texte authentique, court mais dense, qui met en scène une famille camerounaise gérant son budget mensuel. Ce texte, \zh{我们家的预算} (\emph{Wǒmen jiā de yùsuàn} — « Le budget de notre famille »), servira de fil conducteur pour le vocabulaire, la grammaire (\zh{先…再…}) et la compréhension.

\courssub{Première approche : écoute globale et compréhension générale}

L'enseignant lit le texte à vitesse normale, sans s'arrêter. Les élèves ne prennent pas de notes, ils laissent la mélodie de la phrase (l'intonation, le rythme, les tons) leur parvenir. L'objectif n'est pas de tout comprendre mot à mot, mais de saisir \textbf{l'information globale} : qui parle ? De quoi s'agit-il ? Quel est le ton général (sérieux, léger, inquiet, rassurant) ?

\begin{textebox}[Texte support : 我们家的预算]
\textbf{我们家的预算}

我们家的预算：我爸爸每个月工作，工资是五十万块。我们家先付房租和水电费，再买菜和吃的东西。妈妈每个月存一点钱，她说：“省钱很重要。”我也想做预算：我不乱花钱，把钱存起来买书。

\emph{Wǒmen jiā de yùsuàn: Wǒ bàba měi ge yuè gōngzuò, gōngzī shì wǔshí wàn kuài. Wǒmen jiā xiān fù fángzū hé shuǐdiànfèi, zài mǎi cài hé chī de dōngxi. Māma měi ge yuè cún yìdiǎn qián, tā shuō: “Shěng qián hěn zhòngyào.” Wǒ yě xiǎng zuò yùsuàn: wǒ bù luàn huā qián, bǎ qián cún qǐlái mǎi shū.}

\textbf{Traduction : Le budget de notre famille}

Le budget de notre famille : mon père travaille chaque mois, son salaire est de 500\,000 francs. Notre famille paie d'abord le loyer et les factures d'eau et d'électricité, puis achète les légumes et la nourriture. Maman économise un peu d'argent chaque mois, elle dit : « Économiser est très important. » Moi aussi je veux faire un budget : je ne dépense pas n'importe comment, je mets l'argent de côté pour acheter des livres.
\end{textebox}

\begin{aretenir}[Consignes pour l'écoute globale]
\begin{enumerate}
    \item \textbf{Fermez les yeux} pendant la première lecture pour vous concentrer sur les sons.
    \item Identifiez les \textbf{mots-clés} entendus (ex: \zh{爸爸}、\zh{工资}、\zh{房租}、\zh{存钱}、\zh{书}).
    \item Répondez oralement aux questions simples : \emph{Combien de personnes dans la famille ? Qui gagne l'argent ? Que fait la mère ? Que veut faire l'enfant ?}
\end{enumerate}
\end{aretenir}

\courssub{Analyse phrase par phrase : prononciation, tons et structure}

Nous passons maintenant à la \textbf{lecture analytique}. L'enseignant relit phrase par phrase (ou segment par segment). Les élèves répètent en chœur, en imitant exactement la courbe des tons. L'enseignant corrige les erreurs fréquentes (tons 2 vs 3, \zh{ü} vs \zh{u}, nasales \zh{an}/\zh{ang}, \zh{en}/\zh{eng}).

\begin{center}
\begin{tabularx}{\linewidth}{|X|X|X|X|}
\hline
\rowcolor{popblueL}
\textbf{Phrase chinoise} & \textbf{Pinyin (mots séparés)} & \textbf{Point de grammaire / Vocabulaire clé} & \textbf{Traduction} \\
\hline
我爸爸每个月工作 & Wǒ bàba měi ge yuè gōngzuò & \zh{每个月} (chaque mois) = fréquence ; \zh{工作} (verbe) & Mon père travaille chaque mois \\
\hline
工资是五十万块 & Gōngzī shì wǔshí wàn kuài & \zh{工资} (nom) salaire ; \zh{块} classif. oral pour Yuan/Franc & Le salaire est de 500\,000 \\
\hline
我们家先付房租和水电费 & Wǒmen jiā xiān fù fángzū hé shuǐdiànfèi & Structure \zh{先…} (d'abord) ; \zh{房租} loyer ; \zh{水电费} factures eau/élec. & Notre famille paie d'abord le loyer et l'eau/élec. \\
\hline
再买菜和吃的东西 & Zài mǎi cài hé chī de dōngxi & Structure \zh{再…} (ensuite) ; \zh{吃的东西} (nominalisé) nourriture & Puis achète légumes et nourriture \\
\hline
妈妈每个月存一点钱 & Māma měi ge yuè cún yìdiǎn qián & \zh{存} économiser/déposer ; \zh{一点} un peu (après verbe) & Maman économise un peu chaque mois \\
\hline
她说：“省钱很重要。” & Tā shuō: “Shěng qián hěn zhòngyào.” & \zh{省钱} (verbe composé) économiser ; \zh{重要} important & Elle dit : « Économiser est très important. » \\
\hline
我也想做预算 & Wǒ yě xiǎng zuò yùsuàn & \zh{做预算} faire un budget (collocation) & Moi aussi je veux faire un budget \\
\hline
我不乱花钱 & Wǒ bù luàn huā qián & \zh{乱} (adv) n'importe comment ; \zh{花钱} dépenser & Je ne dépense pas n'importe comment \\
\hline
把钱存起来买书 & Bǎ qián cún qǐlái mǎi shū & Structure \zh{把} (objet antéposé) ; \zh{存起来} mettre de côté & Je mets l'argent de côté pour acheter des livres \\
\hline
\end{tabularx}
\end{center}

\begin{attention}[Piège majeur : \zh{工资} (gōngzī) vs \zh{工作} (gōngzuò)]
Ces deux mots partagent le même premier caractère \zh{工} (travail, ouvrage, industriel) et une prononciation très proche (seul le ton du second caractère change : \zh{zī} ton 1 vs \zh{zuò} ton 4). C'est une source d'erreur \textbf{très fréquente} chez les francophones.

\begin{itemize}
    \item \zh{工作} (gōngzuò) = \textbf{Verbe} « travailler » ou \textbf{Nom} « travail, emploi, job ».
    \item \zh{工资} (gōngzī) = \textbf{Nom} « salaire, paye, rémunération ».
\end{itemize}

\textbf{Astuce mnémotechnique} : \zh{资} (zī) contient \zh{贝} (bèi), la clé « coquillage / argent / valeur » (présente dans \zh{买} acheter, \zh{卖} vendre, \zh{贵} cher, \zh{贫} pauvre, \zh{费} frais). \textbf{Là où il y a \zh{贝}, il y a de l'argent.} Donc \zh{工资} = l'argent du travail.
\end{attention}

\begin{methode}[Comment lire un texte chinois : la démarche en 3 temps]
\begin{enumerate}
    \item \textbf{Lecture globale (Skimming)} : Écouter/Lire vite pour le \emph{sens général} (Qui ? Quoi ? Où ? Quand ?). Ne pas s'arrêter aux mots inconnus.
    \item \textbf{Lecture analytique (Scanning)} : Phrase par phrase. \textbf{Repérer la structure} (Sujet - Verbe - Objet), \textbf{surligner les mots-outils} (\zh{了}、\zh{的}、\zh{把}、\zh{被}、\zh{先}、\zh{再}), \textbf{encercler le vocabulaire nouveau}, \textbf{vérifier les tons} sur le pinyin.
    \item \textbf{Relecture fluide (Reading aloud)} : Lire à voix haute \textbf{sans s'arrêter}, en respectant les groupes rythmiques (pauses) et l'intonation finale (descendante pour l'affirmation, montante pour la question). C'est ici que la prononciation se fixe en mémoire à long terme.
\end{enumerate}
\end{methode}

\courssub{Lecture à voix haute : fluidité et intonation}

Après l'analyse, place à la \textbf{production orale}.
\begin{enumerate}
    \item \textbf{Lecture en chœur} : Toute la classe lit ensemble derrière l'enseignant (modèle). Cela permet aux timides de s'entraîner sans exposition.
    \item \textbf{Lecture par groupes} : Un groupe lit les phrases du père, un autre celles de la mère, un autre celles de l'enfant.
    \item \textbf{Lecture individuelle} : 3 à 4 volontaires lisent le texte entier. L'enseignant note sur une grille : Tons (justesse), Liaison (fluidité), Intonation (naturelle), Volume/Regard.
\end{enumerate}

\begin{aretenir}[Rappel des groupes rythmiques (pauses) dans le texte]
En chinois, on ne fait pas de pause après chaque mot, mais après chaque \textbf{groupe sémantique} (sujet / verbe + objet / complément de phrase).
\zh{我爸爸 / 每个月 / 工作， // 工资 / 是 / 五十万块。 // 我们家 / 先 / 付房租和水电费， // 再 / 买菜和吃的东西。 // 妈妈 / 每个月 / 存 / 一点钱， // 她说 / ： / “省钱 / 很重要。” // 我也 / 想 / 做预算： // 我不 / 乱花钱， // 把钱 / 存起来 / 买书。}
\end{aretenir}

\courssub{Repérage lexical : les mots du glossaire dans le texte}

Les élèves surlignent dans le texte (ou sur leur fiche) les mots qui figurent dans le glossaire officiel du thème « Gestion du budget ». Voici le tableau de référence pour cette leçon.

\begin{center}
\begin{tabularx}{\linewidth}{|X|X|X|X|}
\hline
\rowcolor{popblueL}
\textbf{Caractères} & \textbf{Pinyin} & \textbf{Nature} & \textbf{Traduction} \\
\hline
预算 & yùsuàn & nom / verbe & budget ; établir un budget \\
\hline
工资 & gōngzī & nom & salaire, paye \\
\hline
发工资 & fā gōngzī & verbe (VO) & verser le salaire, toucher sa paye \\
\hline
房租 & fángzū & nom & loyer \\
\hline
水电费 & shuǐdiànfèi & nom & factures d'eau et d'électricité \\
\hline
存 & cún & verbe & économiser, déposer (banque), conserver \\
\hline
省钱 & shěng qián & verbe (VO) & économiser de l'argent (litt. « épargner argent ») \\
\hline
乱花钱 & luàn huā qián & verbe (VO) & dépenser n'importe comment, gaspiller \\
\hline
存起来 & cún qǐlái & verbe + compl. de résultat & mettre de côté, mettre de côté (résultat achevé) \\
\hline
先 & xiān & adverbe & d'abord, en premier \\
\hline
再 & zài & adverbe & ensuite, puis (action future/séquentielle) \\
\hline
然后 & ránhòu & adverbe & ensuite, puis (narration, plus formel) \\
\hline
最后 & zuìhòu & adverbe & enfin, en dernier lieu \\
\hline
块 & kuài & classif. & pièce, franc (oral), yuan (oral) \\
\hline
万 & wàn & nombre & dix mille (10\,000) \\
\hline
\end{tabularx}
\end{center}

\courssub{Structure grammaticale observée : \zh{先…，再…} (D'abord…, ensuite…)}

Le texte offre deux occurrences parfaites de cette structure séquentielle indispensable pour raconter une routine, donner un mode d'emploi ou décrire un processus logique.

\begin{exemplebox}[Exemples de \zh{先…，再…} dans et hors du texte]
\begin{itemize}
    \item \zh{我们家先付房租和水电费，再买菜和吃的东西。} (Texte) \\
    \emph{Wǒmen jiā xiān fù fángzū hé shuǐdiànfèi, zài mǎi cài hé chī de dōngxi.} \\
    Notre famille paie d'abord le loyer et l'eau/élec., puis achète la nourriture.
    \item \zh{我先写作业，再看电视。} \\
    \emph{Wǒ xiān xiě zuòyè, zài kàn diànshì.} \\
    Je fais d'abord mes devoirs, puis je regarde la télé.
    \item \zh{请先洗手，再吃饭。} \\
    \emph{Qǐng xiān xǐ shǒu, zài chī fàn.} \\
    Lavez-vous d'abord les mains, puis mangez.
    \item \zh{我们要先学拼音，再学汉字。} \\
    \emph{Wǒmen yào xiān xué pīnyīn, zài xué hànzì.} \\
    Nous devons d'abord apprendre le pinyin, puis les caractères.
\end{itemize}
\end{exemplebox}

\begin{propriete}[Règle : \zh{先…，再…} — Séquence chronologique stricte]
\textbf{Structure} : \zh{Sujet + 先 + Verbe 1 (+ Objet), 再 + Verbe 2 (+ Objet).}

\textbf{Emploi} : Indique que l'action 2 ne commence \textbf{qu'après la fin} de l'action 1. C'est une succession logique et temporelle. Le sujet est généralement le même pour les deux actions (ou un groupe collectif comme \zh{我们家}).

\textbf{Négation} : On nie généralement la seconde action ou l'ensemble : \zh{不要先吃，再洗} (Ne mange pas d'abord, puis lave [= ne fais pas dans cet ordre]). Pour dire « je ne fais pas A puis B », on préfère souvent \zh{不先…再…} ou changer l'ordre.

\textbf{Comparaison \zh{先…再…} vs \zh{先…然后…} (\emph{ránhòu})} :
\begin{itemize}
    \item \zh{再} (zài) : Plus courant à l'oral, actions proches, routine, instructions courtes.
    \item \zh{然后} (ránhòu) : Plus formel, narration, actions plus espacées ou complexes. « D'abord..., puis après..., ensuite... ».
\end{itemize}

\textbf{Attention francophones} : En français, on dit « Avant de faire B, je fais A » ou « Je fais A, puis B ». En chinois, l'ordre syntaxique \textbf{respecte toujours l'ordre chronologique réel} : \zh{先 A，再 B}. On ne dit pas *\zh{再 B，先 A}.
\end{propriete}

\begin{center}
\begin{tabularx}{\linewidth}{|X|X|X|}
\hline
\rowcolor{popblueL}
\textbf{Français} & \textbf{Chinois (Structure 先…再…)} & \textbf{Pinyin} \\
\hline
D'abord je mange, puis je pars. & 我先吃饭，再走。 & Wǒ xiān chī fàn, zài zǒu. \\
\hline
Il faut d'abord apprendre, puis comprendre. & 要先学，再懂。 & Yào xiān xué, zài dǒng. \\
\hline
Paie d'abord, mange ensuite. & 先付钱，再吃。 & Xiān fù qián, zài chī. \\
\hline
\end{tabularx}
\end{center}

\courssub{Schéma visuel : les priorités du budget familial}

La structure \zh{先…再…} du texte dessine une hiérarchie claire des dépenses. Ce schéma en escalier (TikZ) permet de visualiser l'ordre des priorités : les charges fixes incompressibles viennent avant la consommation courante, elle-même avant l'épargne et les loisirs.

\begin{popfigure}
\begin{tikzpicture}[
    node distance=0.8cm and 1.2cm,
    box/.style={rectangle, draw=popblue, thick, fill=popblueL, rounded corners=5pt, minimum width=3.2cm, minimum height=1.2cm, align=center, font=\small},
    arrow/.style={-Stealth, thick, popdark},
    labelbox/.style={rectangle, draw=poporange, thick, fill=poporangeL, rounded corners=3pt, minimum width=1.5cm, minimum height=0.8cm, align=center, font=\tiny, text=popdark}
]
% Nœuds principaux (escalier)
\node[box, align=center] (step1){\textbf{1. Charges fixes}\\\zh{房租 / 水电费}\\(Fángzū / Shuǐdiànfèi)\\\emph{Loyer / Eau-Élec.}};
\node[box, below right=of step1, align=center] (step2){\textbf{2. Vie courante}\\\zh{买菜 / 吃的东西}\\(Mǎi cài / Chī de dōngxi)\\\emph{Courses / Nourriture}};
\node[box, below right=of step2, align=center] (step3){\textbf{3. Épargne}\\\zh{存钱 / 省钱}\\(Cún qián / Shěng qián)\\\emph{Mettre de côté}};
\node[box, below right=of step3, align=center] (step4){\textbf{4. Projets / Loisirs}\\\zh{买书 / 其他}\\(Mǎi shū / Qítā)\\\emph{Livres / Autres}};

% Flèches
\draw[arrow] (step1) -- (step2) node[midway, above left, font=\tiny, text=poppurple] {\zh{先}};
\draw[arrow] (step2) -- (step3) node[midway, above left, font=\tiny, text=poppurple] {\zh{再}};
\draw[arrow] (step3) -- (step4) node[midway, above left, font=\tiny, text=poppurple] {\zh{然后}};

% Étiquettes "Priorité"
\node[labelbox, left=of step1, xshift=-0.5cm, align=center] (l1){Priorité\\MAXIMALE};
\node[labelbox, left=of step2, xshift=-0.5cm, align=center] (l2){Priorité\\FORTE};
\node[labelbox, left=of step3, xshift=-0.5cm, align=center] (l3){Priorité\\MOYENNE};
\node[labelbox, left=of step4, xshift=-0.5cm, align=center] (l4){Priorité\\LIBRE};
\end{tikzpicture}
\legende{Ordre des priorités budgétaires extrait du texte : \zh{先} (d'abord) charges fixes $\rightarrow$ \zh{再} (ensuite) vie courante $\rightarrow$ \zh{然后} (ensuite) épargne/projets.}
\end{popfigure}

\courssub{Clés (radicaux) et ordre des traits : zoom sur trois caractères du texte}

Pour ancrer l'écriture, nous analysons trois caractères clés du texte : \zh{预} (de \zh{预算}), \zh{存} et \zh{房} (de \zh{房租}).

\begin{center}
\begin{tabularx}{\linewidth}{|c|X|X|X|X|}
\hline
\rowcolor{popblueL}
\textbf{Car.} & \textbf{Clé (Radical) + Sens} & \textbf{Structure / Décomposition} & \textbf{Nb traits} & \textbf{Ordre des traits (règles clés)} \\
\hline
\zh{预} & \zh{页} (yè, page/tête) — en haut. \textbf{Sens} : à l'avance, prévoir. & Haut-Bas : \zh{页} + partie inférieure (forme simplifiée de \zh{予} avec traits additionnels). & 16 & 1. Écrire \zh{页} (haut→bas, 6 traits). 2. Écrire la partie basse (10 traits). Règle : \textbf{Haut avant bas}. \\
\hline
\zh{存} & \zh{子} (zǐ, enfant) — en bas. \textbf{Sens} : conserver, déposer. & Haut-Bas : \zh{才} (cái, talent) au-dessus de \zh{子}. & 6 & 1. \zh{才} (3 traits : \zh{丨} \zh{丿} \zh{丶}). 2. \zh{子} (3 traits : \zh{一} \zh{丨} \zh{乀}). Règle : \textbf{Haut avant bas}. \\
\hline
\zh{房} & \zh{户} (hù, porte/maison) — à gauche. \textbf{Sens} : maison, pièce. & Gauche-Droite : \zh{户} + \zh{方} (fāng, carré/côté). & 8 & 1. \zh{户} (4 traits : \zh{丶} \zh{一} \zh{ノ} \zh{一}). 2. \zh{方} (4 traits : \zh{丶} \zh{一} \zh{丨} \zh{一}). Règle : \textbf{Gauche avant droite}. \\
\hline
\end{tabularx}
\end{center}

\begin{savaistu}[Le caractère \zh{贝} (bèi) : la monnaie antique]
Vous l'avez croisé dans \zh{工资} (\zh{资}), \zh{买} (mǎi), \zh{卖} (mài), \zh{贵} (guì), \zh{贫} (pín), \zh{费} (fèi). \zh{贝} représente une \textbf{coquille de cauri} (\emph{Cypraea moneta}), utilisée comme monnaie d'échange dans la Chine antique (dynastie Shang, env. 1600-1046 av. J.-C.) et... dans plusieurs royaumes africains (dont le Cameroun actuel, région côtière et septentrionale) jusqu'au XIXe siècle ! C'est un pont culturel fascinant : la même « clé de l'argent » structure le vocabulaire financier en chinois comme les cauris structuraient l'économie traditionnelle camerounaise.
\end{savaistu}

\courssub{Exercice résolu : Compréhension fine et réemploi}

\begin{exoresolu}[Questions de compréhension sur le texte]
\textbf{Énoncé :} Répondez en français (ou en chinois simple) aux questions suivantes en vous basant \textbf{uniquement} sur le texte.

\begin{enumerate}
    \item Quel est le montant exact du salaire mensuel du père ? Quelle unité monétaire est utilisée (contexte camerounais) ?
    \item Citez dans l'ordre les quatre postes de dépense/épargne mentionnés pour la famille.
    \item Quelle est la citation exacte de la mère ? Quel proverbe ou conseil cela vous rappelle-t-il au Cameroun ?
    \item Quelle structure grammaticale exprime la chronologie des dépenses ? Donnez un autre exemple personnel.
    \item L'enfant dit : \zh{我不乱花钱}. Traduisez \zh{乱} ici. Quelle est la structure \zh{把} utilisée à la phrase suivante ? Décomposez-la.
\end{enumerate}

\tcblower

\textbf{Solution détaillée :}

\begin{enumerate}
    \item \textbf{500\,000 francs} (\zh{五十万块} — \emph{wǔshí wàn kuài}). \zh{万} = 10\,000. 50 $\times$ 10\,000 = 500\,000. \zh{块} (kuài) est le classifieur oral pour l'unité monétaire (Yuan en Chine, Franc CFA au Cameroun dans nos manuels).
    \item Ordre strict (\zh{先…再…}) : 1. \zh{房租} (Loyer) + \zh{水电费} (Eau/Électricité) — \textbf{Charges fixes}. 2. \zh{买菜和吃的东西} (Légumes + Nourriture) — \textbf{Alimentation}. 3. \zh{存一点钱} (Économiser un peu) — \textbf{Épargne mère}. 4. \zh{存起来买书} (Mettre de côté pour livres) — \textbf{Projet enfant}.
    \item \zh{“省钱很重要。”} (\emph{Shěng qián hěn zhòngyào.}) — « Économiser est très important. » Au Cameroun : « \emph{Grain par grain, la poule remplit son gésier} » ou « \emph{Petit à petit, l'oiseau fait son nid} » / « \emph{On ne bâtit pas une case en un jour} ». L'idée d'accumulation patiente (\zh{存}) est universelle.
    \item Structure : \zh{先…，再…} (\emph{xiān…, zài…}). Exemple personnel : \zh{我先复习中文，再玩手机。} (Je révise d'abord le chinois, puis je joue au téléphone).
    \item \zh{乱} (luàn) = « en désordre, n'importe comment, au hasard ». \zh{乱花钱} = « dépenser n'importe comment / gaspiller ».
    Structure \zh{把} : \zh{把 + 钱 (Objet) + 存起来 (Verbe + Complément de résultat)}.
    Schéma : \zh{Sujet (Wǒ) + 把 + Objet (qián) + Verbe (cún) + Compl. Résultat (qǐlái) + But (mǎi shū)}.
    Sens : « Je prends l'argent, je le mets-de-côté [résultat acquis], pour acheter des livres. » La structure \zh{把} met l'accent sur le \textbf{sort de l'objet} (l'argent est mis de côté).
\end{enumerate}
\end{exoresolu}

\courssub{Entraînement guidé : Reconstitution de phrases (ordre chronologique)}

Remettez les segments dans l'ordre logique en utilisant \zh{先} et \zh{再} (ou \zh{然后}). Écrivez les caractères et le pinyin.

\begin{exercice}[Mise en ordre chronologique]
\begin{enumerate}
    \item Segments : \zh{吃饭} (chī fàn) — \zh{洗手} (xǐ shǒu) — \zh{做饭} (zuò fàn).
    \item Segments : \zh{买票} (mǎi piào) — \zh{进站} (jìn zhàn) — \zh{上车} (shàng chē).
    \item Segments (Budget) : \zh{存钱} (cún qián) — \zh{发工资} (fā gōngzī) — \zh{付房租} (fù fángzū) — \zh{买菜} (mǎi cài).
\end{enumerate}
\end{exercice}

\begin{corrige}[Exercice : Mise en ordre chronologique]
\begin{enumerate}
    \item \zh{先洗手，再做饭，然后吃饭。} (\emph{Xiān xǐ shǒu, zài zuò fàn, ránhòu chī fàn.}) — Logique d'hygiène et de chronologie : on se lave les mains, on cuisine, on mange.
    \item \zh{先买票，再进站，然后上车。} (\emph{Xiān mǎi piào, zài jìn zhàn, ránhòu shàng chē.})
    \item \zh{先发工资，再付房租，然后买菜，最后存钱。} (\emph{Xiān fā gōngzī, zài fù fángzū, ránhòu mǎi cài, zuìhòu cún qián.}) — \zh{最后} (zuìhòu) « enfin / en dernier » s'ajoute naturellement pour clore la séquence.
\end{enumerate}
\end{corrige}

\begin{aretenir}[Synthèse de la section « Texte support »]
\begin{itemize}
    \item Le texte \zh{我们家的预算} est un modèle de \textbf{paragraphe narratif/descriptif} simple (HSK 2-3) structuré par \zh{先…再…}.
    \item Le vocabulaire clé (\zh{工资、房租、水电费、存、省钱、乱花钱、预算、发工资}) est \textbf{à haute fréquence} pour l'examen (écrit/oral).
    \item La prononciation des tons sur \zh{工资} (1-1), \zh{工作} (1-4), \zh{预算} (4-4), \zh{水电费} (3-4-4) doit être automatisée.
    \item La structure \zh{把钱存起来} introduit la construction \zh{把} (disposition) + complément de résultat \zh{起来}, point de grammaire majeur du Module 4 (repris en Leçon 26 « Grammaire »).
\end{itemize}
\end{aretenir}

\coursec{Mise en route et découverte du thème}

Bienvenue dans cette leçon consacrée à la \cle{gestion du budget} (\zh{预算管理}, \emph{yùsuàn guǎnlǐ}). Que vous soyez lycéen à Yaoundé, étudiant à Douala ou jeune actif à Buea, savoir gérer son argent — \zh{收入} (revenus), \zh{支出} (dépenses), \zh{储蓄} (épargne) — est une compétence vitale. En chinois, ce domaine lexical est d'une logique limpide : les mots se construisent comme des briques. Aujourd'hui, nous allons lire un texte authentique, maîtriser le vocabulaire clé, disséquer les caractères (clés et ordre des traits), répondre à des questions de compréhension et produire nos premières phrases budgétaires.

\coursec{Texte support : lecture et écoute}

Lisez le texte suivant à voix haute, en soignant les tons. Il met en scène un élève de Terminale, \zh{李明} (\emph{Lǐ Míng}), qui gère son argent de poche et observe le budget familial. Le texte est enregistré dans l'application OnBuch+ (piste audio ZH25-Texte).

\begin{textebox}[我的家庭预算 — Wǒ de jiātíng yùsuàn]
\zh{我叫李明，今年十八岁，是喀麦隆雅温得的一名高三学生。} \\
\zh{我爸爸每个月领工资八千块钱。} \\
\zh{我们家每个月要付房租三千块钱，水电费五百块钱。} \\
\zh{生活费大概两千块钱。} \\
\zh{爸爸常说：「做预算很重要，不要乱花钱，要存钱。」} \\
\zh{我听爸爸的话，每个月存五百块钱，省下来买书。} \\
\zh{我觉得管理预算很有意思，也很重要。}
\end{textebox}

\begin{center}
\begin{tabularx}{\linewidth}{|X|X|X|}
\hline
\rowcolor{popblueL} Caractères & Pinyin & Traduction \\
\hline
\zh{我叫李明，今年十八岁，是喀麦隆雅温得的一名高三学生。} & \emph{Wǒ jiào Lǐ Míng, jīnnián shíbā suì, shì Kāmàilóng Yǎwēndé de yī míng gāosān xuésheng.} & Je m'appelle Li Ming, j'ai 18 ans, je suis un élève de Terminale à Yaoundé, au Cameroun. \\
\hline
\zh{我爸爸每个月领工资八千块钱。} & \emph{Wǒ bàba měi ge yuè lǐng gōngzī bāqiān kuài qián.} & Mon père touche un salaire de 8000 yuans par mois. \\
\hline
\zh{我们家每个月要付房租三千块钱，水电费五百块钱。} & \emph{Wǒmen jiā měi ge yuè yào fù fángzū sānqiān kuài qián, shuǐdiànfèi wǔbǎi kuài qián.} & Notre famille doit payer 3000 yuans de loyer et 500 yuans de factures d'eau/électricité par mois. \\
\hline
\zh{生活费大概两千块钱。} & \emph{Shēnghuófèi dàgài liǎngqiān kuài qián.} & Les frais de vie s'élèvent à environ 2000 yuans. \\
\hline
\zh{爸爸常说：「做预算很重要，不要乱花钱，要存钱。」} & \emph{Bàba cháng shuō: « Zuò yùsuàn hěn zhòngyào, bú yào luàn huā qián, yào cún qián. »} & Papa dit souvent : « Faire un budget est important, ne dépense pas n'importe comment, il faut épargner. » \\
\hline
\zh{我听爸爸的话，每个月存五百块钱，省下来买书。} & \emph{Wǒ tīng bàba de huà, měi ge yuè cún wǔbǎi kuài qián, shěng xiàlái mǎi shū.} & J'écoute mon père, je mets 500 yuans de côté chaque mois, j'économise pour acheter des livres. \\
\hline
\zh{我觉得管理预算很有意思，也很重要。} & \emph{Wǒ juéde guǎnlǐ yùsuàn hěn yǒu yìsi, yě hěn zhòngyào.} & Je trouve que gérer un budget est intéressant et aussi très important. \\
\hline
\end{tabularx}
\end{center}

\begin{aretenir}[Repères culturels : l'argent au quotidien]
\begin{itemize}
\item \zh{块钱} (\emph{kuài qián}) : unité courante à l'oral pour \zh{元} (\emph{yuán}), comme « balles » ou « euros » en français familier. \zh{八千块钱} = 8000 yuans (environ 500 000 FCFA selon le taux 2024-2025).
\item \zh{高三} (\emph{gāosān}) : Terminale (3\up{e} année du lycée). \zh{高一}, \zh{高二}, \zh{高三}.
\item \zh{生活费} (\emph{shēnghuófèi}) : argent pour la vie courante (nourriture, transport, petites dépenses). Au Cameroun comme en Chine, c'est le poste principal après le loyer.
\end{itemize}
\end{aretenir}

\coursec{Glossaire du thème}

Dans la section précédente, nous avons lu et écouté le texte support \zh{我的家庭预算} (\emph{Wǒ de jiātíng yùsuàn}). Vous avez rencontré des mots comme \zh{工资}, \zh{房租} ou \zh{省钱}. Il est temps maintenant de nous arrêter sur chacun de ces mots : les observer, les prononcer correctement, comprendre leur construction et apprendre à les employer. Un bon glossaire n'est pas une liste à réciter : c'est une carte du territoire que nous allons explorer ensemble.

\courssub{Observation : les mots du budget dans leur contexte}

Avant tout tableau, observons quatre phrases tirées de situations réelles. Lisez-les à voix haute et repérez les mots qui parlent d'argent.

\begin{exemplebox}[Observer avant d'apprendre]
\zh{我爸爸每个月领工资。} \emph{Wǒ bàba měi ge yuè lǐng gōngzī.} « Mon père reçoit son salaire chaque mois. »\\
\zh{我们家要付房租和水电费。} \emph{Wǒmen jiā yào fù fángzū hé shuǐdiànfèi.} « Notre famille doit payer le loyer et les factures d'eau et d'électricité. »\\
\zh{这个东西太贵了！} \emph{Zhège dōngxi tài guì le !} « Cette chose est trop chère ! »\\
\zh{不要乱花钱，要存钱。} \emph{Bú yào luàn huā qián, yào cún qián.} « Ne dépense pas l'argent n'importe comment, il faut épargner. »
\end{exemplebox}

Que remarquez-vous ? Le mot \zh{钱} (\emph{qián}, argent) se combine avec des verbes pour former des expressions très fréquentes : \zh{花钱} (dépenser de l'argent), \zh{省钱} (économiser de l'argent), \zh{存钱} (mettre de l'argent de côté). Le chinois construit son vocabulaire par \cle{combinaison de caractères} : chaque caractère apporte sa brique de sens. C'est une excellente nouvelle pour vous : apprendre un caractère, c'est souvent débloquer plusieurs mots.

\courssub{Le mot central : 预算}

\begin{definition}[预算 yùsuàn — le budget]
\zh{预算} (\emph{yùsuàn}, nom) désigne le \cle{plan qui prévoit les revenus et les dépenses} d'une personne, d'une famille ou d'une organisation sur une période donnée (semaine, mois, année). Le caractère \zh{预} signifie « à l'avance » et \zh{算} signifie « calculer » : un budget, c'est littéralement « calculer à l'avance ».\\
Exemple : \zh{我们家每个月做预算。} \emph{Wǒmen jiā měi ge yuè zuò yùsuàn.} « Notre famille fait un budget chaque mois. »
\end{definition}

Notez le verbe qui accompagne le budget : on dit \zh{做预算} (\emph{zuò yùsuàn}, « faire un budget »), comme en français. En chinois, le choix du verbe qui accompagne un nom (la \cle{collocation}) est essentiel : on « fait » un budget, on « reçoit » un salaire (\zh{领工资}), on « paie » un loyer (\zh{付房租}).

\courssub{Tableau de vocabulaire du thème}

Voici le glossaire complet de la leçon. Lisez chaque ligne en trois temps : les caractères, le pinyin (à voix haute), puis la traduction.

\begin{center}
\begin{tabularx}{\linewidth}{|X|X|X|X|}
\hline
\rowcolor{popblueL} Caractères & Pinyin & Nature & Traduction \\
\hline
预算 & yùsuàn & n. & budget \\
\hline
钱 & qián & n. & argent \\
\hline
工资 & gōngzī & n. & salaire \\
\hline
收入 & shōurù & n. & revenu \\
\hline
支出 & zhīchū & n. & dépense (nom) \\
\hline
花 & huā & v. & dépenser \\
\hline
省 & shěng & v. & économiser \\
\hline
存 & cún & v. & déposer, épargner \\
\hline
付 & fù & v. & payer \\
\hline
房租 & fángzū & n. & loyer \\
\hline
水电费 & shuǐdiànfèi & n. & factures d'eau et d'électricité \\
\hline
费用 & fèiyòng & n. & frais \\
\hline
贵 & guì & adj. & cher \\
\hline
便宜 & piányi & adj. & bon marché \\
\hline
乱 & luàn & adv. & n'importe comment, en désordre \\
\hline
重要 & zhòngyào & adj. & important \\
\hline
\end{tabularx}
\end{center}

\begin{aretenir}[Les paires à retenir ensemble]
Le vocabulaire du budget fonctionne par couples logiques. Mémorisez-les deux par deux, c'est deux fois plus efficace :
\begin{itemize}
\item \zh{收入} (revenu) ↔ \zh{支出} (dépense) : remarquez les caractères \zh{入} (entrer) et \zh{出} (sortir) — l'argent qui entre, l'argent qui sort.
\item \zh{贵} (cher) ↔ \zh{便宜} (bon marché) : les deux adjectifs du marché de Mokolo ou de Douala.
\item \zh{花} (dépenser) ↔ \zh{省} / \zh{存} (économiser / épargner) : les deux gestes possibles face à l'argent.
\end{itemize}
\end{aretenir}

\courssub{Anatomie des mots composés}

Regardons de près comment ces mots sont construits. Comprendre la logique interne d'un mot chinois, c'est ne plus jamais l'oublier.

\begin{itemize}
\item \zh{工资} \emph{gōngzī} : \zh{工} (travail) + \zh{资} (ressources, capital) → « la ressource du travail », le salaire.
\item \zh{房租} \emph{fángzū} : \zh{房} (maison, pièce) + \zh{租} (louer) → « la location de la maison », le loyer.
\item \zh{水电费} \emph{shuǐdiànfèi} : \zh{水} (eau) + \zh{电} (électricité) + \zh{费} (frais) → « les frais d'eau et d'électricité ». Trois caractères, un seul mot, une logique limpide.
\item \zh{费用} \emph{fèiyòng} : \zh{费} (frais) + \zh{用} (utiliser) → « les frais d'utilisation », les frais en général.
\item \zh{收入} / \zh{支出} : \zh{收} (recevoir) + \zh{入} (entrer) ; \zh{支} (verser) + \zh{出} (sortir). Deux mots miroirs.
\end{itemize}

\begin{savaistu}[Le caractère 费 fèi, la machine à mots]
Le caractère \zh{费} (\emph{fèi}, frais, dépense) entre dans de très nombreux mots de la vie quotidienne : \zh{学费} (\emph{xuéfèi}, frais de scolarité), \zh{车费} (\emph{chēfèi}, frais de transport), \zh{生活费} (\emph{shēnghuófèi}, frais de vie), \zh{电话费} (\emph{diànhuàfèi}, facture de téléphone). Si vous connaissez \zh{费}, vous devinez le sens de tous ces mots au premier regard. C'est la force du chinois : le vocabulaire est \emph{cumulatif}. Au Cameroun comme en Chine, la question des \zh{学费} est d'ailleurs au cœur du budget des familles : en Chine, les parents commencent souvent à épargner pour les études des enfants dès leur naissance.
\end{savaistu}

\courssub{Travail de prononciation : les tons difficiles}

\begin{attention}[Pièges de prononciation pour les francophones]
\begin{itemize}
\item \zh{省} se prononce \emph{shěng}, \textbf{3e ton} : la voix descend puis remonte. Ne le confondez pas avec \zh{生} (\emph{shēng}, 1er ton, naître). Dites : \emph{shěng qián} (économiser de l'argent), en creusant bien la voix.
\item \zh{便宜} se prononce \emph{piányi} (2e ton + ton neutre) quand il signifie « bon marché ». \textbf{Ne dites jamais \emph{biànyí}} : cette prononciation existe mais appartient à un autre emploi (\zh{便宜} lu « biànyí » signifie « commode, avantageux », emploi rare et plutôt écrit). Au marché, c'est \emph{piányi}, toujours.
\item \zh{付} \emph{fù} : 4e ton, sec et descendant, comme un ordre bref. Ne le laissez pas traîner.
\item \zh{水电费} \emph{shuǐdiànfèi} : enchaînement 3-4-4. Le 3e ton de \zh{水} devient presque un demi-ton devant le 4e ton : dites-le bas, sans remonter.
\end{itemize}
\end{attention}

\begin{experience}[Gymnastique des tons en classe]
\begin{enumerate}
\item \textbf{Lecture en chœur} : le professeur lit chaque mot du tableau, la classe répète deux fois, d'abord lentement puis à vitesse normale.
\item \textbf{Minimal pairs} : le professeur alterne \emph{shěng} / \emph{shēng}, \emph{fù} / \emph{fǔ} ; les élèves lèvent la main gauche ou droite selon le ton entendu.
\item \textbf{Phrase cible} : répétez en chœur \zh{这个东西太贵了，便宜一点儿吧！} (\emph{Zhège dōngxi tài guì le, piányi yìdiǎnr ba !} « C'est trop cher, faites un prix ! ») — la phrase du marchandage par excellence.
\end{enumerate}
\end{experience}

\courssub{Classer pour mémoriser : revenus, dépenses, épargne}

Une liste de seize mots est difficile à retenir. Trois petites familles de mots, c'est facile. Classons notre glossaire selon la logique du budget : l'argent qui \cle{entre}, l'argent qui \cle{sort}, l'argent qu'on \cle{garde}.

\begin{popfigure}
\begin{tikzpicture}[
  box/.style={draw=popdark, rounded corners=4pt, thick, align=center, inner sep=6pt},
  head/.style={font=\bfseries, align=center}
]
\node[box, fill=popgreenL, minimum width=3.6cm, align=center] (rev) at (0,2){\textbf{收入 shōurù}\\[-2pt] \small les revenus};
\node[box, fill=popgreenL!60, minimum width=3.6cm, align=center] at (0,0.8){工资 gōngzī\\ salaire};
\node[box, fill=popgreenL!60, minimum width=3.6cm, align=center] at (0,-0.2){收入 shōurù\\ revenu};

\node[box, fill=poporangeL, minimum width=3.6cm, align=center] (dep) at (5,2){\textbf{支出 zhīchū}\\[-2pt] \small les dépenses};
\node[box, fill=poporangeL!60, minimum width=3.6cm, align=center] at (5,0.8){房租 fángzū / 水电费\\ loyer, factures};
\node[box, fill=poporangeL!60, minimum width=3.6cm, align=center] at (5,-0.2){花钱 huā qián\\ dépenser};

\node[box, fill=popblueL, minimum width=3.6cm, align=center] (epa) at (10,2){\textbf{储蓄 chǔxù}\\[-2pt] \small l'épargne};
\node[box, fill=popblueL!60, minimum width=3.6cm, align=center] at (10,0.8){存钱 cún qián\\ épargner};
\node[box, fill=popblueL!60, minimum width=3.6cm, align=center] at (10,-0.2){省钱 shěng qián\\ économiser};
\end{tikzpicture}
\legende{Les trois familles du vocabulaire du budget : ce qui entre, ce qui sort, ce qu'on garde.}
\end{popfigure}

\begin{methode}[Mémorisation guidée en trois étapes]
\begin{enumerate}
\item \textbf{Lecture en chœur} : lisez le tableau de vocabulaire colonne par colonne, deux fois, en exagérant les tons.
\item \textbf{Cache-cache du pinyin} : recouvrez la colonne pinyin avec une feuille. Regardez les caractères et prononcez. Vérifiez. Recommencez en cachant les caractères et en lisant seulement le pinyin : savez-vous encore dire la traduction ?
\item \textbf{Quiz rapide} : par binômes, un élève dit la traduction française (« le loyer »), l'autre répond en chinois avec le bon ton (\zh{房租}, \emph{fángzū}). Dix mots chacun, puis on inverse les rôles.
\end{enumerate}
\end{methode}

\courssub{Employer les mots : exemples et collocations}

Un mot n'est vraiment appris que s'il est employé dans une phrase. Voici les collocations essentielles du thème, c'est-à-dire les combinaisons verbe + nom que les Chinois utilisent naturellement.

\begin{center}
\begin{tabularx}{\linewidth}{|X|X|X|}
\hline
\rowcolor{popblueL} Collocation & Pinyin & Traduction \\
\hline
做预算 & zuò yùsuàn & faire un budget \\
\hline
领工资 & lǐng gōngzī & recevoir son salaire \\
\hline
付房租 & fù fángzū & payer le loyer \\
\hline
花钱 & huā qián & dépenser de l'argent \\
\hline
乱花钱 & luàn huā qián & dépenser n'importe comment \\
\hline
省钱 & shěng qián & économiser de l'argent \\
\hline
存钱 & cún qián & mettre de l'argent de côté \\
\hline
太贵了 & tài guì le & c'est trop cher \\
\hline
很便宜 & hěn piányi & c'est bon marché \\
\hline
\end{tabularx}
\end{center}

\begin{exemplebox}[Exemples traduits]
\zh{这个东西太贵了！} \emph{Zhège dōngxi tài guì le !} « Cette chose est trop chère ! »\\
\zh{这个月的水电费很高。} \emph{Zhège yuè de shuǐdiànfèi hěn gāo.} « Les factures d'eau et d'électricité sont élevées ce mois-ci. »\\
\zh{妈妈每个月存一点儿钱。} \emph{Māma měi ge yuè cún yìdiǎnr qián.} « Maman met un peu d'argent de côté chaque mois. »\\
\zh{管理预算很重要。} \emph{Guǎnlǐ yùsuàn hěn zhòngyào.} « Gérer son budget est très important. »
\end{exemplebox}

Remarquez la place de l'adverbe \zh{乱} (\emph{luàn}) : il se place \cle{avant le verbe} qu'il modifie — \zh{乱花钱} (dépenser n'importe comment), \zh{乱说} (dire n'importe quoi), \zh{乱买} (acheter n'importe quoi). En français, « n'importe comment » se met après le verbe ; en chinois, l'adverbe précède toujours.

\begin{attention}[Erreurs fréquentes des francophones]
\begin{itemize}
\item Ne dites pas \zh{付钱房租} : un seul verbe suffit. On dit \zh{付房租} (payer le loyer) ou \zh{付钱} (payer de l'argent), mais on n'empile pas les verbes.
\item \zh{花} s'emploie avec l'argent \emph{et} le temps : \zh{花钱} (dépenser de l'argent), \zh{花时间} (passer du temps). Ne traduisez pas « passer du temps » par \zh{过时间}.
\item Attention au sens de \zh{花} : c'est aussi un nom, « fleur » (\zh{一朵花}, une fleur). Le contexte tranche toujours.
\item \zh{贵} est un adjectif, donc on dit \zh{很贵} ou \zh{太贵了}, jamais \zh{是贵} : les adjectifs chinois ne prennent pas le verbe \zh{是}.
\end{itemize}
\end{attention}

\coursec{Clés (radicaux) et ordre des traits}

Écrire un caractère, c'est respecter un ordre précis : \cle{haut → bas}, \cle{gauche → droite}, \cle{horizontal avant vertical}, \cle{extérieur avant intérieur}, \cle{traits centraux avant les ailes}. Voici l'analyse de quatre caractères clés du thème.

\begin{definition}[钱 qián — argent]
\textbf{Radical (clé) :} \zh{钅} (\emph{jīn}, radical « or/métal », 5 traits), à gauche. Il indique la catégorie : l'argent, la monnaie, les métaux.\\
\textbf{Partie phonétique :} \zh{戈} (\emph{gē}, hallebarde), à droite.\\
\textbf{Nombre de traits :} 10 (simplifié).\\
\textbf{Ordre des traits :} 1.横 (héng) 2.竖 (shù) 3.提 (tí) 4.横 (héng) 5.竖 (shù) 6.横折 (héngzhé) 7.横 (héng) 8.撇 (piě) 9.横撇弯钩 (héngpiěwāngōu) — pour le radical \zh{钅} ; puis 10.横 (héng) 11.撇 (piě) 12.横折钩 (héngzhégōu) — pour \zh{戈}. \textit{Note : on écrit le radical gauche complètement avant la partie droite.}
\end{definition}

\begin{definition}[费 fèi — frais, dépense]
\textbf{Radical (clé) :} \zh{贝} (\emph{bèi}, radical « coquillage/monnaie », 4 traits), en bas. Les coquillages servaient de monnaie dans la Chine antique.\\
\textbf{Partie haute :} \zh{弗} (\emph{fú}), indicateur phonétique.\\
\textbf{Nombre de traits :} 9.\\
\textbf{Ordre des traits :} On écrit la partie haute \zh{弗} d'abord (haut → bas) : 1.横 2.竖 3.横 4.撇 5.竖弯钩 ; puis le radical \zh{贝} (extérieur → intérieur) : 6.竖 7.横折 8.横 9.竖弯钩.
\end{definition}

\begin{definition}[省 shěng — économiser ; province]
\textbf{Radical (clé) :} \zh{目} (\emph{mù}, radical « œil », 5 traits), en haut. L'idée originelle : « examiner de près » → « réduire, simplifier ».\\
\textbf{Partie basse :} \zh{少} (\emph{shǎo}, peu).\\
\textbf{Nombre de traits :} 9.\\
\textbf{Ordre des traits :} Haut → bas. D'abord \zh{目} : 1.竖 2.横折 3.横 4.横 5.横 ; puis \zh{少} : 6.竖 7.横 8.撇 9.捺 (ou point selon les normes).
\end{definition}

\begin{definition}[贵 guì — cher,贵重 ; honorifique]
\textbf{Radical (clé) :} \zh{贝} (\emph{bèi}, coquillage/monnaie), en bas. Quelque chose de « cher » a grande valeur monétaire.\\
\textbf{Partie haute :} \zh{无} (\emph{wú}, sans) + \zh{口} (\emph{kǒu}, bouche) + \zh{田} (\emph{tián}, champ) — structure complexe, à mémoriser par blocs.\\
\textbf{Nombre de traits :} 12.\\
\textbf{Ordre des traits :} Haut → bas. Bloc supérieur (\zh{无} + \zh{口} + \zh{田}) : 1.横 2.撇 3.横折钩 4.点 (\zh{无}) ; 5.竖 6.横折 7.横 (\zh{口}) ; 8.竖 9.横折 10.横 11.横 (\zh{田}) ; puis radical \zh{贝} : 12.竖 13.横折 14.横 15.竖弯钩. \textit{Total 12 traits en norme GB (le comptage varie selon les standards).}
\end{definition}

\begin{experience}[Atelier écriture : 10 minutes]
\begin{enumerate}
\item Sur papier quadrillé (grille 米字格), écrivez chaque caractère 5 fois en respectant l'ordre des traits.
\item Vérification croisée : votre binôme surveille l'ordre (pas de traits « sautés », pas de sens inverse).
\item Défi : écrivez \zh{水电费} et \zh{房租} en caractères, sans pinyin, de mémoire.
\end{enumerate}
\end{experience}

\coursec{Compréhension du texte et collocations}

\courssub{Questions de compréhension (texte : 我的家庭预算)}

Répondez en français (niveau 1) ou en chinois (niveau 2) d'après le texte support.

\begin{exercice}[Compréhension détaillée]
\begin{enumerate}
\item Quel est le salaire mensuel du père de Li Ming ? (Donnez le chiffre en yuans.)
\item Quelles sont les trois dépenses fixes mensuelles de la famille ? Citez-les avec leurs montants.
\item Quel conseil le père donne-t-il à Li Ming ? Citez les trois verbes d'action chinois correspondants.
\item Combien Li Ming met-il de côté chaque mois ? À quoi destine-t-il cette épargne ?
\item Trouvez dans le texte l'adjectif qui signifie « important » et l'adverbe qui signifie « souvent ».
\end{enumerate}
\end{exercice}

\begin{corrige}[Exercice « Compréhension détaillée »]
1. 8000 yuans (\zh{八千块钱}).\\
2. Loyer (\zh{房租}) : 3000 yuans ; Factures eau/électricité (\zh{水电费}) : 500 yuans ; Frais de vie (\zh{生活费}) : environ 2000 yuans.\\
3. \zh{做预算} (faire un budget), \zh{不要乱花钱} (ne pas dépenser n'importe comment), \zh{要存钱} (il faut épargner).\\
4. 500 yuans (\zh{五百块钱}) ; pour acheter des livres (\zh{买书}).\\
5. Adjectif : \zh{重要} (\emph{zhòngyào}) ; Adverbe : \zh{常} (\emph{cháng}, dans \zh{常说}).
\end{corrige}

\courssub{Collocations et structures utiles}

Le tableau ci-dessous résume les structures « verbe + objet » ou « adjectif + degré » indispensables pour l'expression écrite et orale sur ce thème.

\begin{center}
\begin{tabularx}{\linewidth}{|X|X|X|}
\hline
\rowcolor{popblueL} Structure & Exemple en caractères + pinyin & Traduction \\
\hline
Sujet + 领 + 工资 & 爸爸领工资。 \emph{Bàba lǐng gōngzī.} & Papa reçoit son salaire. \\
\hline
Sujet + 付 + 费用/房租 & 我们付房租。 \emph{Wǒmen fù fángzū.} & Nous payons le loyer. \\
\hline
Sujet + 花 + 钱/时间 & 别乱花钱。 \emph{Bié luàn huā qián.} & Ne dépense pas l'argent n'importe comment. \\
\hline
Sujet + 省 + 钱/力 & 我们要省钱。 \emph{Wǒmen yào shěng qián.} & Nous devons économiser. \\
\hline
Sujet + 存 + 钱 & 他存五百块。 \emph{Tā cún wǔbǎi kuài.} & Il met 500 yuans de côté. \\
\hline
Sujet + 很/太 + 贵/便宜 + 了 & 这太贵了！ \emph{Zhè tài guì le!} & C'est trop cher ! \\
\hline
做 + 预算 & 做预算很重要。 \emph{Zuò yùsuàn hěn zhòngyào.} & Faire un budget est important. \\
\hline
\end{tabularx}
\end{center}

\coursec{Production guidée et évaluation}

\courssub{Exercice résolu : classer, compléter, transformer}

\begin{exoresolu}[Du vocabulaire à la phrase]
Énoncé — (1) Classez : \zh{工资}, \zh{房租}, \zh{存钱}, \zh{收入}, \zh{水电费}, \zh{省钱}, \zh{花钱}, \zh{付房租} en 3 colonnes : Revenus / Dépenses / Épargne. (2) Complétez : \zh{这个月\_\_\_\_很高，我们要\_\_\_\_。} (Les factures sont hautes, il faut économiser.) (3) Transformez en phrase négative : \zh{这个东西很便宜。} → \zh{这个东西不\_\_\_\_。}
\tcblower
Solution — (1) Revenus : \zh{工资}, \zh{收入}. Dépenses : \zh{房租}, \zh{水电费}, \zh{花钱}, \zh{付房租}. Épargne : \zh{存钱}, \zh{省钱}. (2) \zh{这个月水电费很高，我们要省钱。} (3) \zh{这个东西不便宜。} (\emph{Zhège dōngxi bù piányi.}) — Négation d'adjectif : \zh{不} + adjectif (pas de \zh{很}).
\end{exoresolu}

\courssub{Exercices d'entraînement (devoirs / classe)}

\begin{exercice}[À vous de jouer — Niveau HSK 3-4]
\begin{enumerate}
\item \textbf{Traduction thème (FR → ZH)} : Écrivez en caractères + pinyin.
      \begin{enumerate}
      \item Le loyer de ce mois-ci est trop cher.
      \item Ne dépense pas ton argent n'importe comment.
      \item Faire un budget est très important pour les étudiants.
      \item Maman dépose 1000 yuans à la banque chaque mois. (Indice : \zh{存钱} ou \zh{存} + montant).
      \end{enumerate}
\item \textbf{Complétion lexicale} : Choisissez parmi \zh{花}, \zh{付}, \zh{省}, \zh{存}, \zh{领}, \zh{做}.
      \begin{enumerate}
      \item 爸爸每个月\_\_\_\_工资。
      \item 我们要\_\_\_\_预算。
      \item 房租太贵了，我想\_\_\_\_钱。
      \item 不要乱\_\_\_\_钱！
      \item 他每个月\_\_\_\_一千块钱买书。 (Il met de côté / Il dépense — les deux possibles, précisez le sens).
      \end{enumerate}
\item \textbf{Production écrite guidée} : Rédigez un court paragraphe (5-6 phrases) en chinois : « Mon budget de lycéen ». Utilisez au moins 8 mots du glossaire. Structure suggérée : Revenus (argent de poche, petit job) → Dépenses fixes (transport, cantine) → Épargne → Avis personnel (\zh{我觉得...}).
\item \textbf{Expression orale} : En binôme, jouez la scène : « Au marché de Mokolo / Marché central ». L'un est l'acheteur (trouve \zh{太贵了}, demande \zh{便宜一点儿}), l'autre est le vendeur (refuse, explique \zh{费用} élevés). 1 minute par rôle.
\end{enumerate}
\end{exercice}

\begin{corrige}[Exercice « À vous de jouer »]
1. a) \zh{这个月的房租太贵了。} (\emph{Zhège yuè de fángzū tài guì le.})\\
   b) \zh{不要乱花钱。} (\emph{Bú yào luàn huā qián.})\\
   c) \zh{做预算对学生很重要。} (\emph{Zuò yùsuàn duì xuésheng hěn zhòngyào.})\\
   d) \zh{妈妈每个月存一千块钱。} (\emph{Māma měi ge yuè cún yīqiān kuài qián.})\\
2. a) \zh{领} (\emph{lǐng}) — collocation \zh{领工资}.\\
   b) \zh{做} (\emph{zuò}) — collocation \zh{做预算}.\\
   c) \zh{省} (\emph{shěng}) — \zh{省钱} (réduire les dépenses) ou \zh{存} (\emph{cún}) — \zh{存钱} (mettre de côté). Les deux acceptés, sens nuancé.\\
   d) \zh{花} (\emph{huā}) — \zh{乱花钱}.\\
   e) \zh{存} (\emph{cún}) = il épargne 1000 ; \zh{花} (\emph{huā}) = il dépense 1000. Précisez à l'oral.\\
3. \textit{Exemple de production attendue :}\\
   \zh{我是高三学生。我每个月领零花钱两千块。我付车费五百块，花食堂费八百块。我存五百块钱，省钱买词典。我觉得做预算很重要，不乱花钱很好。}\\
   (\emph{Wǒ shì gāosān xuésheng. Wǒ měi ge yuè lǐng línghuāqián liǎngqiān kuài. Wǒ fù chēfèi wǔbǎi kuài, huā shítángfèi bābǎi kuài. Wǒ cún wǔbǎi kuài qián, shěng qián mǎi cídiǎn. Wǒ juéde zuò yùsuàn hěn zhòngyào, bù luàn huā qián hěn hǎo.})\\
4. Évaluation orale : prononciation des tons (3e ton \zh{省}, 4e ton \zh{付}, \zh{贵}), fluidité, respect des collocations (\zh{付房租} et non \zh{付钱房租}), utilisation de \zh{太...了} / \zh{不...}.
\end{corrige}

\begin{aretenir}[Bilan de la leçon ZH25 — Ce qu'il faut avoir acquis]
\begin{itemize}
\item \textbf{Vocabulaire} : 16 mots du thème « budget » (noms, verbes, adjectifs, adverbe) — reconnaissance, prononciation (tons), écriture (4 clés analysées).
\item \textbf{Grammaire / Structures} : Collocations verbo-nominales (\zh{领工资}, \zh{付房租}, \zh{存钱}, \zh{省钱}, \zh{做预算}) ; place de \zh{乱} (adv + V) ; adjectifs prédicatifs (\zh{很/太 + Adj + 了}) ; négation \zh{不} + Adj.
\item \textbf{Compréhension} : Lire un texte de 80-100 caractères sur la gestion budgétaire familiale, en extraire des informations chiffrées et des conseils.
\item \textbf{Production} : Écrire un paragraphe personnel (budget d'élève) ; jouer un dialogue de marchandage.
\item \textbf{Culture} : Unités monétaires (\zh{块钱} vs \zh{元}) ; réalités comparées loyer / frais d'études Cameroun-Chine ; caractère \zh{贝} comme trace historique de la monnaie-coquillage.
\end{itemize}
\end{aretenir}

\courssub{Prolongement : vers la leçon ZH26 (Grammaire)}

Dans la prochaine leçon, nous exploiterons ce vocabulaire pour maîtriser deux points grammaticaux clés du programme de Terminale : \textbf{la structure comparative avec \zh{比} (\emph{bǐ})} (\zh{房租比工资贵} — « Le loyer est plus cher que le salaire ») et \textbf{l'aspect accompli avec \zh{了} (\emph{le})} (\zh{我已经付了房租} — « J'ai déjà payé le loyer »). Préparez vos cartes de vocabulaire !

\coursec{Clés et ordre des traits}

Après avoir découvert le vocabulaire thématique dans le glossaire (section 3), nous allons maintenant « ouvrir le capot » des caractères pour comprendre leur architecture interne. En chinois, \cle{connaître la clé (radical) et l'ordre des traits} n'est pas un exercice calligraphique gratuit : c'est la condition \emph{sine qua non} pour retrouver un caractère dans un dictionnaire, pour l'écrire correctement (vitesse, lisibilité, proportion) et surtout pour en deviner le sens ou la prononciation. Les quatre caractères ci-dessous sont les « piliers » du thème « gestion du budget » ; les maîtriser, c'est s'ouvrir la porte de dizaines de mots composés.

\courssub{Analyse détaillée de quatre caractères fondamentaux}

Nous procédons pour chaque caractère selon la même méthode rigoureuse : \textbf{identification de la clé (sémantique) + identification de la partie phonétique (souvent) + décomposition en traits + ordre d'écriture normalisé}. Rappelons la règle d'or de l'ordre des traits : \textbf{haut $\to$ bas, gauche $\to$ droite, horizontal avant vertical, trait central avant les ailes, enveloppe avant contenu, fermeture en dernier}.

\begin{definition}[钱 \zh{qián} — argent, monnaie]
\textbf{Clé (radical) :} \zh{钅} (variante de \zh{金} \emph{jīn}, « métal/or »), placée à \textbf{gauche}. Elle indique la catégorie sémantique : le métal, donc la monnaie métallique ancienne.
\textbf{Partie droite (phonétique) :} \zh{戋} \emph{jiān} (petit, mince), qui donne l'indication phonétique (proximité sonore qián / jiān).
\textbf{Nombre de traits :} \textbf{10 traits} au total (5 pour \zh{钅} + 5 pour \zh{戋}).
\textbf{Ordre des traits (détail) :}
\begin{enumerate}
    \item \zh{钅} (clé gauche, 5 traits) : ① point (丶) en haut à gauche ; ② trait horizontal court (一) ; ③ trait horizontal court (一) ; ④ trait horizontal court (一) ; ⑤ trait relevé (㇀, \emph{tí}) — \emph{attention : ce n'est pas un crochet vers le bas (亅) ni vers la gauche, c'est un trait relevé qui « appelle » la partie droite}.
    \item \zh{戋} (partie droite, 5 traits) : ⑥ trait horizontal (一) en haut ; ⑦ trait vertical (丨) descendant ; ⑧ trait horizontal (一) ; ⑨ trait horizontal (一) ; ⑩ trait horizontal (一) en bas.
\end{enumerate}
\textbf{Structure :} \cle{Gauche-Droite} (左右结构). La clé est fine, la partie droite est large ; on centre l'ensemble dans le carré imaginaire.
\end{definition}

\begin{definition}[费 \zh{fèi} — frais, dépense, coûter, gaspiller]
\textbf{Clé (radical) :} \zh{贝} \emph{bèi} (coquillage), placée en \textbf{bas}. Rappel historique : les coquillages (cauris) servaient de monnaie dans la Chine antique (dynastie Shang). Toute clé \zh{贝} signale un lien avec l'argent, la valeur, le commerce.
\textbf{Partie haute (phonétique) :} \zh{弗} \emph{fú} (non, ne pas), qui donne la prononciation \emph{fèi} (proximité fú / fèi).
\textbf{Nombre de traits :} \textbf{9 traits} (comptage officiel simplifié pour l'examen : 5 pour \zh{弗} + 4 pour la partie basse \zh{贝} dans cette composition). Note : \zh{贝} isolé compte 7 traits, mais dans \zh{费}, le trait horizontal final de \zh{弗} fusionne avec le trait horizontal initial de \zh{贝}.
\textbf{Ordre des traits (détail) :}
\begin{enumerate}
    \item \zh{弗} (haut, 5 traits) : ① horizontal (一) ; ② vertical (丨) ; ③ crochet horizontal (㇆) — \emph{un seul trait : horizontal + crochet vers le bas} ; ④ point (丶) à gauche ; ⑤ point (丶) à droite.
    \item \zh{贝} (bas, 4 traits restants) : ⑥ trait horizontal (一) — \emph{base du 弗 et sommet du 贝} ; ⑦ trait vertical (丨) descendant ; ⑧ point (丶) à gauche ; ⑨ point (丶) à droite.
\end{enumerate}
\textbf{Structure :} \cle{Haut-Bas} (上下结构). Le haut (\zh{弗}) est plus large que le bas (\zh{贝}) ; on aligne l'axe central.
\end{definition}

\begin{definition}[省 \zh{shěng} — économiser ; province]
\textbf{Clé (radical) :} \zh{目} \emph{mù} (œil), placée en \textbf{bas}. Pourquoi l'œil pour « économiser » ? Sens originel : « regarder attentivement, inspecter » $\to$ « examiner ses dépenses » $\to$ « réduire, économiser ».
\textbf{Partie haute :} \zh{少} \emph{shǎo} (peu, jeune), indicateur phonétique (shǎo $\to$ shěng).
\textbf{Nombre de traits :} \textbf{9 traits} (4 pour \zh{少} + 5 pour \zh{目}).
\textbf{Ordre des traits (détail) :}
\begin{enumerate}
    \item \zh{少} (haut, 4 traits) : ① point (丶) ; ② trait oblique (丿) ; ③ point (丶) ; ④ trait vertical (丨) descendant — \emph{ne pas faire de crochet en bas, le trait s'arrête net pour « poser » sur le 目}.
    \item \zh{目} (bas, 5 traits) : ⑤ horizontal (一) — \emph{le toit de l'œil} ; ⑥ vertical (丨) descendant au milieu ; ⑦ horizontal (一) au milieu ; ⑧ horizontal (一) en bas ; ⑨ horizontal (一) de fermeture — \emph{le 目 est un rectangle divisé en deux par un trait horizontal médian, puis fermé par le bas}.
\end{enumerate}
\textbf{Structure :} \cle{Haut-Bas} (上下结构). Le haut \zh{少} dépasse légèrement sur les côtés du \zh{目}.
\end{definition}

\begin{definition}[算 \zh{suàn} — calculer, compter, planifier, estimer]
\textbf{Clé (radical) :} \zh{竹} \emph{zhú} (bambou), écrite \zh{⺮} (clé « toit de bambou ») en \textbf{haut}. Sens : les anciens calculaient avec des baguettes de bambou (算筹 \zh{suànchóu}).
\textbf{Partie basse (complexe) :} \zh{目} (œil) + \zh{廾} \emph{gǒng} (deux mains jointes) $\to$ « l'œil qui regarde les mains compter les baguettes ».
\textbf{Nombre de traits :} \textbf{14 traits} (6 pour \zh{⺮} + 5 pour \zh{目} + 3 pour \zh{廾}).
\textbf{Ordre des traits (détail) :}
\begin{enumerate}
    \item \zh{⺮} (clé bambou, haut, 6 traits) : ① horizontal (一) ; ② oblique (丿) ; ③ horizontal (一) ; ④ oblique (丿) ; ⑤ horizontal (一) ; ⑥ point (丶) — \emph{les six traits du « toit de bambou » s'écrivent comme le caractère 竹 mais compressés en largeur}.
    \item \zh{目} (milieu, 5 traits) : ⑦ horizontal (一) ; ⑧ vertical (丨) ; ⑨ horizontal (一) ; ⑩ horizontal (一) ; (11) horizontal (一) — \emph{même ordre que pour 省}.
    \item \zh{廾} (bas, 3 traits) : (12) horizontal (一) — \emph{bras gauche} ; (13) horizontal (一) — \emph{bras droit} ; (14) vertical (丨) descendant au centre reliant les deux — \emph{attention : l'ordre est horizontal gauche, horizontal droit, puis vertical central qui « soude » le tout}.
\end{enumerate}
\textbf{Structure :} \cle{Haut-Milieu-Bas} (上中下结构). Le \zh{⺮} est plat et large ; le \zh{目} est centré ; le \zh{廾} écarte ses deux horizontaux sous le \zh{目}.
\end{definition}

\begin{propriete}[La clé \zh{贝} (bèi) — le coquillage, racine de l'économie]
Dans l'écriture chinoise, la clé \zh{贝} (coquillage cauri) est le \cle{marqueur sémantique universel de l'argent, de la valeur, du commerce et de la richesse}. Les coquillages étaient une monnaie d'échange courante sous les Shang (商朝, env. 1600-1046 av. J.-C.). Tous les caractères ci-dessous contiennent \zh{贝} et relèvent du champ lexical « argent / valeur / échange » :
\begin{center}
\begin{tabularx}{\linewidth}{|X|X|X|X|}
\hline
\rowcolor{popblueL} \textbf{Caractère} & \textbf{Pinyin} & \textbf{Sens principal} & \textbf{Lien avec \zh{贝}} \\
\hline
费 & fèi & frais, dépenser & Argent qui sort (弗 = non/nier $\to$ perte) \\
\hline
贵 & guì & cher,贵重, noble & Haute valeur (古 = ancien/élevé $\to$ prix haut) \\
\hline
购 & gòu & acheter & Échanger (phon. 戔) contre argent \\
\hline
财 & cái & richesse, biens & Talent (才) + argent = richesse \\
\hline
赚 & zhuàn & gagner de l'argent & Difficile (兼) + argent = gain \\
\hline
赔 & péi & indemniser, perdre & Non (非) + argent = perte/remboursement \\
\hline
贷 & dài & prêter, emprunt & Substitut (代) + argent = prêt \\
\hline
\end{tabularx}
\end{center}
\textbf{Astuce de mémorisation :} Quand vous voyez \zh{贝}, pensez immédiatement « \emph{argent, valeur, transaction} ». Cela permet de deviner le sens d'un caractère inconnu (ex: \zh{赊} \emph{shē} = acheter à crédit, \zh{贿} \emph{huì} = pot-de-vin).
\end{propriete}

\begin{methode}[Comment mémoriser un caractère de façon durable (Méthode « Clé + Phonétique + Geste »)]
Suivez ces 4 étapes à chaque nouveau caractère :
\begin{enumerate}
    \item \textbf{Identifiez la clé (radical) :} Cherchez quel côté (gauche, droite, haut, bas, encadrement) porte le sens. Nommez-la (\emph{ex: « clé métal 钅 », « clé coquillage 贝 », « clé bambou ⺮ », « clé œil 目 »}). Cela ancre le \textbf{sens}.
    \item \textbf{Repérez la partie phonétique :} Souvent à droite (gauche-droite) ou en haut (haut-bas). Lisez-la à voix haute. Même si la prononciation a divergé (ex: \zh{戋} jiān $\to$ \zh{钱} qián), la proximité historique aide. Cela ancre le \textbf{son}.
    \item \textbf{Décomposez en traits et écrivez en comptant :} Prenez un stylo, écrivez le caractère \textbf{3 fois de suite} en énonçant le numéro de chaque trait (\emph{« un, deux, trois... »}). Respectez scrupuleusement l'ordre standard. Cela ancre le \textbf{geste moteur} (mémoire musculaire).
    \item \textbf{Formez 2 mots composés immédiats :} Écrivez le caractère dans deux mots du vocabulaire vu en section 3 (\emph{ex: 钱 $\to$ 钱包, 零钱 ; 算 $\to$ 计算, 打算}). Cela ancre l'\textbf{usage}.
\end{enumerate}
\textbf{Conseil pro :} Utilisez du papier à carreaux chinois (米字格 \zh{mǐzìgé}) pour respecter les proportions : la clé gauche prend 30-40\% de la largeur, la clé haut prend 30-40\% de la hauteur.
\end{methode}

\begin{attention}[Piège majeur : \zh{省} a deux sens mais une seule lecture \emph{shěng}]
\begin{itemize}
    \item \zh{省} \textbf{shěng} (3e ton) = \textbf{verbe} : économiser, réduire (\emph{省钱} \zh{shěng qián} — économiser de l'argent ; \emph{省事} \zh{shěng shì} — éviter des ennuis/travail).
    \item \zh{省} \textbf{shěng} (3e ton) = \textbf{nom} : province (division administrative) (\emph{广东省} \zh{Guǎngdōng Shěng} — province du Guangdong ; \emph{省会} \zh{shěnghuì} — chef-lieu de province).
    \item \textbf{Contexte seul décide :} \zh{我们要省钱} (Wǒmen yào \textbf{shěng qián}) = Nous devons économiser. \zh{他住在云南省} (Tā zhù zài Yúnnán \textbf{Shěng}) = Il habite dans la province du Yunnan.
    \item \textbf{Erreur fréquente :} Lire \emph{xǐng} (comme \zh{醒} « se réveiller ») par confusion visuelle. \zh{省} n'a \textbf{jamais} la lecture \emph{xǐng}.
\end{itemize}
\end{attention}

\begin{exemplebox}[Collocations indispensables du thème (HSK 3-4)]
\begin{center}
\begin{tabularx}{\linewidth}{|X|X|X|}
\hline
\rowcolor{popblueL} \textbf{Chinois} & \textbf{Pinyin} & \textbf{Français} \\
\hline
省钱 & shěng qián & économiser de l'argent (verbe + objet) \\
\hline
省下 & shěng xià & mettre de côté, économiser (résultat) \zh{省下钱来} \\
\hline
打算 & dǎ suàn & avoir l'intention de, projeter de (verbe) \zh{打算去中国} \\
\hline
计算 & jì suàn & calculer, calcul (nom/verbe formel) \zh{计算机} (ordinateur) \\
\hline
预算 & yù suàn & budget (nom) \zh{做预算} (établir un budget) \\
\hline
花费 & huā fèi & dépenser (argent/temps) ; frais (nom) \zh{花费很大} \\
\hline
费用 & fèi yòng & frais, coûts, dépenses (nom) \zh{学费} \zh{生活费} \\
\hline
工资 & gōng zī & salaire (nom) \zh{发工资} \\
\hline
收入 & shōu rù & revenus, recettes (nom) \zh{收入不稳定} \\
\hline
支出 & zhī chū & dépenses, sorties d'argent (nom) \zh{收支平衡} \\
\hline
结余 & jié yú & solde, reste (nom) \zh{年终结余} \\
\hline
\end{tabularx}
\end{center}
\textbf{Note syntaxique :} \zh{打算} et \zh{计算} partagent \zh{算} mais \zh{打算} (verbe d'intention) est suivi d'un verbe (\zh{打算去}), alors que \zh{计算} (verbe de calcul) est suivi d'un nombre ou d'un objet (\zh{计算总数}).
\end{exemplebox}

\courssub{Caractères proches à distinguer (Discrimination visuelle)}

Les apprenants francophones confondent souvent des caractères qui partagent des composants graphiques mais ont des clés (donc des sens) radicalement différents. Voici les deux paires critiques de cette leçon.

\begin{definition}[钱 \zh{qián} (argent) vs 线 \zh{xiàn} (fil, ligne, indice)]
\begin{itemize}
    \item \zh{钱} : Clé \zh{钅} (\textbf{métal/or}) à gauche. \textbf{Sens :} monnaie, argent. \emph{Mnémotechnique :} « L'argent est en \textbf{métal} (or, pièces). »
    \item \zh{线} : Clé \zh{纟} (\textbf{soie/fil}, variante de \zh{糸}) à gauche. \textbf{Sens :} fil, ligne, frontière, indice. \emph{Mnémotechnique :} « Le \textbf{fil} est en \textbf{soie}. »
    \item \textbf{Différence visuelle :} \zh{钅} a un trait relevé final (㇀) ; \zh{纟} a un trait oblique (丿) puis un point (丶) — forme « zigzag » typique de la soie.
    \item \textbf{Exemple piège :} \zh{银线} (fil d'argent / fil métallique) vs \zh{银钱} (argent en espèces, vieilli).
\end{itemize}
\end{definition}

\begin{definition}[费 \zh{fèi} (frais) vs 佛 \zh{fó} (Bouddha)]
\begin{itemize}
    \item \zh{费} : Clé \zh{贝} (\textbf{coquillage/argent}) en bas. Haut : \zh{弗} (phonétique). \textbf{Sens :} dépense, coûter.
    \item \zh{佛} : Clé \zh{亻} (\textbf{personne/homme}, variante de \zh{人}) à gauche. Droite : \zh{弗} (phonétique). \textbf{Sens :} Bouddha, bouddhisme.
    \item \textbf{Différence visuelle :} \zh{费} = structure \textbf{haut-bas} (大头小脚) ; \zh{佛} = structure \textbf{gauche-droite} (人字旁).
    \item \textbf{Contexte :} \zh{学费} (frais de scolarité) vs \zh{佛教} (bouddhisme). Ne confondez jamais « payer les frais » et « prier Bouddha » !
\end{itemize}
\end{definition}

\courssub{Décomposition graphique (Schémas TikZ)}

Les figures ci-dessous visualisent l'architecture « Clé + Phonétique » pour les deux caractères à structure gauche-droite (\zh{钱}) et haut-bas (\zh{费}).

\begin{popfigure}
\begin{tikzpicture}[
    node distance=0.8cm and 0.5cm,
    box/.style={draw=popblue, thick, rounded corners=3pt, fill=popblueL, minimum height=1.2cm, minimum width=2.2cm, align=center, font=\sffamily\small},
    arrow/.style={-Stealth, thick, poporange},
    labelbox/.style={font=\sffamily\small, text=popdark}
]
% 钱 decomposition
\node[box, align=center] (qian){\zh{钱}\\\emph{qián}\\(argent)};
\node[box, left=of qian, xshift=-1cm, align=center] (jin){\zh{钅}\\\emph{jīn}\\(métal/or)\\Clé sémantique};
\node[box, right=of qian, xshift=1cm, align=center] (jian){\zh{戋}\\\emph{jiān}\\(petit)\\Partie phonétique};
\draw[arrow] (jin) -- node[above, labelbox] {+} (qian);
\draw[arrow] (jian) -- node[above, labelbox] {+} (qian);
\node[labelbox, above=0.3cm of qian] {Structure Gauche-Droite (左右)};
\end{tikzpicture}
\legende{Décomposition de \zh{钱} (qián) : Clé \zh{钅} (métal) + Phonétique \zh{戋} (jiān).}
\end{popfigure}

\begin{popfigure}
\begin{tikzpicture}[
    node distance=0.8cm and 0.5cm,
    box/.style={draw=popgreen, thick, rounded corners=3pt, fill=popgreenL, minimum height=1.2cm, minimum width=2.2cm, align=center, font=\sffamily\small},
    arrow/.style={-Stealth, thick, poporange},
    labelbox/.style={font=\sffamily\small, text=popdark}
]
% 费 decomposition
\node[box, align=center] (fei){\zh{费}\\\emph{fèi}\\(frais)};
\node[box, above=of fei, align=center] (fu){\zh{弗}\\\emph{fú}\\(non)\\Partie phonétique};
\node[box, below=of fei, align=center] (bei){\zh{贝}\\\emph{bèi}\\(coquillage)\\Clé sémantique};
\draw[arrow] (fu) -- node[left, labelbox] {+} (fei);
\draw[arrow] (bei) -- node[left, labelbox] {+} (fei);
\node[labelbox, below=0.3cm of bei] {Structure Haut-Bas (上下)};
\end{tikzpicture}
\legende{Décomposition de \zh{费} (fèi) : Phonétique \zh{弗} (fú) + Clé \zh{贝} (bèi, coquillage/argent).}
\end{popfigure}

\courssub{Tableau récapitulatif des 4 caractères piliers}

Ce tableau synthétise les informations essentielles pour la révision et l'examen (écriture dictée, reconnaissance de clé, comptage de traits).

\begin{center}
\begin{tabularx}{\linewidth}{|>{\centering}X|>{\centering}X|>{\centering}X|>{\centering}X|>{\centering\arraybackslash}X|}
\hline
\rowcolor{popblueL} \textbf{Caractère} & \textbf{Pinyin / Sens} & \textbf{Clé (Radical)} & \textbf{Nb Traits} & \textbf{Structure + Ordre clé} \\
\hline
\zh{钱} & qián / argent & \zh{钅} (métal, gauche) & 10 & Gauche-Droite : \zh{钅} (5) puis \zh{戋} (5) \\
\hline
\zh{费} & fèi / frais & \zh{贝} (coquillage, bas) & 9* & Haut-Bas : \zh{弗} (5) puis \zh{贝} (4*) \\
\hline
\zh{省} & shěng / économiser ; province & \zh{目} (œil, bas) & 9 & Haut-Bas : \zh{少} (4) puis \zh{目} (5) \\
\hline
\zh{算} & suàn / calculer & \zh{⺮} (bambou, haut) & 14 & Haut-Milieu-Bas : \zh{⺮} (6) + \zh{目} (5) + \zh{廾} (3) \\
\hline
\end{tabularx}
\end{center}
\emph{* Note sur 费 : Le comptage officiel simplifié pour l'examen retient souvent 9 traits (弗5 + 贝简写4). L'écriture calligraphique complète de 贝 seul est 7 traits.}

\begin{exoresolu}[Entraînement à l'ordre des traits — Modèle guidé]
\textbf{Consigne :} Écrire le caractère \zh{算} \emph{suàn} en respectant l'ordre des 14 traits. Numéroter chaque trait sur votre copie.
\textbf{Solution détaillée (étape par étape) :}
\begin{enumerate}
    \item \textbf{Clé bambou \zh{⺮} (6 traits) :} ① — (horiz.) ② \textbackslash (oblique) ③ — (horiz.) ④ \textbackslash (oblique) ⑤ — (horiz.) ⑥ \textbullet{} (point). \emph{Astuce : les 6 traits forment un « M » aplati avec un point final.}
    \item \textbf{Œil \zh{目} (5 traits) :} ⑦ — (horiz. haut) ⑧ | (vert. centre) ⑨ — (horiz. milieu) ⑩ — (horiz. bas) (11) — (horiz. fermeture). \emph{Astuce : un rectangle divisé en deux par le trait ⑨, fermé par (11).}
    \item \textbf{Mains \zh{廾} (3 traits) :} (12) — (horiz. gauche, long) (13) — (horiz. droit, long) (14) | (vert. centre, relie les deux). \emph{Astuce : comme un « H » majuscule mais avec le trait vertical qui descend plus bas.}
\end{enumerate}
\textbf{Vérification :} Comptez : 6 + 5 + 3 = 14 traits. Le caractère final doit être bien centré, le \zh{⺮} large et plat, le \zh{目} bien carré, le \zh{廾} évasé vers le bas.
\end{exoresolu}

\begin{exercice}[Exercices d'écriture et de discrimination (À faire sur cahier / papier 米字格)]
\begin{enumerate}
    \item \textbf{Exercice de tracé (Obligatoire) :} Écrire chaque caractère des 4 piliers (\zh{钱}, \zh{费}, \zh{省}, \zh{算}) \textbf{5 fois de suite} en respectant l'ordre des traits. Cercler votre plus belle occurrence pour chaque caractère.
    \item \textbf{Dictée de caractères (Oral $\to$ Écrit) :} L'enseignant dicte : \emph{qián, fèi, shěng, suàn, guì, gòu, cái}. Écrire le caractère correct + pinyin + traduction.
    \item \textbf{Chasse à la clé :} Dans la liste ci-dessous, entourez la clé \zh{贝} et donnez le sens du caractère : \zh{贵}, \zh{赔}, \zh{贷}, \zh{资}, \zh{贝}, \zh{贿}, \zh{责}, \zh{贸}. (\emph{Indice : \zh{资} et \zh{责} ont \zh{贝} à droite ; \zh{贸} a \zh{贝} en bas}).
    \item \textbf{Discrimination visuelle :} Choisissez le bon caractère pour compléter :
    \begin{enumerate}
        \item Je n'ai pas assez d'\underline{\hspace{2cm}} (argent) pour acheter ce téléphone. (\zh{钱} / \zh{线})
        \item Les \underline{\hspace{2cm}} (frais) de scolarité augmentent cette année. (\zh{费} / \zh{佛})
        \item Nous devons \underline{\hspace{2cm}} (économiser) l'eau et l'électricité. (\zh{省} / \zh{眠})
        \item As-tu \underline{\hspace{2cm}} (calculé) le total ? (\zh{算} / \zh{管})
    \end{enumerate}
\end{enumerate}
\end{exercice}

\begin{corrige}[Exercices d'écriture et de discrimination]
\begin{enumerate}
    \item \textbf{Tracé :} Vérification par l'enseignant / auto-correction avec les schémas TikZ ci-dessus. Critères : ordre respecté, proportions équilibrées (clé fine/large selon structure), traits nets.
    \item \textbf{Dictée :}
    \begin{itemize}
        \item qián $\to$ \zh{钱} (argent)
        \item fèi $\to$ \zh{费} (frais)
        \item shěng $\to$ \zh{省} (économiser/province)
        \item suàn $\to$ \zh{算} (calculer)
        \item guì $\to$ \zh{贵} (cher/précieux) — clé \zh{贝}
        \item gòu $\to$ \zh{购} (acheter) — clé \zh{贝}
        \item cái $\to$ \zh{财} (richesse) — clé \zh{贝}
    \end{itemize}
    \item \textbf{Chasse à la clé \zh{贝} :}
    \begin{center}
    \begin{tabularx}{\linewidth}{|X|X|X|}
    \hline
    \rowcolor{popblueL} Caractère & Clé \zh{贝} position & Sens \\
    \hline
    \zh{贵} & Bas (sous \zh{古}) & Cher,贵重 \\
    \zh{赔} & Bas (sous \zh{非}) & Indemniser, perdre argent \\
    \zh{贷} & Bas (sous \zh{代}) & Prêter, emprunt \\
    \zh{资} & Droite (après \zh{次}) & Capital, fonds, ressources \\
    \zh{贝} & Unique & Coquillage, coquille \\
    \zh{贿} & Bas (sous \zh{丰}) & Pot-de-vin, corruption \\
    \zh{责} & Droite (après \zh{朿}) & Réprimander, responsabilité \\
    \zh{贸} & Bas (sous partie haute) & Commerce, négoce \\
    \hline
    \end{tabularx}
    \end{center}
    \item \textbf{Discrimination :}
    \begin{enumerate}
        \item \zh{钱} (argent — clé métal). \zh{线} = fil.
        \item \zh{费} (frais — clé coquillage). \zh{佛} = Bouddha.
        \item \zh{省} (économiser — clé œil). \zh{眠} = dormir (clé œil aussi mais sens diff.).
        \item \zh{算} (calculer — clé bambou). \zh{管} = gérer/tube (clé bambou aussi).
    \end{enumerate}
\end{enumerate}
\end{corrige}

\begin{savaistu}[Le cauri (贝 bèi) : la première monnaie mondiale ?]
Bien avant les pièces de métal et le papier-monnaie (inventé en Chine sous les Tang/Song), les \textbf{cauris} (coquillages marins blancs, brillants, de la famille des Cypraeidae) circulaient comme monnaie d'échange à travers l'Afrique, l'Asie du Sud, l'Océanie et la Chine antique. Au Cameroun, comme en Afrique de l'Ouest, le cauri (\emph{nkumu} en bassa, \emph{kauri} en peul) a servi de monnaie traditionnelle jusqu'à la colonisation. Le caractère \zh{贝} est donc un \textbf{témoin historique universel} : il relie la Chine ancienne aux sociétés africaines précoloniales par le même objet monétaire. Quand vous écrivez \zh{费}, \zh{贵}, \zh{财}, vous tracez un symbole qui a traversé les océans et les millénaires !
\end{savaistu}

\coursec{Compréhension du texte et collocations}

Après avoir analysé l'écriture et les radicaux des mots-clés du thème, nous passons maintenant à l'exploitation communicative du texte. Cette section vise à vérifier votre compréhension fine, à fixer les collocations lexicales indispensables pour parler d'argent et de budget, et à maîtriser la structure syntaxique majeure \zh{花 + somme + 买 + objet}. Nous suivrons la méthode « observer $\rightarrow$ comprendre $\rightarrow$ s'entraîner $\rightarrow$ produire ».

\courssub{Relecture active du texte support}

Relisez attentivement le texte ci-dessous (déjà vu en section 2). Concentrez-vous sur les chiffres, les verbes d'action financière (\zh{付}, \zh{存}, \zh{花}) et les connecteurs logiques (\zh{先…然后…}, \zh{因为…所以…}).

\begin{textebox}[Texte support : La gestion du budget familial]
我爸爸每个月工资五千块钱。每次发工资的时候，他们先付房租和电费，然后买菜、买米、买日用品。妈妈很节俭，她每个月都把剩下的钱存起来，存在银行里。我不想乱花钱，我想找一份兼职工作，因为我想买一个新手机，也不想让爸爸妈妈太累。
\\
Wǒ bàba měi gè yuè gōngzī wǔ qiān kuài qián. Měi cì fā gōngzī de shíhou, tāmen xiān fù fángzū hé diànfèi, ránhòu mǎi cài, mǎi mǐ, mǎi rìyòngpǐn. Māma hěn jiéjiǎn, tā měi gè yuè dōu bǎ shèngxià de qián cún qǐlai, cún zài yínháng lǐ. Wǒ bù xiǎng luàn huā qián, wǒ xiǎng zhǎo yī fèn jiānzhí gōngzuò, yīnwèi wǒ xiǎng mǎi yī gè xīn shǒujī, yě bù xiǎng ràng bàba māma tài lèi.
\\
Mon père gagne cinq mille francs par mois. Chaque fois que le salaire est versé, ils paient d'abord le loyer et l'électricité, puis achètent des légumes, du riz et des articles de première nécessité. Maman est très économe : chaque mois, elle met l'argent restant de côté, à la banque. Je ne veux pas dépenser n'importe comment ; je veux trouver un travail à temps partiel, parce que je veux acheter un nouveau téléphone, et je ne veux pas non plus que papa et maman soient trop fatigués.
\end{textebox}

\begin{center}
\begin{tabularx}{\linewidth}{|X|X|X|X|}
\hline
\rowcolor{popblueL} \textbf{Caractères} & \textbf{Pinyin} & \textbf{Nature} & \textbf{Traduction} \\
\hline
工资 & gōngzī & n. & salaire \\
\hline
发工资 & fā gōngzī & loc. v. & toucher sa paie ; verser le salaire \\
\hline
房租 & fángzū & n. & loyer \\
\hline
电费 & diànfèi & n. & facture d'électricité \\
\hline
日用品 & rìyòngpǐn & n. & articles de première nécessité \\
\hline
节俭 & jiéjiǎn & adj. & économe, frugal \\
\hline
剩下 & shèngxià & v. & rester (ce qui reste) \\
\hline
存起来 & cún qǐlai & v. & mettre de côté, épargner \\
\hline
乱花钱 & luàn huā qián & loc. v. & dépenser n'importe comment, gaspiller \\
\hline
兼职 & jiānzhí & n. & travail à temps partiel \\
\hline
累 & lèi & adj. & fatigué \\
\hline
\end{tabularx}
\end{center}

\begin{aretenir}[Le classificateur de l'argent : \zh{块}]
Pour compter l'argent en chinois courant, on utilise \zh{块} (kuài), équivalent oral de « francs » : \zh{五千块钱} (wǔ qiān kuài qián) « cinq mille francs ». Structure : \textbf{nombre + \zh{块} + \zh{钱}}. À l'écrit formel, on utilise \zh{元} (yuán), la monnaie chinoise.
\end{aretenir}

\courssub{Questions de compréhension détaillée}

Répondez en phrases complètes en chinois. La méthode est rappelée dans la boîte \emph{Méthode} plus bas. Voici les réponses modèles attendues.

\begin{exemplebox}[Réponses guidées aux 4 questions officielles]
\textbf{1. 爸爸每个月的工资是多少？} (Bàba měi gè yuè de gōngzī shì duōshao ?)\\
\emph{Repérage :} « 我爸爸每个月工资五千块钱 ».\\
\emph{Réponse :} \zh{爸爸每个月的工资是五千块钱。} (Bàba měi gè yuè de gōngzī shì wǔ qiān kuài qián.) — Le salaire de papa est de cinq mille francs par mois.

\textbf{2. 他们家先付什么？} (Tāmen jiā xiān fù shénme ?)\\
\emph{Repérage :} « 他们先付房租和电费 ».\\
\emph{Réponse :} \zh{他们先付房租和电费。} (Tāmen xiān fù fángzū hé diànfèi.) — Ils paient d'abord le loyer et l'électricité.

\textbf{3. 妈妈每个月做什么？} (Māma měi gè yuè zuò shénme ?)\\
\emph{Repérage :} « 她每个月都把剩下的钱存起来 ».\\
\emph{Réponse :} \zh{妈妈每个月把剩下的钱存起来。} (Māma měi gè yuè bǎ shèngxià de qián cún qǐlai.) — Maman met l'argent restant de côté chaque mois.

\textbf{4. “我”想做什么？为什么？} (« Wǒ » xiǎng zuò shénme ? Wèishénme ?)\\
\emph{Repérage :} « 我想找一份兼职工作，因为我想买一个新手机 ».\\
\emph{Réponse :} \zh{我想找一份兼职工作，因为我想买一个新手机。} (Wǒ xiǎng zhǎo yī fèn jiānzhí gōngzuò, yīnwèi wǒ xiǎng mǎi yī gè xīn shǒujī.) — Je veux trouver un travail à temps partiel parce que je veux acheter un nouveau téléphone.
\end{exemplebox}

\courssub{Collocations essentielles : l'argent en action}

Ces associations verbe + objet sont des « blocs » lexicaux à mémoriser par cœur. Elles permettent de passer du mot isolé à la phrase correcte.

\begin{center}
\begin{tabularx}{\linewidth}{|X|X|X|X|}
\hline
\rowcolor{poporangeL} \textbf{Collocation} & \textbf{Pinyin} & \textbf{Structure} & \textbf{Traduction} \\
\hline
花钱 & huā qián & v. + obj. & dépenser de l'argent \\
\hline
省钱 & shěng qián & v. + obj. & économiser de l'argent \\
\hline
存钱 & cún qián & v. + obj. & épargner, mettre de l'argent à la banque \\
\hline
付钱 & fù qián & v. + obj. & payer (l'acte de payer) \\
\hline
做预算 & zuò yùsuàn & v. + obj. & faire un budget \\
\hline
乱花钱 & luàn huā qián & adv. + v. + obj. & dépenser n'importe comment, gaspiller \\
\hline
赚钱 & zhuàn qián & v. + obj. & gagner de l'argent \\
\hline
借钱 & jiè qián & v. + obj. & emprunter de l'argent \\
\hline
还钱 & huán qián & v. + obj. & rembourser de l'argent \\
\hline
\end{tabularx}
\end{center}

\begin{propriete}[{Structure clé : \zh{花 + [somme d'argent] + 买 + [objet]}}]
Le verbe \zh{花} (huā) signifie « dépenser (de l'argent, du temps) ». Quand il s'agit d'argent, la \textbf{somme dépensée} se place \textbf{immédiatement après le verbe} \zh{花}, avant le verbe d'achat \zh{买}.
\par
\textbf{Schéma :} Sujet + \zh{花} + somme + (\zh{买}) + objet
\par
\textbf{Exemples :}
\begin{itemize}
    \item \zh{我花五百块买书。} (Wǒ huā wǔ bǎi kuài mǎi shū.) — Je dépense cinq cents francs pour acheter des livres.
    \item \zh{他花了两千块买手机。} (Tā huā le liǎng qiān kuài mǎi shǒujī.) — Il a dépensé deux mille francs pour acheter un téléphone (aspect accompli \zh{了}).
    \item \zh{别花太多钱买衣服。} (Bié huā tài duō qián mǎi yīfu.) — Ne dépense pas trop d'argent pour des vêtements.
\end{itemize}
\par
\textbf{Comparaison FR/CN :} en français, on dit « dépenser cinq cents francs \textbf{pour} acheter ». En chinois, la préposition « pour » disparaît : la somme devient complément d'objet direct de \zh{花}, et \zh{买} introduit l'objet visé dans une construction en série verbale.
\par
\textbf{Attention au nombre :} pour « deux mille », on dit \zh{两千} (liǎng qiān) et non \zh{二千} : devant un classificateur ou une unité de mesure, « deux » se dit \zh{两}.
\end{propriete}

\begin{popfigure}
\begin{tikzpicture}[
    node distance=1.2cm and 0.8cm,
    box/.style={draw=popblue, thick, rounded corners, fill=popblueL, text width=2.2cm, align=center, font=\small},
    arrow/.style={-Stealth, thick, poporange},
    labelnode/.style={font=\footnotesize, text=popdark, align=center}
]
\node[box, align=center] (hua){\zh{花}\\(huā)\\dépenser};
\node[box, right=of hua, align=center] (somme){\zh{五百块}\\(wǔ bǎi kuài)\\somme : 500 F};
\node[box, right=of somme, align=center] (mai){\zh{买}\\(mǎi)\\acheter};
\node[box, right=of mai, align=center] (obj){\zh{书}\\(shū)\\objet : livres};
\draw[arrow] (hua.east) -- (somme.west) node[midway, above, labelnode] {combien ?};
\draw[arrow] (somme.east) -- (mai.west) node[midway, above, labelnode] {action};
\draw[arrow] (mai.east) -- (obj.west) node[midway, above, labelnode] {quoi ?};
\node[above=0.3cm of hua.north, anchor=south, font=\bfseries, text=popblue] {Structure : \zh{花 + somme + 买 + objet}};
\end{tikzpicture}
\legende{Schéma de la construction en série verbale avec \zh{花} : la somme est le complément d'objet direct de \zh{花}.}
\end{popfigure}

\begin{attention}[Piège majeur : \zh{花} — fleur ou dépenser ?]
Le caractère \zh{花} (huā, 1\ier{} ton) a deux emplois aux sens radicalement différents :
\begin{itemize}
    \item \zh{花} : \textbf{nom} = « fleur » (\zh{鲜花} xiānhuā, fleurs fraîches ; \zh{花园} huāyuán, jardin).
    \item \zh{花} : \textbf{verbe} = « dépenser » (argent, temps) : \zh{花钱} huā qián, \zh{花时间} huā shíjian.
\end{itemize}
\par
\textbf{Erreur fréquente :} traduire \zh{花钱} par « fleur-argent ».
\par
\textbf{Astuce :} regardez la place dans la phrase. Sujet + \zh{花} + \textbf{somme} $\rightarrow$ verbe « dépenser ». Nombre + classificateur + \zh{花} (ex. \zh{一朵花} yī duǒ huā, une fleur) $\rightarrow$ nom « fleur ».
\par
\textbf{Autre piège :} ne pas confondre \zh{存钱} (cún qián, épargner, mettre à la banque) et \zh{省钱} (shěng qián, économiser, ne pas dépenser). \zh{妈妈存钱。} = Maman met de l'argent à la banque. \zh{妈妈省钱。} = Maman fait des économies (elle dépense moins).
\end{attention}

\courssub{Méthode : répondre aux questions de compréhension écrite (HSK 3-4)}

\begin{methode}[Stratégie en 4 étapes pour les questions en \zh{什么}/\zh{谁}/\zh{为什么}/\zh{多少}]
\begin{enumerate}
    \item \textbf{Identifiez le mot interrogatif} : \zh{什么} (shénme, quoi), \zh{谁} (shéi, qui), \zh{为什么} (wèishénme, pourquoi), \zh{多少} (duōshao, combien), \zh{怎么} (zěnme, comment).
    \item \textbf{Localisez la zone du texte} : cherchez les mots-clés de la question (ex. \zh{爸爸}, \zh{工资}, \zh{先}, \zh{妈妈}, \zh{我}) dans le texte et soulignez la phrase cible.
    \item \textbf{Recopiez en adaptant} : reprenez la structure de la phrase du texte, supprimez le mot interrogatif et insérez l'information trouvée. Gardez la personne demandée par la question.
    \item \textbf{Vérifiez la cohérence} : la réponse est-elle une phrase complète (Sujet + Verbe + Complément) ? La marque d'aspect (\zh{了}, \zh{过}) est-elle justifiée ?
\end{enumerate}
\par
\textbf{Exemple d'application (Q4) :} question \zh{“我”想做什么？为什么？} $\rightarrow$ mots-clés : \zh{我}, \zh{想}, \zh{为什么} $\rightarrow$ zone : dernière phrase du texte $\rightarrow$ adaptation : \zh{我想找一份兼职工作，因为我想买一个新手机。}
\end{methode}

\courssub{Exemple résolu : production écrite guidée}

\begin{exoresolu}[Rédaction d'un mini-paragraphe « Mon budget »]
\textbf{Consigne :} en vous inspirant du texte support, écrivez 5 phrases en chinois (caractères + pinyin) pour décrire votre gestion de budget mensuel (argent de poche, petit travail, ou budget familial fictif). Utilisez au moins 3 collocations du tableau et la structure \zh{花…买…}.
\tcblower
\textbf{Production proposée :}
\begin{enumerate}
    \item \zh{我每个月有一千块钱生活费。} (Wǒ měi gè yuè yǒu yī qiān kuài qián shēnghuófèi.) — J'ai mille francs d'argent de poche par mois.
    \item \zh{我先存两百块钱。} (Wǒ xiān cún liǎng bǎi kuài qián.) — D'abord, je mets deux cents francs de côté.
    \item \zh{我花三百块买吃的。} (Wǒ huā sān bǎi kuài mǎi chī de.) — Je dépense trois cents francs pour acheter à manger.
    \item \zh{我不乱花钱买衣服。} (Wǒ bù luàn huā qián mǎi yīfu.) — Je ne gaspille pas mon argent en vêtements.
    \item \zh{我想省钱买电脑。} (Wǒ xiǎng shěng qián mǎi diànnǎo.) — Je veux économiser pour acheter un ordinateur.
\end{enumerate}
\textbf{Analyse :} phrases courtes et correctes ; vocabulaire du thème réemployé (\zh{存}, \zh{花}, \zh{省}, \zh{乱花钱}) ; structure \zh{花 + somme + 买} respectée (phrase 3) ; connecteur de séquence \zh{先} utilisé.
\end{exoresolu}

\courssub{Mini-dialogue de réemploi : au marché}

Entraînez-vous à l'oral par paires. Alternez les rôles A et B.

\begin{textebox}[Dialogue : Négocier et gérer son budget]
\textbf{A :} 这个书包多少钱？\\
Zhè gè shūbāo duōshao qián ? — Combien coûte ce sac à dos ?
\\
\textbf{B :} 一百五十块。\\
Yī bǎi wǔshí kuài. — Cent cinquante francs.
\\
\textbf{A :} 太贵了！我想省钱。能不能便宜一点？\\
Tài guì le ! Wǒ xiǎng shěng qián. Néng bu néng piányi yīdiǎn ? — C'est trop cher ! Je veux économiser. Pouvez-vous faire un prix ?
\\
\textbf{B :} 好吧，一百二十块。\\
Hǎo ba, yī bǎi èrshí kuài. — D'accord, cent vingt francs.
\\
\textbf{A :} 好，我花一百二十块买这个书包。\\
Hǎo, wǒ huā yī bǎi èrshí kuài mǎi zhè gè shūbāo. — Très bien, je dépense cent vingt francs pour acheter ce sac.
\end{textebox}

\begin{savaistu}[Négocier au marché : Cameroun et Chine]
Au Cameroun comme en Chine, marchander au marché est une pratique courante et attendue : demander \zh{能不能便宜一点？} (« pouvez-vous baisser un peu ? ») est tout à fait naturel. En revanche, dans les supermarchés des deux pays, les prix sont fixes. En Chine, le mot \zh{便宜} (piányi, « bon marché ») est l'un des premiers mots que les étrangers apprennent pour faire leurs courses !
\end{savaistu}

\courssub{Exercices d'application}

\begin{exercice}[Entraînement : compréhension, lexique, structure]
\textbf{Exercice 1. Complétez avec la collocation qui convient :} \zh{花钱}, \zh{省钱}, \zh{存钱}, \zh{付钱}, \zh{做预算}, \zh{乱花钱}
\begin{enumerate}
    \item 每个月初，我们都要 \_\_\_\_\_\_\_\_\_\_\_\_，计划开支。
    \item 别 \_\_\_\_\_\_\_\_\_\_\_\_，别买不需要的东西。
    \item 吃饭以后，请去前台 \_\_\_\_\_\_\_\_\_\_\_\_。
    \item 妈妈很节俭，每个月都 \_\_\_\_\_\_\_\_\_\_\_\_。
    \item 买房子要 \_\_\_\_\_\_\_\_\_\_\_\_ 很多钱。
\end{enumerate}

\textbf{Exercice 2. Transformez les phrases en utilisant la structure \zh{花 + somme + 买 + objet} :}
\begin{enumerate}
    \item 我买了一辆自行车，用了八百块。 $\rightarrow$ \_\_\_\_\_\_\_\_\_\_\_\_
    \item 妹妹想买裙子，大概需要两百块。 $\rightarrow$ \_\_\_\_\_\_\_\_\_\_\_\_ (utilisez \zh{想})
    \item 他们付了五万块买这辆车。 $\rightarrow$ \_\_\_\_\_\_\_\_\_\_\_\_
\end{enumerate}

\textbf{Exercice 3. Traduisez en chinois (caractères + pinyin) :}
\begin{enumerate}
    \item Mon père paie l'électricité et l'eau chaque mois.
    \item Ma sœur est très économe, elle met de l'argent de côté tous les mois.
    \item Je veux trouver un travail à temps partiel pendant les vacances d'été.
    \item Ne dépense pas n'importe comment ! Nous devons faire un budget.
\end{enumerate}
\end{exercice}

\begin{corrige}[Exercices d'application]
\textbf{Exercice 1.}
\begin{enumerate}
    \item \zh{做预算} — \zh{每个月初，我们都要做预算，计划开支。} (Měi gè yuè chū, wǒmen dōu yào zuò yùsuàn, jìhuà kāizhī.) — Au début de chaque mois, nous devons faire un budget et planifier les dépenses.
    \item \zh{乱花钱} — \zh{别乱花钱，别买不需要的东西。} (Bié luàn huā qián, bié mǎi bù xūyào de dōngxi.) — Ne gaspille pas ton argent, n'achète pas de choses inutiles.
    \item \zh{付钱} — \zh{吃饭以后，请去前台付钱。} (Chī fàn yǐhòu, qǐng qù qiántái fù qián.) — Après le repas, allez payer à la caisse.
    \item \zh{存钱} — \zh{妈妈很节俭，每个月都存钱。} (Māma hěn jiéjiǎn, měi gè yuè dōu cún qián.) — Maman est très économe, elle épargne chaque mois.
    \item \zh{花} — \zh{买房子要花很多钱。} (Mǎi fángzi yào huā hěn duō qián.) — Acheter une maison coûte beaucoup d'argent. (Ici, seul le verbe \zh{花} convient, car \zh{很多钱} est déjà son complément.)
\end{enumerate}

\textbf{Exercice 2.}
\begin{enumerate}
    \item \zh{我花八百块买了一辆自行车。} (Wǒ huā bā bǎi kuài mǎi le yī liàng zìxíngchē.) — J'ai dépensé huit cents francs pour acheter un vélo. (La marque d'accompli \zh{了} se place après \zh{买}, verbe qui exprime le résultat de l'achat.)
    \item \zh{妹妹想花两百块买裙子。} (Mèimei xiǎng huā liǎng bǎi kuài mǎi qúnzi.) — Ma petite sœur veut dépenser deux cents francs pour acheter une jupe.
    \item \zh{他们花五万块买了这辆车。} (Tāmen huā wǔ wàn kuài mǎi le zhè liàng chē.) — Ils ont dépensé cinquante mille francs pour acheter cette voiture.
\end{enumerate}

\textbf{Exercice 3.}
\begin{enumerate}
    \item \zh{我爸爸每个月付电费和水费。} (Wǒ bàba měi gè yuè fù diànfèi hé shuǐfèi.)
    \item \zh{我姐姐很节俭，每个月都存钱。} (Wǒ jiějie hěn jiéjiǎn, měi gè yuè dōu cún qián.) — On peut aussi dire \zh{把剩下的钱存起来} (mettre l'argent restant de côté).
    \item \zh{我想暑假找一份兼职工作。} (Wǒ xiǎng shǔjià zhǎo yī fèn jiānzhí gōngzuò.) — Le complément de temps \zh{暑假} se place avant le verbe.
    \item \zh{别乱花钱！我们要做预算。} (Bié luàn huā qián ! Wǒmen yào zuò yùsuàn.)
\end{enumerate}
\end{corrige}

\coursec{Production guidée et évaluation}

Après avoir travaillé la compréhension du texte et les collocations du thème « gestion du budget », nous passons à la mise en œuvre active. Cette section vise à transformer le vocabulaire et les structures observées en compétences productives : écrire un court paragraphe cohérent, le présenter à l'oral, simuler une situation de négociation réelle, et s'auto-évaluer grâce à une grille claire. Chaque activité est conçue pour préparer l'évaluation certificative (épreuves écrites et orales du baccalauréat).

\courssub{Production écrite guidée : « 我的预算 »}

Nous commençons par l'observation d'un modèle de production attendu. Ce texte court (5 phrases) mobilise le lexique de la leçon (revenus, dépenses classées, épargne, projet) et la structure temporelle \cle{先…再…} (d'abord… ensuite…).

\begin{textebox}[Modèle de production — 我的预算]
我每个月有两千块。我先付车费，再买菜。我不乱花钱，我想存钱买书。
\\
Wǒ měi ge yuè yǒu liǎng qiān kuài. Wǒ xiān fù chēfèi, zài mǎi cài. Wǒ bù luàn huā qián, wǒ xiǎng cún qián mǎi shū.
\\
Chaque mois, j'ai 2000 yuans. Je paie d'abord les frais de transport, ensuite j'achète à manger. Je ne dépense pas n'importe comment, je veux économiser pour acheter des livres.
\end{textebox}

\begin{methode}[Réussir sa production écrite « 我的预算 »]
\textbf{Étape 1 :} Choisir un montant mensuel réaliste (ex. : 1500, 2000, 3000 块).\\
\textbf{Étape 2 :} Lister 2 à 3 dépenses prioritaires en utilisant \cle{先…，再…，最后…} (d'abord…, ensuite…, enfin…). Utiliser le vocabulaire : \zh{房租} (fángzū, loyer), \zh{车费} (chēfèi, transport), \zh{吃饭} (chīfàn, repas), \zh{水电费} (shuǐdiànfèi, eau/électricité).\\
\textbf{Étape 3 :} Exprimer la discipline budgétaire avec \zh{不乱花钱} (bù luàn huā qián).\\
\textbf{Étape 4 :} Formuler un objectif d'épargne avec \cle{想存钱 + [verbe but]} (想存钱买书 / 想存钱去旅游).\\
\textbf{Étape 5 :} Relire : vérifier les tons (surtout \zh{省} shěng, \zh{算} suàn), l'ordre des mots (Sujet + Temps + Lieu/Manière + Verbe + Objet) et la présence d'au moins \textbf{5 mots du glossaire}.
\end{methode}

\begin{propriete}[Structures clés pour le budget]
\textbf{1. Exprimer la fréquence et le revenu :} \\
Schéma : \cle{Sujet + 每个月 + 有 + 数量 + 块/元} \\
Exemple : \zh{我每个月有一千五百元。} (Wǒ měi ge yuè yǒu yī qiān wǔ bǎi yuán.) — Je gagne 1500 yuans par mois.

\textbf{2. Ordonner les actions (chronologie) :} \\
Schéma : \cle{Sujet + 先 + Verbe 1 + Objet 1，再 + Verbe 2 + Objet 2} \\
Exemple : \zh{我先付房租，再买菜。} (Wǒ xiān fù fángzū, zài mǎi cài.) — Je paie d'abord le loyer, ensuite j'achète des légumes.

\textbf{3. Dépenser / Économiser (Attention à la structure verbale) :} \\
\cle{Sujet + 花 + Montant + Verbe d'action + Objet} $\rightarrow$ \zh{我花五十块买饭。} (Wǒ huā wǔshí kuài mǎi fàn.) \\
\cle{Sujet + 存 + Montant / 钱 + But} $\rightarrow$ \zh{我存五百块钱买手机。} (Wǒ cún wǔbǎi kuài qián mǎi shǒujī.)
\end{propriete}

\begin{center}
\begin{tabularx}{\linewidth}{|X|X|X|X|}
\hline
\rowcolor{popblueL} \textbf{Caractères} & \textbf{Pinyin} & \textbf{Nature} & \textbf{Traduction} \\
\hline
收入 & shōurù & nom/verbe & revenu / percevoir des revenus \\
\hline
支出 & zhīchū & nom/verbe & dépense / dépenser \\
\hline
房租 & fángzū & nom & loyer \\
\hline
车费 & chēfèi & nom & frais de transport \\
\hline
水电费 & shuǐdiànfèi & nom & facture eau/électricité \\
\hline
生活费 & shēnghuófèi & nom & argent de poche / frais de vie \\
\hline
存钱 & cún qián & verbe + objet & économiser de l'argent \\
\hline
乱花钱 & luàn huā qián & verbe + objet & dépenser n'importe comment \\
\hline
预算 & yùsuàn & nom/verbe & budget / budgétiser \\
\hline
结余 & jiéyú & nom/verbe & solde / reste (budget) \\
\hline
\end{tabularx}
\end{center}

\begin{exoresolu}[Production écrite : « 我的预算 »]
\textbf{Consigne :} Rédigez votre budget mensuel (4 à 6 phrases) en vous inspirant du modèle. Montant imaginaire : 1800 块. Dépenses : loyer 800, transport 200, nourriture 400. Projet : acheter un vélo (自行车 zìxíngchē).

\textbf{Production élève (brouillon) :}
我每个月有一千八百块。我先付房租八百块，再付车费两百块，最后买吃的四百块。我不乱花钱。我想存钱买自行车。

\textbf{Analyse et correction :}
\begin{itemize}
    \item \textbf{Phrase 1 :} Correcte. \zh{一千八百块} (yī qiān bā bǎi kuài).
    \item \textbf{Phrase 2 :} Structure \cle{先…再…最后…} bien utilisée. Attention : \zh{付房租} (payer le loyer) est correct. L'ajout du montant après le verbe (\zh{付房租八百块}) est une structure courante (Verbe + Objet + Complément de quantité) mais on préfère souvent : \zh{付八百块房租} (payer 800 de loyer) ou \zh{房租八百块} (le loyer [est] 800). \textbf{Version améliorée :} \zh{我先付八百块房租，再付两百块车费，最后花四百块买菜。}
    \item \textbf{Phrase 3 :} \zh{我不乱花钱} parfait.
    \item \textbf{Phrase 4 :} \zh{我想存钱买自行车} parfait. Utilisation de \zh{存钱} + but.
\end{itemize}
\textbf{Version finale validée :}
\zh{我每个月有一千八百块。我先付八百块房租，再付两百块车费，最后花四百块买菜。我不乱花钱，我想存四百块钱买自行车。}
\\
Wǒ měi ge yuè yǒu yī qiān bā bǎi kuài. Wǒ xiān fù bā bǎi kuài fángzū, zài fù liǎng bǎi kuài chēfèi, zuìhòu huā sì bǎi kuài mǎi cài. Wǒ bù luàn huā qián, wǒ xiǎng cún sì bǎi kuài qián mǎi zìxíngchē.
\end{exoresolu}

\begin{exercice}[À vous d'écrire : Mon budget]
Rédigez votre propre paragraphe « 我的预算 » (5 phrases minimum) sur votre cahier.
\begin{enumerate}
    \item Montant mensuel : \zh{我每个月有 \_\_\_\_\_ 块。}
    \item 3 dépenses prioritaires avec \zh{先…，再…，最后…} (utilisez \zh{付}, \zh{花}, \zh{买}).
    \item Discipline : \zh{我不乱花钱。}
    \item Projet d'épargne : \zh{我想存钱 \_\_\_\_\_。} (acheter un téléphone \zh{手机}, voyager \zh{旅游}, payer les frais d'inscription \zh{学费}...).
\end{enumerate}
\textit{Contrainte : Utilisez au moins 6 mots du tableau de vocabulaire ci-dessus.}
\end{exercice}

\courssub{Production orale : présentation du budget}

L'expression orale continue valide la prononciation des tons (notamment les tons 2 et 3 sur \zh{存} cún, \zh{省} shěng, \zh{算} suàn) et la fluidité de la structure \cle{先…再…}.

\begin{experience}[Activité orale : « Mon budget à mon voisin »]
\textbf{Déroulement (10 min) :}
\begin{enumerate}
    \item \textbf{Préparation individuelle (2 min) :} Chaque élève relit sa production écrite, surligne les mots-clés (montants, verbes d'action) et s'entraîne à voix basse (murmure) en pointant les caractères.
    \item \textbf{Échange en binôme (4 min) :} L'élève A présente son budget à l'élève B \textbf{sans lire}, en le regardant dans les yeux. L'élève B écoute et coche sur une fiche : « Montant clair ? », « 先/再/最后 entendus ? », « Projet d'épargne compris ? », « Tons corrects (échantillon) ? ». Inversion des rôles.
    \item \textbf{Retour classe (4 min) :} 3 à 4 volontaires présentent leur budget au tableau. La classe valide la grille collective (voir grille ci-après). Le professeur note les points forts et les erreurs récurrentes (ex. : confusion \zh{付} fù / \zh{花} huā, ton de \zh{乱} luàn).
\end{enumerate}
\end{experience}

\courssub{Jeu de rôle : négociation au marché Mokolo}

Situation authentique camerounaise. Le marché Mokolo (Yaoundé) est un lieu incontournable où la négociation (讨价还价 tǎojià huánjià) est la norme. Ce jeu de rôle prépare l'épreuve orale d'interaction.

\begin{experience}[Jeu de rôle : « Au marché Mokolo »]
\textbf{Contexte :} Vous voulez acheter un sac à dos (书包 shūbāo) ou des baskets (运动鞋 yùndòngxié). Le vendeur annonce un prix élevé.
\textbf{Rôles :} A = Acheteur (étudiant, budget serré) | B = Vendeur (commerçant de Mokolo).
\textbf{Contraintes linguistiques obligatoires :}
\begin{itemize}
    \item Acheteur : \zh{太贵了！} (Tài guì le !) — \zh{便宜一点吧！} (Piányi yìdiǎn ba !) — \zh{我只有...块。} (Wǒ zhǐ yǒu... kuài.)
    \item Vendeur : \zh{这个质量好。} (Zhège zhìliàng hǎo.) — \zh{最低...块，不能再少了。} (Zuìdī... kuài, bù néng zài shǎo le.)
    \item Conclusion : Accord (\zh{成交！} Chéngjiāo !) ou rupture polie (\zh{下次再来。} Xiàcì zài lái.)
\end{itemize}
\textbf{Déroulement :} 3 min de préparation (notes autorisées en pinyin seulement), 2 min de jeu par binôme. Changement de rôles. Le professeur circule, note les interactions réussies pour débriefing.
\end{experience}

\begin{savaistu}[Le secret de la négociation en Chine]
Dans les marchés chinois (如 \zh{秀水街} Xiùshuǐ Jiē à Pékin ou les marchés de nuit), dire \cle{便宜一点吧！} (Piányi yìdiǎn ba ! — « Un peu moins cher ! ») avec le sourire et le particle \zh{吧} (ba, suggestion douce) est la clé. On ne dit pas « C'est trop cher » de façon agressive, on propose un prix : \zh{一百块行吗？} (Yībǎi kuài xíng ma ? — « 100 yuans, ça va ? »). Le marchand répond souvent \zh{不行，二百。} (Bùxíng, èrbǎi.) Le jeu continue jusqu'au \zh{成交} (accord) ou au départ feint (\zh{走着瞧} — « on verra » / je m'en vais) qui fait souvent baisser le prix final.
\end{savaistu}

\begin{attention}[Piège majeur : « 花 » n'est pas « dépenser \textbf{dans} »]
En français : « Je dépense \textbf{dans} la nourriture » / « Je dépense 5000 francs \textbf{pour} le loyer ».
En chinois : \textbf{Pas de préposition} après \zh{花} (huā).
\begin{center}
\begin{tabularx}{\linewidth}{|X|X|}
\hline
\rowcolor{poporangeL} \textbf{Incorrect (calque français)} & \textbf{Correct (chinois standard)} \\
\hline
\zh{我花在吃饭五十块。} & \zh{我花五十块买饭。} (Wǒ huā wǔshí kuài mǎi fàn.) \\
\hline
\zh{我花钱在书上。} & \zh{我花钱买书。} (Wǒ huā qián mǎi shū.) \\
\hline
\zh{他在衣服上花了很多钱。} & \zh{他在衣服上花了很多钱。} ✓ (Structure \zh{在...上花钱} acceptée pour le focus sur le domaine, mais \zh{花钱买衣服} plus naturel). \\
\hline
\end{tabularx}
\end{center}
\textbf{Règle d'or :} \cle{Sujet + 花 + Montant + Verbe d'action + Objet}. Le montant est complément d'objet direct de \zh{花}.
\end{attention}

\courssub{Correction collective : grille d'évaluation simple}

Cette grille permet une auto-évaluation et une co-évaluation rapides. Elle est alignée sur les critères du baccalauréat (LV2) : cohérence lexicale, correction morphosyntaxique, prononciation/tonalité.

\begin{center}
\begin{tabularx}{\linewidth}{|X|X|X|X|}
\hline
\rowcolor{poppurpleL} \textbf{Critère} & \textbf{4 pts : Maîtrisé} & \textbf{2 pts : En acquisition} & \textbf{0 pt : Non acquis} \\
\hline
\textbf{Vocabulaire thème} & $\ge$ 6 mots du glossaire employés correctement & 3--5 mots employés & $<$ 3 mots ou hors thème \\
\hline
\textbf{Structure \zh{先…再…}} & Utilisée correctement 2 fois min. (chronologie claire) & Utilisée 1 fois ou ordre confus & Absente ou ordre illogique \\
\hline
\textbf{Verbes \zh{花} / \zh{存} / \zh{付}} & Structure \zh{花+金额+V} et \zh{存+金额} parfaites & 1 erreur de structure (ex. préposition) & Confusion \zh{花}/\zh{付} ou structure fautive \\
\hline
\textbf{Tons (échantillon 5 mots)} & 0 erreur sur \zh{钱、费、省、算、预、结} & 1--2 erreurs de ton & $\ge$ 3 erreurs de ton \\
\hline
\textbf{Cohérence globale} & Budget réaliste, projet logique, phrases liées & Budget incohérent ou projet absent & Texte incompréhensible \\
\hline
\end{tabularx}
\end{center}

\begin{methode}[Utilisation de la grille en classe]
\begin{enumerate}
    \item L'élève s'auto-évalue sur 20 après sa production orale/écrite.
    \item Le binôme co-évalue (note sur 20).
    \item Le professeur valide/ajuste (note finale sur 20).
    \item Objectif Terminale : viser \textbf{$\ge$ 16/20} pour être prêt à l'examen.
\end{enumerate}
\end{methode}

\courssub{Synthèse de la leçon : carte mentale complétée}

Les élèves complètent collectivement cette carte mentale au tableau (ou sur fiche) pour visualiser les liens logiques du thème « Gestion du budget ». Elle structure la révision du Module 4.

\begin{popfigure}
\begin{tikzpicture}[
    mindmap,
    grow cyclic,
    every node/.style={concept, execute at begin node=\hskip0pt},
    concept color=popblue!80,
    level 1/.style={level distance=4.5cm, sibling angle=90, concept color=popblue!60},
    level 2/.style={level distance=3cm, sibling angle=45, concept color=popgreen!60},
    level 3/.style={level distance=2.5cm, sibling angle=30, concept color=poporange!60},
    text=popdark,
    font=\small
]
\node[concept, minimum size=2.2cm, align=center]{Gestion du budget\\预算管理}
    child[concept color=popblue!70] {node[concept, minimum size=1.8cm, align=center]{Revenus\\收入}
        child {node[concept, align=center]{Salaire\\工资}}
        child {node[concept, align=center]{Bourse\\奖学金}}
        child {node[concept, align=center]{Argent de poche\\零花钱}}}
    child[concept color=popgreen!70] {node[concept, minimum size=1.8cm, align=center]{Dépenses prioritaires\\必要支出}
        child {node[concept, align=center]{Loyer\\房租}}
        child {node[concept, align=center]{Transport\\车费}}
        child {node[concept, align=center]{Nourriture\\吃饭/买菜}}
        child {node[concept, align=center]{Eau/Élec\\水电费}}}
    child[concept color=poporange!70] {node[concept, minimum size=1.8cm, align=center]{Épargne\\存钱/储蓄}
        child {node[concept, align=center]{Ne pas gaspiller\\不乱花钱}}
        child {node[concept, align=center]{Calculer\\算账/预算}}
        child {node[concept, align=center]{Économiser\\省钱}}}
    child[concept color=poppink!70] {node[concept, minimum size=1.8cm, align=center]{Projets\\目标/计划}
        child {node[concept, align=center]{Acheter livres\\买书}}
        child {node[concept, align=center]{Voyager\\旅游}}
        child {node[concept, align=center]{Études\\学费}}
        child {node[concept, align=center]{Vélo/Téléphone\\买车/手机}}};
\end{tikzpicture}
\legende{Carte mentale synthétique : Revenu $\rightarrow$ Dépenses prioritaires $\rightarrow$ Épargne $\rightarrow$ Projet (ex. 买书). À compléter par les élèves avec pinyin et traduction.}
\end{popfigure}

\courssub{Devoirs}

\begin{enumerate}
    \item \textbf{Apprentissage lexical :} Réviser l'intégralité du glossaire de la leçon (tableau Section 3 + tableau ci-dessus). Test de vocabulaire (chinois $\rightarrow$ français et français $\rightarrow$ chinois + pinyin) en début de prochaine séance.
    \item \textbf{Écriture des caractères (traits) :} Tracer \textbf{5 fois chacun} les 4 caractères suivants sur votre cahier d'écriture (carreau), en respectant l'ordre des traits et les proportions :
    \begin{center}
    \begin{tabular}{cccc}
    \zh{钱} (qián) & \zh{费} (fèi) & \zh{省} (shěng) & \zh{算} (suàn) \\
    \textit{Radical :} \zh{钅} (or) & \textit{Radical :} \zh{扌} (main) & \textit{Radical :} \zh{目} (œil) & \textit{Radical :} \zh{竹} (bambou) \\
    \textit{Ordre :} gauche $\rightarrow$ droite & \textit{Ordre :} gauche $\rightarrow$ droite & \textit{Ordre :} haut $\rightarrow$ bas & \textit{Ordre :} haut $\rightarrow$ bas \\
    (钅 + 戈) & (扌 + 弗) & (少 + 目) & (竹 + 具)
    \end{tabular}
    \end{center}
    \item \textbf{Production écrite finale :} Reprendre votre texte « 我的预算 » corrigé en classe, l'écrire proprement (caractères + pinyin + traduction) pour le portfolio de chinois (note continue).
    \item \textbf{Préparation orale :} S'enregistrer (téléphone) en lisant votre budget + une phrase de négociation type Mokolo (\zh{太贵了，便宜一点吧！}). S'écouter, corriger les tons, réenregistrer si nécessaire.
\end{enumerate}

\begin{aretenir}[Bilan de la Leçon 25 — Compétences validées]
\begin{itemize}
    \item \textbf{Lexique :} Maîtrise du vocabulaire budgétaire (收入, 支出, 房租, 车费, 存钱, 乱花钱, 预算, 结余).
    \item \textbf{Grammaire :} Structure temporelle \cle{先…再…(最后…)} ; Verbe \cle{花} + montant direct (sans préposition) ; Verbe \cle{存} + but.
    \item \textbf{Écriture :} Caractères \zh{钱, 费, 省, 算} (radicaux, ordre des traits, composition).
    \item \textbf{Compréhension :} Lecture d'un texte authentique sur la gestion étudiante, réponse aux questions fermées/ouvertes.
    \item \textbf{Expression :} Production écrite guidée (5 phrases cohérentes) ; Interaction orale (présentation + négociation marchande).
    \item \textbf{Socioculturel :} Comparaison modes de vie étudiant Cameroun/Chine ; Codes de négociation au marché (Mokolo / Chine).
\end{itemize}
\end{aretenir}

\newpage

\coursec{Activité d'intégration}

\coursec{Leçon 25 — Vocabulaire et compréhension de texte : gestion du budget}

\courssub{Mise en route et découverte du thème}
\begin{experience}[Brainstorming : « Argent et vie quotidienne »]
En classe entière, au tableau. L'enseignant écrit \zh{钱} (qián) au centre. Les élèves proposent en chinois (caractères + pinyin) tous les mots liés à l'argent qu'ils connaissent : gagner, dépenser, économiser, billets, banque, prix, cher, bon marché, etc. L'enseignant classe les propositions en trois colonnes au fur et à mesure : \textbf{Revenus} (\zh{收入}、\zh{工资}、\zh{零用钱}), \textbf{Dépenses} (\zh{支出}、\zh{房租}、\zh{吃饭}、\zh{买衣服}), \textbf{Gestion} (\zh{存钱}、\zh{预算}、\zh{储蓄}). Cette carte mentale active le schéma lexical et prépare la lecture.
\end{experience}

\courssub{Texte support : lecture et écoute}
\begin{textebox}[Texte original : Le budget de Marie, étudiante à Yaoundé]
\zh{我叫玛丽，住在雅温得。我每月有一份兼职工作，收入五万法郎。} \\
\zh{我先付房租一万五千法郎，再交水电费五千法郎。} \\
\zh{吃饭和生活用品花一万二千法郎，交通费花八千法郎。} \\
\zh{学费比较贵，我花六千法郎。} \\
\zh{食品和房租是最重要的支出。} \\
\zh{最后我存了四千法郎，备用。我觉得这个预算比较合理。} \\
\emph{Wǒ jiào Mǎlì, zhù zài Yǎwēndé. Wǒ měiyuè yǒu yī fèn jiānzhí gōngzuò, shōurù wǔ wàn fǎláng.} \\
\emph{Wǒ xiān fù fángzū yī wàn wǔ qiān fǎláng, zài jiāo shuǐdiànfèi wǔ qiān fǎláng.} \\
\emph{Chīfàn hé shēnghuó yòngpǐn huā yī wàn èr qiān fǎláng, jiāotōngfèi huā bā qiān fǎláng.} \\
\emph{Xuéfei bǐjiào guì, wǒ huā liù qiān fǎláng.} \\
\emph{Shípǐn hé fángzū shì zuì zhòngyào de zhīchū.} \\
\emph{Zuìhòu wǒ cún le sì qiān fǎláng, bèiyòng. Wǒ juéde zhège yùsuàn bǐjiào hélǐ.}
« Je m'appelle Marie, j'habite à Yaoundé. J'ai un job étudiant chaque mois, revenu 50 000 francs. Je paie d'abord le loyer 15 000 francs, ensuite je paie l'eau/électricité 5 000 francs. Manger et produits de vie courante coûtent 12 000 francs, transport 8 000 francs. La scolarité est assez chère, je paie 6 000 francs. Nourriture et loyer sont les dépenses les plus importantes. Finalement j'ai épargné 4 000 francs, pour imprévus. Je trouve ce budget assez raisonnable. »
\end{textebox}

\courssub{Glossaire du thème}
\begin{definition}[Vocabulaire de base — Gestion du budget (HSK 3-4)]
\begin{center}
\begin{tabularx}{\linewidth}{|X|X|X|X|}\hline
\rowcolor{popblueL} \textbf{Caractères} & \textbf{Pinyin} & \textbf{Nature} & \textbf{Traduction} \\ \hline
收入 & shōurù & n./v. & revenu / gagner (de l'argent) \\ \hline
支出 & zhīchū & n./v. & dépense / dépenser \\ \hline
预算 & yùsuàn & n. & budget \\ \hline
储蓄 & chǔxù & n./v. & épargne / épargner \\ \hline
存钱 & cúnqián & v. (sep.) & mettre de l'argent de côté (épargner) \\ \hline
花费 & huāfèi & v./n. & dépenser / frais, coût \\ \hline
房租 & fángzū & n. & loyer \\ \hline
水电费 & shuǐdiànfèi & n. & facture eau + électricité \\ \hline
交通费 & jiāotōngfèi & n. & frais de transport \\ \hline
生活费 & shēnghuófèi & n. & frais de vie (nourriture, etc.) \\ \hline
学费 & xuéfèi & n. & frais de scolarité \\ \hline
兼职 & jiānzhí & n./adj. & job étudiant / à temps partiel \\ \hline
法郎 & fǎláng & n. & franc (FCFA) \\ \hline
万 & wàn & num. & dix mille (10 000) \\ \hline
千 & qiān & num. & mille (1 000) \\ \hline
合理 & hélǐ & adj. & raisonnable, logique \\ \hline
重要 & zhòngyào & adj. & important \\ \hline
贵 / 便宜 & guì / piányi & adj. & cher / bon marché \\ \hline
备用 & bèiyòng & v./adj. & réserver pour imprévus / de réserve \\ \hline
付 & fù & v. & payer (souvent + loyer, facture) \\ \hline
交 & jiāo & v. & payer / s'acquitter (factures, impôts) \\ \hline
所以 & suǒyǐ & conj. & donc, par conséquent \\ \hline
最后 & zuìhòu & adv. & finalement, en dernier \\ \hline
\end{tabularx}
\end{center}
\end{definition}

\courssub{Clés et ordre des traits}
\begin{propriete}[Analyse de caractères clés du thème]
\begin{center}
\begin{tabularx}{\linewidth}{|X|X|X|X|}\hline
\rowcolor{popblueL} \textbf{Caractère} & \textbf{Clé (Radical)} & \textbf{Nb traits} & \textbf{Ordre des traits (règles principales)} \\ \hline
房 (fáng) & \zh{户} (hù, porte) & 8 & Haut → bas (clé \zh{户} d'abord), puis \zh{方} à l'intérieur (horizontal avant vertical croisé) \\ \hline
水 (shuǐ) & \zh{水} (shuǐ, eau) & 4 & Trait central vertical, puis deux gouttes latérales gauche/droite, trait final horizontal \\ \hline
电 (diàn) & \zh{田} (tián, champ) & 5 & Cadre \zh{田} (extérieur avant intérieur), puis trait final \zh{乙} (bas) \\ \hline
交 (jiāo) & \zh{亠} (tóu, toit) & 6 & Haut \zh{亠}, puis \zh{父} (haut → bas : \zh{亠} + \zh{八} + horizontales) \\ \hline
学 (xué) & \zh{子} (zǐ, enfant) & 8 & Haut \zh{宀} (toit), puis \zh{子} (horizontal, vertical, horizontal, horizontal) \\ \hline
存 (cún) & \zh{子} (zǐ, enfant) & 6 & Gauche \zh{仓} (cāng, dépôt) simplifié : haut \zh{亠}, \zh{口}, \zh{一}, \zh{人} ; puis \zh{子} à droite \\ \hline
\end{tabularx}
\end{center}
\emph{Note :} L'ordre des traits suit les règles standards : horizontal avant vertical (\zh{十}), gauche avant droite (\zh{人}), extérieur avant intérieur (\zh{同}), fermer en dernier (\zh{口}).
\end{propriete}

\courssub{Compréhension du texte et collocations}
\begin{exercice}[Questions sur le texte « Le budget de Marie »]
\begin{enumerate}
    \item \textbf{Vrai ou Faux ?} Justifiez par une phrase du texte.
    \begin{itemize}
        \item[a)] Marie habite à Douala.
        \item[b)] Son revenu mensuel est de 50 000 FCFA.
        \item[c)] Elle paie l'eau/électricité avant le loyer.
        \item[d)] Le transport coûte 12 000 FCFA.
        \item[e)] Elle n'épargne rien.
    \end{itemize}
    \item \textbf{Questions ouvertes (réponse en chinois) :}
    \begin{itemize}
        \item[a)] 玛丽每月收入多少钱？
        \item[b)] 玛丽先付什么？再交什么？
        \item[c)] 为什么玛丽花六千法郎付学费？(Utilisez \zh{因为…所以…})
        \item[d)] 玛丽存了多少钱？为了什么？
        \item[e)] Selon Marie, quelles sont les deux dépenses les plus importantes ?
    \end{itemize}
    \item \textbf{Collocations :} Reliez le verbe à son objet logique (une flèche par ligne).
    \begin{center}
    \begin{tabular}{c c c}
    \zh{付} & $\rightarrow$ & \zh{房租 / 水电费 / 学费} \\
    \zh{交} & $\rightarrow$ & \zh{水电费 / 学费 / 罚款} \\
    \zh{买} & $\rightarrow$ & \zh{菜 / 生活用品 / 车票} \\
    \zh{存} & $\rightarrow$ & \zh{钱 / 钱在银行} \\
    \end{tabular}
    \end{center}
\end{enumerate}
\end{exercice}

\begin{corrige}[Exercice Compréhension]
\begin{enumerate}
    \item \textbf{V/F :} a) F (\zh{住在雅温得}) — b) V (\zh{收入五万法郎}) — c) F (\zh{先付房租，再交水电费}) — d) F (\zh{交通费花八千法郎}) — e) F (\zh{存了四千法郎}).
    \item \textbf{Réponses :}
    a) \zh{玛丽每月收入五万法郎。} \\
    b) \zh{先付房租，再交水电费。} \\
    c) \zh{因为学费比较贵，所以花六千法郎。} \\
    d) \zh{存了四千法郎，备用。} \\
    e) \zh{食品和房租是最重要的支出。}
    \item \textbf{Collocations :} \zh{付房租/学费}, \zh{交水电费}, \zh{买菜/生活用品}, \zh{存钱}. \textit{Note :} \zh{付} et \zh{交} s'échangent souvent pour factures ; \zh{付} insiste sur le versement, \zh{交} sur l'obligation administrative.
\end{enumerate}
\end{corrige}

\courssub{Structures grammaticales clés}
\begin{propriete}[Synthèse des structures pour « Gestion du budget »]
\begin{center}
\begin{tabularx}{\linewidth}{|X|X|X|}\hline
\rowcolor{popblueL} \textbf{Structure} & \textbf{Exemple (caractères + pinyin)} & \textbf{Traduction} \\ \hline
先 A 再 B (séquence) & 我们\cle{先}付房租\cle{再}交水电费。 & Nous payons \textbf{d'abord} le loyer, \textbf{ensuite} l'eau/électricité. \\ \hline
花 + montant + 买/付/在 + objet & 我\cle{花}八千法郎\cle{买}菜。 / \zh{花八千法郎在交通上。} & Je \textbf{dépense} 8 000 FCFA \textbf{pour acheter} à manger / \textbf{sur} le transport. \\ \hline
存 + 了 + montant (+ 钱) & 我们\cle{存了}四千法郎。 & Nous \textbf{avons épargné} 4 000 FCFA. (\zh{了} = action achevée) \\ \hline
比较 + adj. (comparatif léger) & 房租\cle{比较}贵。 & Le loyer est \textbf{assez} cher. \\ \hline
最 + adj. (superlatif) & 食品\cle{是最}重要的支出。 & La nourriture \textbf{est la} dépense \textbf{la plus} importante. \\ \hline
为什么…因为…/所以… & \cle{为什么}存钱少？\cle{因为}房租太贵。 & \textbf{Pourquoi} peu d'épargne ? \textbf{Parce que} le loyer est trop cher. \\ \hline
Subject + 时间 + Verb + Montant & 我们\cle{每月}收入\cle{五万法郎}。 & Nous gagnons \textbf{chaque mois} \textbf{50 000 francs}. \\ \hline
\end{tabularx}
\end{center}
\end{propriete}

\courssub{Activité d'intégration : Élaborer un budget familial}
\begin{textebox}[Fiche de travail — Budget mensuel d'une famille à Yaoundé]
\textbf{Contexte :} Vous êtes une famille habitant à \zh{雅温得} (Yaoundé). Le revenu mensuel net du foyer est de \zh{50 000 法郎} (50 000 FCFA). Vous devez établir le budget du mois prochain.

\textbf{Document de référence — Tableau des dépenses courantes (prix indicatifs Yaoundé 2024) :}
\begin{center}
\begin{tabularx}{\linewidth}{|X|X|X|}\hline
\rowcolor{popblueL} \textbf{Poste de dépense (中文)} & \textbf{Pinyin} & \textbf{Montant estimé (FCFA)} \\ \hline
房租 (loyer, chambre + salon) & fángzū & 15 000 -- 25 000 \\ \hline
水电费 (eau + électricité) & shuǐdiànfèi & 5 000 -- 8 000 \\ \hline
交通费 (transport : taxi, moto, bus) & jiāotōngfèi & 8 000 -- 12 000 \\ \hline
食品/生活费 (nourriture : riz, ndolé, poisson, marché) & shípǐn / shēnghuófèi & 12 000 -- 18 000 \\ \hline
学费 / 学杂费 (scolarité, fournitures) & xuéfèi / xuázáfèi & 5 000 -- 10 000 \\ \hline
医疗/意外 (pharmacie, consultation) & yīliáo / yìwài & 2 000 -- 5 000 \\ \hline
通讯费 (forfait téléphone/internet) & tōngxùnfèi & 2 000 -- 3 000 \\ \hline
娱乐/其他 (loisirs, coiffure, cadeaux) & yúlè / qítā & 2 000 -- 5 000 \\ \hline
\end{tabularx}
\end{center}

\textbf{Votre mission :} En groupe de 3, décidez de la répartition exacte (montants entiers, total = 50 000). Rédigez un \textbf{mini-rapport budgétaire} en chinois (4 à 6 phrases) qui suit la trame imposée. Préparez ensuite une \textbf{présentation orale de 2 minutes} (jeu de rôle : un parent présente, deux enfants/voisins posent des questions). La classe votera pour le \zh{最合理的预算} (budget le plus raisonnable).
\end{textebox}

\courssub{Consignes}
\begin{enumerate}
    \item \textbf{Analyse du document (5 min) :} En groupe, lisez le tableau. Identifiez les 5 à 6 postes indispensables pour votre famille. Calculez une répartition totale de 50 000 FCFA. Notez les montants choisis en chiffres arabes et en chinois (ex. : \zh{一万五千法郎}).
    \item \textbf{Rédaction du mini-rapport (15 min) :} Rédigez \textbf{ensemble} un texte de 4 à 6 phrases en caractères chinois (avec pinyin au brouillon). Le texte doit \textbf{obligatoirement} contenir :
    \begin{itemize}
        \item Une phrase d'ouverture avec le revenu : \zh{我们家每月收入五万法郎。}
        \item La structure \zh{先…再…} pour ordonner les priorités (ex. : \zh{我们先付房租，再交水电费。}).
        \item La structure \zh{花…买…} ou \zh{花…付…} au moins deux fois pour détailler des postes (ex. : \zh{我们花一万五千法郎付房租，花八千法郎买菜。}).
        \item Le verbe \zh{存} (avec \zh{了}) + montant (ex. : \zh{我们存了五千法郎。}).
        \item Un comparatif ou superlatif pour justifier un choix (ex. : \zh{房租比较贵，所以我们找便宜的房子。} ou \zh{食品是最重要的。}).
    \end{itemize}
    \item \textbf{Préparation de l'oral (10 min) :} Répartissez les rôles : \textbf{Rôle A} (père/mère) lit le rapport ; \textbf{Rôle B \& C} (enfants, voisin, conseiller budgétaire) préparent 3 questions chacun en chinois (ex. : \zh{房租多少钱？为什么不存更多钱？交通费怎么算的？}). Répétez la présentation à voix haute en soignant les tons.
    \item \textbf{Présentation devant la classe (2 min/groupe) :} Le groupe joue la scène. La classe note les montants entendus sur une grille d'écoute. À la fin, vote à main levée pour le \zh{最合理的预算} (critères : réalisme des montants, épargne non nulle, justifications claires).
    \item \textbf{Bilan individuel (5 min) :} Chaque élève écrit dans son cahier : (a) le vocabulaire nouveau qu'il a utilisé, (b) une structure qu'il a trouvée difficile, (c) une différence notable entre un budget type Yaoundé et un budget type Pékin/Shanghai (d'après les connaissances de la leçon).
\end{enumerate}

\courssub{Attention : Pièges fréquents pour francophones}
\begin{attention}[Erreurs types à éviter]
\begin{itemize}
    \item \textbf{Ordre des mots (Complément de quantité) :} En chinois, le montant se place \textbf{avant} le verbe (ou après avec \zh{了}). \textit{Erreur :} \zh{我们收入五万法郎每月。} $\rightarrow$ \textit{Correct :} \zh{我们每月收入五万法郎。} / \zh{我们收入了五万法郎。}
    \item \textbf{万 vs mille :} 50 000 = \zh{五万} (5 × 10 000), pas \zh{五十千}. 8 000 = \zh{八千}. 1 500 = \zh{一千五百} (ou \zh{一千五}).
    \item \textbf{花… + Verbe d'action :} \zh{花} (dépenser) appelle un verbe de réalisation (\zh{买} acheter, \zh{付} payer, \zh{在} sur). Ne pas dire \zh{买…花…}. \textit{Attention :} On ne « achète » pas un frais (\zh{买交通费} $\times$), on « paie » un frais (\zh{付交通费}) ou on « dépense pour » (\zh{在交通上}).
    \item \textbf{先…再… (Sujet unique) :} Les deux verbes concernent le \textbf{même sujet}. \textit{Erreur :} \zh{我先付房租，他再买菜。} $\rightarrow$ Si sujet change, utiliser \zh{然后} (ensuite) ou nouvelle phrase.
    \item \textbf{Verbe séparable 存钱 :} \zh{存} = verbe, \zh{钱} = objet. On n'ajoute pas de montant entre les deux. \textit{Correct :} \zh{存钱} (activité), \zh{存了五千法郎} (montant), \zh{存五千法郎钱} (montant + objet).
    \item \textbf{Tons critiques :} \zh{存 cún} (2e ton, montant) $\neq$ \zh{春 chūn} (1er ton, printemps) ; \zh{预 yù} (4e ton, pré-) $\neq$ \zh{雨 yǔ} (3e ton, pluie) ; \zh{算 suàn} (4e ton, calculer) $\neq$ \zh{酸 suān} (1er ton, acide). Entraînez-vous : \zh{预算 yùsuàn}, \zh{存钱 cúnqián}, \zh{比较 bǐjiào}.
\end{itemize}
\end{attention}

\courssub{Production guidée et évaluation}
\begin{exoresolu}[Mini-rapport budgétaire modèle — Groupe « Famille Mbarga »]
\textbf{Énoncé :} Rédiger un mini-rapport de 4 à 6 phrases respectant la trame imposée, pour un budget de 50 000 FCFA.

\tcblower
\textbf{Production modèle (caractères + pinyin + traduction) :}

\begin{enumerate}
    \item \zh{我们家每月收入五万法郎。} \\
    \emph{Wǒmen jiā měiyuè shōurù wǔ wàn fǎláng.} \\
    « Notre famille a un revenu mensuel de 50 000 francs. »
    
    \item \zh{我们\cle{先}付房租一万五千法郎，\cle{再}交水电费五千法郎。} \\
    \emph{Wǒmen \cle{xiān} fù fángzū yī wàn wǔ qiān fǎláng, \cle{zài} jiāo shuǐdiànfèi wǔ qiān fǎláng.} \\
    « Nous \textbf{payons d'abord} le loyer 15 000 francs, \textbf{ensuite} nous payons l'eau/électricité 5 000 francs. »
    
    \item \zh{我们\cle{花}一万二千法郎\cle{买}菜和生活用品，\cle{花}八千法郎\cle{付}交通费。} \\
    \emph{Wǒmen \cle{huā} yī wàn èr qiān fǎláng \cle{mǎi} cài hé shēnghuó yòngpǐn, \cle{huā} bā qiān fǎláng \cle{fù} jiāotōngfèi.} \\
    « Nous \textbf{dépensons} 12 000 francs \textbf{pour acheter} légumes et produits de première nécessité, \textbf{dépensons} 8 000 francs \textbf{pour payer} le transport. »
    
    \item \zh{学费比较贵，所以我们\cle{花}六千法郎付学费。} \\
    \emph{Xuéfei bǐjiào guì, suǒyǐ wǒmen \cle{huā} liù qiān fǎláng fù xuéfèi.} \\
    « La scolarité est \textbf{assez chère}, \textbf{donc} nous \textbf{dépensons} 6 000 francs pour payer la scolarité. »
    
    \item \zh{食品和房租\cle{是最}重要的支出。} \\
    \emph{Shípǐn hé fángzū \cle{shì zuì} zhòngyào de zhīchū.} \\
    « La nourriture et le loyer \textbf{sont les} dépenses \textbf{les plus} importantes. »
    
    \item \zh{最后我们\cle{存了}四千法郎，备用。} \\
    \emph{Zuìhòu wǒmen \cle{cún le} sì qiān fǎláng, bèiyòng.} \\
    « Finalement nous \textbf{avons épargné} 4 000 francs, pour imprévus. »
\end{enumerate}

\textbf{Vérification du total :} 15 000 + 5 000 + 12 000 + 8 000 + 6 000 + 4 000 = 50 000 FCFA. ✓

\textbf{Commentaire des choix de langue :}
\begin{itemize}
    \item \textbf{Phrase 1} : Structure standard \cle{Sujet + Temps + Verbe + Quantité}. \zh{每月} place le cadre temporel avant le verbe.
    \item \textbf{Phrase 2} : \cle{先…再…} ordonne logiquement les paiements fixes. \zh{交} (s'acquitter) est plus naturel que \zh{付} pour \zh{水电费}. Les montants suivent le verbe (complément de quantité post-verbal).
    \item \textbf{Phrase 3} : Double emploi de \cle{花…V…}. Distinction \zh{买} (biens physiques : \zh{菜}, \zh{生活用品}) vs \zh{付} (services/frais : \zh{交通费}). \zh{和} coordonne les objets.
    \item \textbf{Phrase 4} : Justification cause-effet \zh{比较贵} + \zh{所以} + action. Structure attendue en Terminale pour l'argumentation.
    \item \textbf{Phrase 5} : Superlatif \cle{最} + adjectif + \zh{的} + nom. \zh{是…的} structure focale (c'est bien X qui est…).
    \item \textbf{Phrase 6} : \zh{最后} marque la dernière étape. \zh{存了} (verbe + \zh{了} accompli) + montant direct. \zh{备用} (pour usage futur) ajoute une nuance pragmatique.
    \item \textbf{Tons soignés :} \zh{预算 yùsuàn} (4-4), \zh{存钱 cúnqián} (2-2), \zh{比较 bǐjiào} (3-4), \zh{重要 zhòngyào} (4-4), \zh{合理 hélǐ} (2-3), \zh{备用 bèiyòng} (4-4).
\end{itemize}

\textbf{Exemples de questions orales attendues (Rôles B \& C) :}
\begin{itemize}
    \item \zh{房租多少钱？} (Combien le loyer ?)
    \item \zh{为什么水电费只有五千？} (Pourquoi l'eau/électricité seulement 5 000 ?)
    \item \zh{你们存了四千，够不够？} (Vous avez épargné 4 000, est-ce suffisant ?)
    \item \zh{如果生病了，钱从哪里出？} (Si maladie, l'argent vient d'où ?)
    \item \zh{交通费八千，坐什么车？} (Transport 8 000, vous prenez quoi comme transport ?)
\end{itemize}

\textbf{Réponses types (Rôle A) :}
\begin{itemize}
    \item \zh{房租一万五，因为我们住两室一厅。} (Loyer 15 000, car nous habitons 2 chambres + salon.)
    \item \zh{水电费五千，我们省电，不用空调。} (Eau/élec 5 000, on économise l'électricité, pas de clim.)
    \item \zh{四千不多，但我们没办法，下个月再存。} (4 000 pas beaucoup, mais on ne peut pas faire autrement, le mois prochain on épargnera plus.)
    \item \zh{生病用备用金，或者借亲戚。} (Maladie : fonds d'urgence ou emprunter à la famille.)
    \item \zh{坐摩托车和小巴，比较便宜。} (Moto-taxi et mini-bus, assez bon marché.)
\end{itemize}
\end{exoresolu}

\courssub{Bilan et prolongements}
\begin{aretenir}[Ce que vous devez retenir de cette leçon]
\begin{enumerate}
    \item \textbf{Lexique ancré :} Vous savez nommer les postes d'un budget réel (\zh{房租}、\zh{水电费}、\zh{交通费}、\zh{生活费}、\zh{学费}、\zh{储蓄}) et manipuler les grands nombres en chinois (\zh{万}、\zh{千}、\zh{百}、\zh{法郎}).
    \item \textbf{Structures maîtrisées :} \cle{先…再…} pour prioriser, \cle{花…买/付/在…} pour détailler les dépenses, \cle{存了+montant} pour l'épargne accomplie, \cle{比较/最} pour évaluer. Ces schémas sont transférables à tout exposé sur la gestion financière.
    \item \textbf{Compétence communicative :} Passer d'un texte écrit préparé à une interaction orale spontanée (questions/réponses \zh{为什么}、\zh{多少钱}、\zh{怎么算}) est l'objectif final du programme Terminale. Le jeu de rôle force l'écoute active et la reformulation.
    \item \textbf{Auto-évaluation :} Si vous avez réussi à écrire le rapport sans dictionnaire, à le présenter avec des tons corrects (notamment \zh{存 cún}, \zh{预 yù}, \zh{算 suàn}), et à répondre aux questions \zh{为什么} sans hésiter, vous avez validé la compétence « Vocabulaire et compréhension de texte : gestion du budget ». Sinon, révisez le glossaire et refaites l'exercice à la maison en changeant les montants.
\end{enumerate}
\end{aretenir}

\begin{savaistu}[Regard interculturel : Budget Cameroun vs Chine]
\begin{itemize}
    \item \textbf{Pouvoir d'achat :} 50 000 FCFA ≈ 76 € / 580 CNY. Un étudiant chinois à Pékin/Shanghai dispose souvent de 2 500--4 000 元/mois (≈ 160 000--260 000 FCFA) versés par les parents (\zh{生活费}).
    \item \textbf{Postes majeurs :} Au Cameroun : \zh{房租} (15--25\%), \zh{食品} (25--35\%), \zh{交通} (15--20\%), \zh{学费} (variable). En Chine (grandes villes) : \zh{房租} souvent > 50\% du budget (colocation obligatoire), \zh{外卖} (livraison repas) très fréquent, \zh{地铁/公交} (métro/bus) bon marché et efficace, \zh{社保/医保} (assurance sociale) souvent couverte par l'employeur/université.
    \item \textbf{Épargne :} \zh{存钱} est une valeur confucéenne forte (\zh{积少成多} « accumuler peu pour faire beaucoup »). Les \zh{红包} (hóngbāo, enveloppes rouges) du Nouvel An ou mariages constituent une entrée d'argent ponctuelle importante, sans équivalent direct au Cameroun (où les \cle{cadeaux cérémoniels} / \cle{dot} jouent un rôle social différent).
    \item \textbf{Paiement :} La Chine est quasi \textbf{sans numéraire} (\zh{无现金社会}) : \zh{微信支付} (WeChat Pay) et \zh{支付宝} (Alipay) via QR code. Au Cameroun, le cash (FCFA) et \zh{Mobile Money} (Orange Money, MTN MoMo) cohabitent.
\end{itemize}
\end{savaistu}

\begin{methode}[Prolongements possibles pour l'évaluation ou le club de chinois]
\begin{enumerate}
    \item \textbf{Version numérique (TICE) :} Saisissez le budget dans un tableur avec colonnes en chinois : \zh{项目} (poste), \zh{预算} (prévu), \zh{实际} (réel), \zh{差额} (écart). Créez un \zh{饼图} (camembert) et commentez : \zh{房租占百分之三十，食品占百分之二十四。}
    \item \textbf{Débat contradictoire :} « \zh{存钱重要，还是享受生活重要？} » (Épargner vs Profiter). Deux équipes, 5 min préparation, 10 min débat. Vocabulaire imposé : \zh{未来} (avenir), \zh{现在} (présent), \zh{平衡} (équilibre), \zh{压力} (pression), \zh{幸福} (bonheur), \zh{及时行乐} (profiter de l'instant).
    \item \textbf{Enquête réelle :} Interrogez un parent/tuteur sur le vrai budget familial mensuel (anonymisé), notez les chiffres en chinois, présentez en classe. Comparez avec le modèle 50 000 FCFA.
    \item \textbf{Écriture créative :} Rédigez un dialogue « Au marché de Mfoundi » : marchander le prix du poisson/ndolé (\zh{太贵了，便宜点儿吧！}、\zh{一斤多少钱？}、\zh{给我打个折吧}).
\end{enumerate}
\end{methode}

\newpage

\coursec{Méthodes et exercices résolus}

\courssub{Méthodes et exercices résolus}

\begin{methode}[Lire et comprendre globalement un texte sur la gestion du budget]
\textbf{Objectif :} Saisir l'idée principale, identifier le type de document et le public visé sans connaître tous les mots.
\begin{enumerate}
    \item \textbf{Observer les indices paratextuels :} Titre, mise en page, mots-clés en gras (souvent le vocabulaire cible : \zh{预算}、\zh{收支平衡}).
    \item \textbf{Repérer les connecteurs logiques :} \zh{首先} (d'abord), \zh{然后} (ensuite), \zh{但是} (mais), \zh{因此} (donc) $\rightarrow$ structure du raisonnement.
    \item \textbf{Identifier la structure « Problème $\rightarrow$ Solution » :} Souvent : Constat (dépenses $>$ revenus) $\rightarrow$ Méthode (noter, classer) $\rightarrow$ Résultat (épargne).
    \item \textbf{Inférer le sens des mots inconnus :} Contexte + radicaux (ex: \zh{支} dans \zh{支出} et \zh{支付} $\rightarrow$ idée de « sortir/dépenser »).
    \item \textbf{Résumer en une phrase en français :} « Ce texte explique comment un étudiant gère son budget mensuel pour éviter le découvert. »
\end{enumerate}
\cle{Réflexe :} Ne pas traduire mot à mot. Lire le premier et le dernier paragraphe en priorité.
\end{methode}

\begin{exoresolu}[Compréhension globale — Type Bac]
\textbf{Texte support :} « \zh{小李的月度理财小妙招} » (Les petites astuces de Xiao Li pour la gestion mensuelle de ses finances).
\begin{textebox}[Texte original]
小李是喀麦隆雅温得大学的一名大三学生。每个月初，他都会制定一个详细的\zh{预算}。他的\zh{收入}主要来自奖学金和家里的资助，大约\zh{三十万}非洲法郎。\zh{固定支出}包括房租\zh{十五万}、交通\zh{两万}和网费\zh{一万}。变动支出是餐饮和娱乐。

小李坚持\zh{记账}：每天睡前用手机App记下每一笔\zh{开销}。上个月，他发现\zh{餐饮费}超支了\zh{五万}法郎，因为常和同学去吃\zh{烤鱼}。\zh{因此}，他决定\zh{调整}计划：每周只吃一次外面，其余时间在宿舍用电饭煲做简单的\zh{米饭}和\zh{青菜}。这个月，他终于实现了\zh{收支平衡}，还\zh{存下}了\zh{两万}法郎作为\zh{应急储蓄}。小李说：\zh{「理财不在于钱多，而在于有计划。」}
\end{textebox}
\textbf{Questions :}
\begin{enumerate}
    \item Quel est le type de ce document ? Justifiez par un élément du texte.
    \item Quelle est la nationalité et le statut du personnage principal ?
    \item Résumez la problème rencontré le mois dernier et la solution adoptée.
    \item Quelle est la leçon finale (morale) que tire Xiao Li ?
\end{enumerate}
\tcblower
\textbf{Solution détaillée :}
\begin{enumerate}
    \item \textbf{Type :} Témoignage personnel / Article de blog éducatif (récit à la 3e personne sur une expérience vécue). \textbf{Justification :} Présence de « \zh{小李说} » (Xiao Li dit) et style narratif chronologique (\zh{每个月初}、\zh{上个月}、\zh{这个月}).
    \item \textbf{Nationalité/Statut :} Camerounais (étudie à \zh{雅温得大学} - Université de Yaoundé), étudiant en 3e année (\zh{大三学生}).
    \item \textbf{Problème :} \zh{餐饮费超支了五万法郎} (Dépenses alimentaires dépassées de 50 000 F) à cause des sorties \zh{烤鱼} (poisson grillé). \textbf{Solution :} \zh{调整计划} : limiter resto à 1 fois/semaine, cuisiner soi-même (\zh{电饭煲做米饭青菜}).
    \item \textbf{Morale :} \zh{「理财不在于钱多，而在于有计划。」} (La gestion financière ne dépend pas de la quantité d'argent, mais d'avoir un plan). Structure \zh{不在于A，而在于B}.
\end{enumerate}
\cle{Conclusion :} La lecture globale permet de valider les items 1, 2, 4 sans connaître \zh{应急储蓄} ni \zh{电饭煲}.
\end{exoresolu}

\begin{methode}[Maîtriser le vocabulaire thématique : reconnaître, prononcer, classifier]
\textbf{Objectif :} Transformer le lexique passif (lu dans le texte) en lexique actif (utilisable à l'oral/écrit).
\begin{enumerate}
    \item \textbf{Tableau « Caractères — Pinyin — Nature — Traduction » :} Remplir le tableau ci-dessous (voir Glossaire leçon). Mémoriser les \textbf{tons} (ex: \zh{预算} yù\underline{suàn} 4-4, pas yúsuǎn).
    \item \textbf{Groupement par familles sémantiques (Mindmap) :}
        \begin{itemize}
            \item \textbf{Revenus :} \zh{收入}、\zh{奖学金}、\zh{资助}、\zh{工资}。
            \item \textbf{Dépenses :} \zh{支出}、\zh{开销}、\zh{房租}、\zh{水电费}、\zh{餐饮费}。
            \item \textbf{Verbes d'action :} \zh{制定}、\zh{记账}、\zh{超支}、\zh{调整}、\zh{存下}、\zh{平衡}。
            \item \textbf{Adjectifs/États :} \zh{固定} (fixe) vs \zh{变动} (variable), \zh{必需} vs \zh{奢侈}。
        \end{itemize}
    \item \textbf{Collocations indispensables (结构) :}
        \begin{itemize}
            \item \zh{制定 + 预算/计划} (établir un budget/plan)
            \item \zh{记账 / 记 + 开销 / 流水账} (noter les dépenses)
            \item \zh{收入 / 支出 + 平衡 / 超支} (équilibré / en déficit)
            \item \zh{存 + 钱 / 储蓄} (épargner)
        \end{itemize}
    \item \textbf{Entraînement « Pinyin $\rightarrow$ Caractères » :} Dictée quotidienne de 5 mots (application OnBuch+ ou cahier).
    \item \textbf{Production de micro-phrases :} Écrire 3 phrases personnelles avec les verbes d'action.
\end{enumerate}
\cle{Structure à retenir :} \zh{Subject + 对 + Objet + 进行 + 动词} (ex: \zh{他对预算进行了调整} — structure formelle HSK4).
\end{methode}

\begin{exoresolu}[Vocabulaire \& Collocations — Type Bac]
\textbf{Exercice :} Complétez les phrases avec le mot juste du tableau ci-dessous (forme correcte). Attention aux classificateurs et aux structures verbales.
\begin{center}
\begin{tabularx}{\linewidth}{|X|X|X|X|}
\hline \rowcolor{popblueL} \textbf{Caractères} & \textbf{Pinyin} & \textbf{Nature} & \textbf{Traduction} \\ \hline
预算 & yùsuàn & n. & budget \\ \hline
收入 & shōurù & n./v. & revenu / percevoir \\ \hline
支出 & zhīchū & n./v. & dépense / dépenser \\ \hline
记账 & jìzhàng & v. & tenir les comptes \\ \hline
超支 & chāozhī & v. & dépasser le budget \\ \hline
调整 & tiáozhěng & v./n. & ajuster / ajustement \\ \hline
储蓄 & chǔxù & n./v. & épargne / épargner \\ \hline
固定 & gùdìng & adj. & fixe \\ \hline
变动 & biàndòng & adj. & variable \\ \hline
必需品 & bìxūpǐn & n. & produit de première nécessité \\ \hline
\end{tabularx}
\end{center}
\begin{enumerate}
    \item 为了买新手机，他每个月都要存一部分 \zh{\_\_\_\_\_\_\_\_}。
    \item 房租属于 \zh{\_\_\_\_\_\_\_\_} 支出，每个月金额不变。
    \item 请 \zh{\_\_\_\_\_\_\_\_} 你的消费习惯，少买 \zh{\_\_\_\_\_\_\_\_}。
    \item 如果不 \zh{\_\_\_\_\_\_\_\_}，根本不知道钱花哪儿了。
    \item 这个月 \zh{\_\_\_\_\_\_\_\_} 了两千块，下个月得省着点花。
\end{enumerate}
\tcblower
\textbf{Solution détaillée :}
\begin{enumerate}
    \item \zh{储蓄} (chǔxù). \textbf{Justification :} « Épargner une partie de son \textbf{épargne} » $\rightarrow$ nom objet de \zh{存}. Phrase : \zh{为了买新手机，他每个月都要存一部分储蓄。}
    \item \zh{固定} (gùdìng). \textbf{Justification :} Adjectif qualifiant \zh{支出}. Antonyme de \zh{变动} cité dans le texte. \zh{房租属于固定支出。}
    \item \zh{调整} (tiáozhěng) + \zh{奢侈品} (shēchǐpǐn - hors liste, inféré par opposition à \zh{必需品}). \textbf{Justification :} \zh{调整} prend un objet (habitudes). \zh{少买奢侈品} (acheter moins de luxe). Phrase : \zh{请调整你的消费习惯，少买奢侈品。}
    \item \zh{记账} (jìzhàng). \textbf{Justification :} Verbe intransitif ici (ou \zh{记账} = action générale). \zh{如果不记账，根本不知道钱花哪儿了。}
    \item \zh{超支} (chāozhī). \textbf{Justification :} Verbe d'état/résultat. Sujet : \zh{这个月} (ce mois-ci). Complément de quantité : \zh{了两千块} (aspect \zh{了} + quantité). Phrase : \zh{这个月超支了两千块，下个月得省着点花。}
\end{enumerate}
\cle{Piège :} \zh{收入/支出} sont des noms \textbf{et} des verbes. \zh{超支} est verbe uniquement. \zh{记账} est verbe séparable (\zh{记} = verbe, \zh{账} = objet) $\rightarrow$ \zh{记了一笔账}.
\end{exoresolu}

\begin{methode}[Analyser les caractères : radicaux (clés) et ordre des traits]
\textbf{Objectif :} Comprendre la logique de construction pour mémoriser l'écriture et deviner le sens.
\begin{enumerate}
    \item \textbf{Identifier la clé (部首 bùshǒu) :} Elle donne l'indice sémantique (catégorie).
        \begin{itemize}
            \item \zh{贝} (bèi) = coquillage/argent $\rightarrow$ \zh{财} (richesse), \zh{账} (compte), \zh{资} (ressource), \zh{费} (frais), \zh{赚} (gagner).
            \item \zh{禾} (hé) = riz/récolte $\rightarrow$ \zh{预} (prévu/à l'avance), \zh{利} (profit/avantage).
            \item \zh{走} (zǒu) = marcher/aller $\rightarrow$ \zh{超} (dépasser), \zh{赴} (se rendre).
            \item \zh{亻} (rén) = humain $\rightarrow$ \zh{储} (réserve/épargner), \zh{借} (emprunter), \zh{贷} (prêter).
        \end{itemize}
    \item \textbf{Appliquer l'ordre des traits (笔顺 bǐshùn) :} Règles d'or : \textbf{Haut $\rightarrow$ Bas}, \textbf{Gauche $\rightarrow$ Droite}, \textbf{Horizontal avant Vertical} (si croisement), \textbf{Extérieur avant Intérieur}, \textbf{Traits centraux avant ailes}.
    \item \textbf{Exercices ciblés :}
        \begin{itemize}
            \item \zh{预} : \zh{禾} (5 traits) + \zh{予} (4 traits). Total 9. \textbf{Ordre :} 禾 (横撇竖横) $\rightarrow$ 予 (撇横竖钩).
            \item \zh{算} : \zh{竹} (6 traits, haut) + \zh{工} (3) + \zh{戈} (4). Total 13. \textbf{Ordre :} Bambou (haut) $\rightarrow$ 工 $\rightarrow$ 戈.
            \item \zh{储} : \zh{亻} (2, gauche) + \zh{宁} (5, droite). Total 7. \textbf{Ordre :} 人 (亻) $\rightarrow$ 宁 (宀 + 丁).
        \end{itemize}
    \item \textbf{Jeu « Caractères proches » (辨析) :} \zh{预} (yù, à l'avance) vs \zh{豫} (yù, province Henan / joie) ; \zh{账} (zhàng, compte) vs \zh{帐} (zhàng, tente — trad. \zh{帐篷}) ; \zh{费} (fèi, frais) vs \zh{废} (fèi, déchet).
\end{enumerate}
\cle{Réflexe écriture :} Écrire en disant le nom des traits (\zh{横、竖、撇、捺、钩}) pour ancrer la mémoire motrice.
\end{methode}

\begin{exoresolu}[Écriture \& Analyse de caractères — Type Bac]
\textbf{Exercice :}
\begin{enumerate}
    \item Pour les caractères \zh{预}、\zh{算}、\zh{储}、\zh{账} : indiquez la clé (radical), son sens général, le nombre total de traits.
    \item Parmi les paires suivantes, entourez le caractère correct pour le contexte « budget » et expliquez la différence de sens :
        \begin{itemize}
            \item \zh{预算} / \zh{豫算}
            \item \zh{记账} / \zh{记帐}
            \item \zh{费用} / \zh{废用}
        \end{itemize}
    \item Écrivez l'ordre des traits pour \zh{支} (zhī) et \zh{出} (chū). Combien de traits chacun ?
\end{enumerate}
\tcblower
\textbf{Solution détaillée :}
\begin{enumerate}
    \item \textbf{Tableau récapitulatif :}
    \begin{center}
    \begin{tabularx}{\linewidth}{|X|X|X|X|}
    \hline \rowcolor{popblueL} \textbf{Caractère} & \textbf{Clé (Radical)} & \textbf{Sens de la clé} & \textbf{Nb traits} \\ \hline
    \zh{预} & \zh{禾} (hé) & Riz, grain, récolte & 9 \\ \hline
    \zh{算} & \zh{竹} (zhú) & Bambou (bâtonnets de calcul) & 13 \\ \hline
    \zh{储} & \zh{亻} (rén) & Humain, personne & 7 \\ \hline
    \zh{账} & \zh{贝} (bèi) & Coquillage, monnaie, argent & 12 \\ \hline
    \end{tabularx}
    \end{center}
    \item \textbf{Choix justifiés :}
        \begin{itemize}
            \item \zh{预算} : \zh{预} = à l'avance (禾 = graines préparées). \zh{豫} = province Henan / joie. \textbf{Budget = calcul prévu à l'avance.}
            \item \zh{记账} : \zh{账} = comptes, dette (clé \zh{贝} argent). \zh{帐} = tente, voile (clé \zh{巾} tissu). \textbf{Tenir les comptes = 记账.}
            \item \zh{费用} : \zh{费} = frais, dépenser (clé \zh{贝} argent + \zh{弗} phonétique). \zh{废} = déchet, abandonné (clé \zh{广} toit + \zh{弗}). \textbf{Frais = 费用.}
        \end{itemize}
    \item \textbf{Ordre des traits :}
        \begin{itemize}
            \item \zh{支} (zhī) : 7 traits. Ordre : \zh{一} (横) $\rightarrow$ \zh{丨} (竖) $\rightarrow$ \zh{一} (横) $\rightarrow$ \zh{ノ} (撇) $\rightarrow$ \zh{丶} (点) $\rightarrow$ \zh{一} (横) $\rightarrow$ \zh{丨} (竖钩). \textit{Note : le trait central vertical traverse.}
            \item \zh{出} (chū) : 5 traits. Ordre : \zh{一} (横) $\rightarrow$ \zh{丨} (竖) $\rightarrow$ \zh{一} (横) $\rightarrow$ \zh{ノ} (撇) $\rightarrow$ \zh{乀} (捺/长横). \textit{Règle : cadre extérieur (haut de 山) puis trait final long.}
        \end{itemize}
\end{enumerate}
\cle{Point clé Bac :} La connaissance des clés \zh{贝} (argent) et \zh{禾} (plan/récolte) permet de deviner \zh{预算}、\zh{费用}、\zh{储蓄} même si on oublie le phonétique.
\end{exoresolu}

\begin{methode}[Répondre aux questions de compréhension détaillée (Explicite / Implicite)]
\textbf{Objectif :} Localiser l'information précise, reformuler en chinois ou français, gérer l'inférence.
\begin{enumerate}
    \item \textbf{Typologie des questions :}
        \begin{itemize}
            \item \textbf{Explicite (\zh{找信息}) :} « Selon le texte, quel est le montant du loyer ? » $\rightarrow$ Scanner mots-clés : \zh{房租}、\zh{十五万}. Réponse courte : \zh{十五万法郎} / \zh{150 000 F}.
            \item \textbf{Implicite (\zh{推断}) :} « Pourquoi Xiao Li a-t-il cuisiné lui-même ? » $\rightarrow$ Lien cause-effet : \zh{超支} (cause) $\rightarrow$ \zh{调整} (effet). Réponse : \zh{因为餐饮费超支了，所以他决定自己做饭。}
            \item \textbf{Vocabulaire en contexte (\zh{猜词}) :} « Que veut dire \zh{应急储蓄} ? » $\rightarrow$ \zh{应急} (urgence) + \zh{储蓄} (épargne) = « épargne de précaution ».
            \item \textbf{Opinion/Auteur (\zh{主旨}) :} « Êtes-vous d'accord avec la méthode de Xiao Li ? » $\rightarrow$ Structure : \zh{我同意。因为……另外……} (J'approuve. Parce que... De plus...).
        \end{itemize}
    \item \textbf{Méthode « Mots-clés $\rightarrow$ Phrase source $\rightarrow$ Réponse » :}
        \begin{enumerate}
            \item Surligner la question (mots-clés : \zh{谁}、\zh{什么}、\zh{为什么}、\zh{怎么}、\zh{多少}).
            \item Scanner le texte pour retrouver ces mots ou synonymes (\zh{房租} $\leftrightarrow$ \zh{住宿费}).
            \item Copier/Coller ou Reformuler (changer \zh{他} en \zh{小李}, \zh{这个月} en \zh{上个月}).
            \item Vérifier la cohérence grammaticale (chinois ou français selon consigne).
        \end{enumerate}
    \item \textbf{Structures de réponse types :}
        \begin{itemize}
            \item Cause : \zh{因为……所以……} / \zh{由于……因此……}
            \item But : \zh{为了……} / \zh{以便……}
            \item Contraste : \zh{虽然……但是……} / \zh{可是……}
        \end{itemize}
\end{enumerate}
\cle{Astuce :} Si la question est en français, répondre en français (sauf consigne « en chinois »). Si en chinois, répondre en chinois (phrases complètes).
\end{methode}

\begin{exoresolu}[Compréhension détaillée \& Inférence — Type Bac]
\textbf{Reprenez le texte « \zh{小李的月度理财小妙招} ». Répondez en chinois (phrases complètes).}
\begin{enumerate}
    \item 小李的收入来源有哪两个？(Quelles sont les deux sources de revenus ?)
    \item 上个月小李为什么超支？具体数字是多少？(Pourquoi dépassement le mois dernier ? Montant ?)
    \item 小李采取了什么具体措施来解决超支问题？(Mesures concrètes ?)
    \item 根据课文，判断正误并引用原文就证 : « 小李这个月没有存下钱。 » (Vrai/Faux + citation).
    \item 请用中文总结小李的理财经验 (30-40 caractères). (Résumer l'expérience en 30-40 car.)
\end{enumerate}
\tcblower
\textbf{Solution détaillée :}
\begin{enumerate}
    \item \zh{小李的收入来源是奖学金和家里的资助。} (Localisation : \zh{收入主要来自奖学金和家里的资助}).
    \item \zh{因为常和同学去吃烤鱼，餐饮费超支了五万法郎。} (Lien \zh{因为……所以……} implicite dans texte : \zh{因为...因此...}). \textbf{Justification :} \zh{上个月...发现餐饮费超支了五万法郎，因为常和同学去吃烤鱼。}
    \item \zh{他决定每周只吃一次外面，其余时间在宿舍用电饭煲做米饭和青菜。} (Reprise exacte : \zh{调整计划：每周只吃一次外面...做简单的米饭和青菜}).
    \item \zh{错误。} \textbf{Preuve :} \zh{「这个月，他终于实现了收支平衡，还存下了两万法郎作为应急储蓄。」} (Le texte dit explicitement qu'il a épargné 20 000).
    \item \zh{制定预算，坚持记账，及时调整，量入为出。} (16 car. + ponctuation). \textbf{Analyse :} \zh{制定预算} (planifier), \zh{坚持记账} (suivi), \zh{及时调整} (réagir), \zh{量入为出} (chengyu : dépenser selon ses revenus - HSK4+). \textit{Variante simple :} \zh{有计划地记账、调整开销，实现收支平衡。} (18 car.)
\end{enumerate}
\cle{Note correcteur :} Les réponses 1-3 reprennent le vocabulaire du texte (pas d'invention). Réponse 4 : citation obligatoire avec guillemets chinois \zh{「 」}. Réponse 5 : concision et vocabulaire ciblé.
\end{exoresolu}

\begin{methode}[Réinvestir le lexique et les structures en production guidée (Écrit/Oral)]
\textbf{Objectif :} Passer de la compréhension à l'expression personnelle (mon budget / conseil ami).
\begin{enumerate}
    \item \textbf{Cadre de rédaction « Mon budget mensuel » (\zh{表达个人情况}) :}
        \begin{enumerate}
            \item \textbf{Présentation :} \zh{我是……学生，每月收入约……元/法郎。} (Je suis ét., revenus ~).
            \item \textbf{Structure fixe :} \zh{固定支出有：房租……，交通……，网费……。} (Liste \zh{有} + items).
            \item \textbf{Variable :} \zh{变动支出主要是餐饮和……} (Principalement nourriture...).
            \item \textbf{Action :} \zh{我打算/决定……} (Je prévois/décide) + \zh{制定预算 / 坚持记账 / 减少奢侈品消费 / 自己做饭}).
            \item \textbf{Objectif :} \zh{希望能实现收支平衡，存下……作为储蓄。}
        \end{enumerate}
    \item \textbf{Cadre « Conseil à un ami » (\zh{给建议} - Oral/Écrit) :} Structures \zh{建议/应该/最好} + Verbe.
        \begin{itemize}
            \item \zh{我建议你先制定一个预算。} (Conseil général).
            \item \zh{你应该每天记账，知道钱花哪儿了。} (Obligation morale).
            \item \zh{最好少吃外卖，多自己做饭，既省钱又健康。} (Optimal + double avantage \zh{既……又……}).
            \item \zh{如果超支了，要及时调整计划。} (Conditionnel \zh{如果……就……} / \zh{要}).
        \end{itemize}
    \item \textbf{Connecteurs pour fluidifier (HSK3-4) :} \zh{首先……其次……最后……} ; \zh{一方面……另一方面……} ; \zh{不仅……而且……}.
    \item \textbf{Auto-correction (Check-list) :} Tons ? Classificateurs (\zh{一笔钱、一个月、些建议}) ? \zh{了} (achèvement/état nouveau) vs \zh{过} (expérience) ? Ordre : Sujet - Temps - Lieu - Manière - Verbe - Objet/Complément.
\end{enumerate}
\cle{Modèle oral (1 min) :} « Bonjour, je m'appelle X. Moi aussi je gère mon budget. Mon revenu c'est bourse 200 000 F. Loyer 100 000. Je note tout sur WeChat. Le mois dernier j'ai trop dépensé pour les jeux. Maintenant je cuisine. Je veux économiser 50 000. Merci. »
\end{methode}

\begin{exoresolu}[Production écrite guidée — Type Bac]
\textbf{Sujet :} Vous êtes \zh{王明} (Wáng Míng), étudiant à Douala. Écrivez un message WeChat (80-100 caractères chinois) à votre ami \zh{老张} (Lǎo Zhāng) pour lui expliquer votre nouvelle méthode de gestion de budget et l'inviter à faire de même.
\textbf{Contraintes :} Utiliser au moins 5 mots du vocabulaire de la leçon (\zh{预算、收入、支出、记账、超支、调整、储蓄、固定、变动、必需品}). Utiliser \zh{建议/应该/最好} et \zh{如果……就……}.
\tcblower
\textbf{Proposition de rédaction (Modèle corrigé) :}
\begin{textebox}[Message WeChat]
老张，你好！
最近我学会了理财。\zh{首先}，我\zh{制定}了详细的\zh{预算}：\zh{收入}是奖学金，\zh{固定支出}是房租和网费，\zh{变动支出}是吃饭。\zh{其次}，我天天\zh{记账}。上个月我\zh{超支}了，\zh{所以}这个月我\zh{调整}了计划，少买\zh{奢侈品}，多买\zh{必需品}，自己做饭。\zh{如果}你也想\zh{存钱}，\zh{建议}你也\zh{试试}。\zh{希望}我们都能有\zh{储蓄}！
王明
\end{textebox}
\textbf{Analyse linguistique (Pourquoi ça marche) :}
\begin{itemize}
    \item \textbf{Vocabulaire (6/5 requis) :} \zh{预算、收入、固定支出、变动支出、记账、超支、调整、必需品、储蓄} (9 mots). \textbf{Validé.}
    \item \textbf{Grammaire :} \zh{首先/其次} (structure logique). \zh{是...} (identification). \zh{所以} (conséquence). \zh{如果...建议...} (condition + conseil). \zh{也} (aussi) bien placé avant verbe.
    \item \textbf{Caractères :} ~95 caractères (dans fourchette 80-100).
    \item \textbf{Ton :} Amical (\zh{老张、你好、试试}).
\end{itemize}
\cle{Variante orale :} Enregistrer ce message (1 min) en respectant les tons : \zh{预算} (yù\underline{suàn}), \zh{记账} (jì\underline{zhàng}), \zh{超支} (chāo\underline{zhī}), \zh{储蓄} (chǔ\underline{xù}).
\end{exoresolu}

\begin{attention}[Les 5 erreurs qui coûtent des points au Bac — Spécial « Gestion du budget »]
\begin{enumerate}
    \item \textbf{Tons : Confusion \zh{支} (zhī, 1er ton) vs \zh{知} (zhī) / \zh{只} (zhǐ, 3e ton).}
    \cle{Piège :} \zh{支出} (zhīchū - dépenses) se prononce \textbf{1er ton + 1er ton}. \zh{支持} (zhīchí - soutenir) aussi. Mais \zh{只有} (zhǐyǒu - seulement) = 3e ton. \textbf{Exercice :} Lire à haute voix : \zh{收入支出平衡} (shōurù \textbf{zhīchū} pínghéng).
    \item \textbf{Caractères confondus : \zh{账} (compte, clé \zh{贝}) vs \zh{帐} (tente, clé \zh{巾}).}
    \cle{Erreur fréquente :} Écrire \zh{记帐} ou \zh{帐单}. \textbf{Règle :} Argent/Économie $\rightarrow$ \zh{贝} (\zh{账单、账户、账目}). Tissu/Toile $\rightarrow$ \zh{巾} (\zh{帐篷、幕帐}).
    \item \textbf{Ordre des mots : Place du complément de lieu/manière AVANT le verbe.}
    \cle{Calque français :} *\zh{我记账每天} (Je note comptes chaque jour). \textbf{Correct :} \zh{我每天记账} (Sujet + Temps + Verbe). *\zh{他在宿舍做饭} (Correct) vs *\zh{他做饭在宿舍} (Incorrect). \textbf{Règle d'or :} \textbf{Sujet + Temps + Lieu + Manière + Verbe + Objet}.
    \item \textbf{Mots-outils : \zh{了} (aspect accompli/état nouveau) vs \zh{过} (expérience passée).}
    \cle{Contexte budget :} \zh{我超支了} (Cette fois-ci, c'est fait, état actuel « je suis dans le rouge »). \zh{我超支过} (J'ai déjà connu ça dans ma vie, expérience). \textbf{Bac :} « Le mois dernier il a dépassé » $\rightarrow$ \zh{上个月他超支了} (événement passé daté + \zh{了}). « Il a déjà dépassé » $\rightarrow$ \zh{他超支过}.
    \item \textbf{Classificateurs : Oubli ou erreur sur \zh{笔} (traits/écritures/comptes), \zh{项} (articles/projets), \zh{份} (copies/parts).}
    \cle{Collocations budget :} \zh{一笔收入/支出/储蓄/开销/借款} (somme d'argent). \zh{一项预算/计划/开支} (ligne budgétaire). \zh{一份预算表/工资单/报告} (document). \textbf{Erreur :} *\zh{一个预算、一个收入}. \textbf{Correct :} \zh{一份预算、一笔收入}.
\end{enumerate}
\cle{Conseil final :} Relisez votre copie en traçant du doigt : 1. Tons (marques au-dessus voyelles) 2. Clé \zh{贝} sur mots argent 3. Ordre T/L avant V 4. \zh{了}/\zh{过} selon contexte 5. Classifieur \zh{笔/项/份}.
\end{attention}

\newpage

\coursec{Exercices}

\courssub{Je vérifie mes connaissances}

\begin{exercice}[QCM : Vocabulaire et structures du budget]
Chaque question ne comporte qu'une seule réponse correcte. Choisis la bonne proposition.
\begin{enumerate}
    \item Quel mot signifie « salaire / revenu mensuel » ?
    \begin{itemize}
        \item[A.] \zh{房租} (fángzū)
        \item[B.] \zh{工资} (gōngzī)
        \item[C.] \zh{水电费} (shuǐdiànfèi)
        \item[D.] \zh{余额} (yú'é)
    \end{itemize}
    \item Pour dire « économiser de l'argent / mettre de l'argent de côté », on utilise le verbe :
    \begin{itemize}
        \item[A.] \zh{花} (huā)
        \item[B.] \zh{付} (fù)
        \item[C.] \zh{存} (cún)
        \item[D.] \zh{算} (suàn)
    \end{itemize}
    \item La structure correcte pour exprimer « D'abord payer le loyer, ensuite acheter à manger » est :
    \begin{itemize}
        \item[A.] \zh{先买菜，再付房租。}
        \item[B.] \zh{先付房租，再买菜。}
        \item[C.] \zh{付房租先，买菜再。}
        \item[D.] \zh{再付房租，先买菜。}
    \end{itemize}
    \item Le radical (clé) du caractère \zh{费} (fèi, frais / dépense) est :
    \begin{itemize}
        \item[A.] \zh{贝} (bèi, coquillage / argent)
        \item[B.] \zh{氵} (sān diǎn shuǐ, eau)
        \item[C.] \zh{扌} (tí shǒu páng, main)
        \item[D.] \zh{宀} (bǎo gài tóu, toit)
    \end{itemize}
\end{enumerate}
\end{exercice}

Lis le court message de Li Ming à son ami Wang Peng, puis fais l'exercice qui suit.

\begin{textebox}[Message WeChat de Li Ming]
小王，这个月工资到了。我先付了房租和水电费，又买了一些菜。剩下的钱我想存一点，买一张去喀麦隆的机票。\\
Xiǎo Wáng, zhège yuè gōngzī dào le. Wǒ xiān fù le fángzū hé shuǐdiànfèi, yòu mǎi le yìxiē cài. Shèngxià de qián wǒ xiǎng cún yìdiǎn, mǎi yì zhāng qù Kāmàilóng de jīpiào.\\
\emph{« Petit Wang, le salaire de ce mois est arrivé. J'ai d'abord payé le loyer et les factures d'eau et d'électricité, puis j'ai acheté quelques provisions. Avec l'argent restant, je voudrais en mettre un peu de côté pour acheter un billet d'avion pour le Cameroun. »}
\end{textebox}

\begin{exercice}[Vrai ou Faux justifié : Lecture rapide]
Détermine si les affirmations sont vraies (V) ou fausses (F) et justifie ta réponse en citant un extrait du texte (caractères + pinyin).
\begin{enumerate}
    \item Li Ming a reçu son salaire ce mois-ci.
    \item Il a acheté des vêtements avant de payer le loyer.
    \item Il veut mettre de l'argent de côté pour un billet d'avion pour le Cameroun.
    \item Il a dépensé tout son argent pour la nourriture.
\end{enumerate}
\end{exercice}

\begin{exercice}[Vocabulaire : Catégorisation et collocations]
Complète le tableau en classant les mots du glossaire dans la bonne colonne (Revenu / Dépense / Épargne / Action). Puis associe chaque verbe à son objet fréquent (collocation) en traçant une flèche.
\begin{center}
\begin{tabularx}{\linewidth}{|X|X|X|X|}
\hline
\rowcolor{popblueL} \textbf{Revenu (\zh{收入})} & \textbf{Dépense (\zh{支出})} & \textbf{Épargne (\zh{储蓄})} & \textbf{Action (\zh{动作})} \\
\hline
 &  &  &  \\
\hline
\end{tabularx}
\end{center}
\vspace{0.3cm}
\textbf{Mots à classer :} \zh{工资、房租、水电费、存钱、花钱、付账、余额、预算、理财、买菜}\\
\textbf{Collocations à former :}
\begin{itemize}
    \item \zh{付} $\rightarrow$ \zh{房租 / 水电费 / 账单}
    \item \zh{存} $\rightarrow$ \zh{钱 / 钱进银行}
    \item \zh{花} $\rightarrow$ \zh{钱 / 时间}
    \item \zh{做} $\rightarrow$ \zh{预算 / 理财}
\end{itemize}
\end{exercice}

\begin{exercice}[Pinyin et Caractères : Reconnaissance visuelle]
Relie chaque mot en pinyin (avec tons) au caractère chinois correspondant, puis écris la traduction française.
\begin{center}
\begin{tabularx}{\linewidth}{|X|X|X|}
\hline
\rowcolor{popgreenL} \textbf{Pinyin} & \textbf{Caractère} & \textbf{Traduction} \\
\hline
gōngzī &  &  \\
\hline
fángzū &  &  \\
\hline
shuǐdiànfèi &  &  \\
\hline
cún &  &  \\
\hline
yú'é &  &  \\
\hline
\end{tabularx}
\end{center}
\textbf{Caractères à disposition :} \zh{工资、房租、水电费、存、余额}
\end{exercice}

\courssub{Je m'entraîne}

\begin{exercice}[Traduction : Thème (Français $\rightarrow$ Chinois)]
Traduis les phrases suivantes en chinois (caractères + pinyin). Utilise le vocabulaire de la leçon (\zh{先...再...}, \zh{工资、房租、水电费、买菜、存、想}).
\begin{enumerate}
    \item Mon père paie d'abord le loyer, puis il achète la nourriture.
    \item Je ne dépense pas n'importe comment, je veux économiser de l'argent.
    \item Ce mois-ci, les factures d'eau et d'électricité sont très chères.
    \item Après avoir reçu son salaire, elle fait un budget.
\end{enumerate}
\end{exercice}

\begin{exercice}[Transformation de phrases : Négation et Aspect]
Transforme les phrases affirmatives suivantes en phrases négatives (selon l'indication) ou en ajoutant la particule d'aspect appropriée (\zh{了、过、着}). Écris le résultat en caractères et pinyin.
\begin{enumerate}
    \item \zh{我付了房租。} $\rightarrow$ (Négation au passé : \zh{没})
    \item \zh{他存钱。} $\rightarrow$ (Action accomplie : \zh{了})
    \item \zh{我们做预算。} $\rightarrow$ (Expérience vécue : \zh{过})
    \item \zh{水电费很贵。} $\rightarrow$ (Négation : \zh{不})
    \item \zh{他记账。} $\rightarrow$ (Action en cours / duratif : \zh{着})
\end{enumerate}
\end{exercice}

\begin{attention}[Négation au passé : \zh{没} et non \zh{不}]
Pour nier une action \emph{accomplie}, le chinois utilise \zh{没(有) + verbe}, et la particule \zh{了} disparaît :\\
\zh{我付了房租。} $\rightarrow$ \zh{我没付房租。} (et jamais \zh{我没付了房租} ni \zh{我不付了房租}).\\
\zh{不} sert pour le présent, l'habitude ou la volonté : \zh{我不乱花钱。} « Je ne dépense pas n'importe comment. »
\end{attention}

\begin{exercice}[Texte à trous : Mots-outils et vocabulaire clé]
Complète le texte suivant avec les mots manquants choisis dans la liste (chaque mot n'est utilisé qu'une fois). Écris les caractères chinois.
\begin{center}
\zh{我每个月有零花钱。我 \_\_\_\_\_\_ 买文具，\_\_\_\_\_\_ 存一点钱。上个月我 \_\_\_\_\_\_ 了五十元，\_\_\_\_\_\_ 了一本中文书。现在我的余额 \_\_\_\_\_\_ 多，我得 \_\_\_\_\_\_ 着花钱。}\\
\emph{(Wǒ měige yuè yǒu línghuāqián. Wǒ \_\_\_\_ mǎi wénjù, \_\_\_\_ cún yìdiǎn qián. Shàngge yuè wǒ \_\_\_\_ le wǔshí yuán, \_\_\_\_ le yì běn Zhōngwén shū. Xiànzài wǒ de yú'é \_\_\_\_ duō, wǒ děi \_\_\_\_ zhe huā qián.)}
\end{center}
\textbf{Liste :} \zh{先、再、花、买、不、省}
\end{exercice}

\begin{exercice}[Appariement et Remise en ordre]
\begin{enumerate}
    \item \textbf{Appariement :} Classe chaque terme dans la bonne catégorie (Revenu / Dépense / Épargne).
    \begin{center}
    \begin{tabular}{lll}
    \zh{工资} & \zh{房租} & \zh{存款} \\
    \zh{奖金} & \zh{水电费} & \zh{理财产品} \\
    \zh{兼职收入} & \zh{买菜钱} & \zh{零花钱}
    \end{tabular}
    \end{center}
    \item \textbf{Remise en ordre :} Remets les mots dans l'ordre pour former une phrase correcte (caractères + pinyin).
    \begin{itemize}
        \item \zh{钱 / 我 / 存 / 每月 / 想 / 一点}
        \item \zh{很 / 房租 / 今年 / 贵}
        \item \zh{预算 / 做 / 我们 / 应该 / 一个}
    \end{itemize}
\end{enumerate}
\end{exercice}

\courssub{Je me prépare au Bac}

Lis le texte suivant, puis réponds aux questions de l'exercice de compréhension.

\begin{textebox}[Le budget mensuel de Xiao Chen]
陈晨是北京的一名大学生。每个月父母给他两千元生活费。他住在学校宿舍，不用付房租。但是伙食费很贵，食堂一顿饭大概二三十元。他每天吃三顿，一个月伙食费要一千五百元左右。陈晨喜欢打篮球，买球鞋花了五百元。剩下的钱他存进银行卡里，准备暑假去喀麦隆旅游。陈晨说：「记账很重要，钱要花在刀刃上。」\\
Chénchén shì Běijīng de yì míng dàxuéshēng. Měige yuè fùmǔ gěi tā liǎngqiān yuán shēnghuófèi. Tā zhù zài xuéxiào sùshè, bùyòng fù fángzū. Dànshì huǒshífèi hěn guì, shítáng yí dùn fàn dàgài èr-sānshí yuán. Tā měitiān chī sān dùn, yí ge yuè huǒshífèi yào yīqiān wǔbǎi yuán zuǒyòu. Chénchén xǐhuan dǎ lánqiú, mǎi qiúxié huā le wǔbǎi yuán. Shèngxià de qián tā cún jìn yínhángkǎ lǐ, zhǔnbèi shǔjià qù Kāmàilóng lǚyóu. Chénchén shuō : « Jìzhàng hěn zhòngyào, qián yào huā zài dāorèn shàng. »\\
\emph{« Chenchen est un étudiant de Pékin. Chaque mois, ses parents lui donnent deux mille yuans pour ses frais de vie. Il loge dans la résidence universitaire, il n'a pas de loyer à payer. Mais les frais de nourriture sont élevés : un repas à la cantine coûte environ vingt à trente yuans. Il mange trois repas par jour, soit environ mille cinq cents yuans de nourriture par mois. Chenchen aime jouer au basket ; il a dépensé cinq cents yuans pour des chaussures de sport. L'argent restant, il le dépose sur sa carte bancaire, en préparation d'un voyage au Cameroun pendant les vacances d'été. Chenchen dit : "Tenir ses comptes est très important, l'argent doit être dépensé là où c'est utile." »}
\end{textebox}

\begin{exercice}[Compréhension écrite type Baccalauréat : Le budget de Xiao Chen]
Réponds aux questions en français (Q1--Q3) et en chinois (Q4--Q5).
\begin{enumerate}
    \item (V/F + justification) Xiao Chen paie un loyer tous les mois.
    \item (Question ouverte) Combien Xiao Chen dépense-t-il environ pour la nourriture par mois ? D'où vient cet argent ?
    \item (Question ouverte) Que compte-t-il faire de ses économies pendant les vacances d'été ?
    \item \textbf{En chinois :} \zh{陈晨为什么觉得记账很重要？}
    \item \textbf{En chinois :} Réécris avec \zh{先...再...} : \zh{他买球鞋，存剩下的钱。}
\end{enumerate}
\end{exercice}

\begin{savaistu}[« 钱要花在刀刃上 »]
L'expression \zh{钱要花在刀刃上} (qián yào huā zài dāorèn shàng), littéralement « l'argent doit être dépensé sur le tranchant de la lame », signifie qu'il faut dépenser là où c'est le plus utile, sans gaspillage. Au Cameroun comme en Chine, les familles apprennent aux jeunes à \zh{记账} (jìzhàng, tenir ses comptes) : cahier de dépenses à Yaoundé, application mobile à Pékin, la prudence budgétaire est une valeur partagée.
\end{savaistu}

\begin{exercice}[Traduction : Thème et Version]
\begin{enumerate}
    \item \textbf{Thème (Français $\rightarrow$ Chinois) :} Traduis en chinois (caractères + pinyin). \textit{Attendu : 2 phrases, environ 50 caractères au total.}
    \begin{itemize}
        \item[a.] « Chaque mois, je reçois mon salaire, je paie d'abord le loyer et l'électricité, ensuite je fais les courses. »
        \item[b.] « Cette année, la vie est chère, donc je dois économiser pour voyager au Cameroun. »
    \end{itemize}
    \item \textbf{Version (Chinois $\rightarrow$ Français) :} Traduis le texte suivant en français courant.
    \begin{center}
    \zh{老张工资不高，但他不乱花钱。菜市场的菜比较便宜，他常去那里买。他每个月都能存下五百元。他说：「省钱不是小气，是为了将来不愁钱。」}\\
    \emph{(Lǎo Zhāng gōngzī bù gāo, dàn tā bù luàn huā qián. Càishìchǎng de cài bǐjiào piányi, tā cháng qù nàlǐ mǎi. Tā měige yuè dōu néng cún xia wǔbǎi yuán. Tā shuō : « Shěng qián bù shì xiǎoqì, shì wèile jiānglái bù chóu qián. »)}
    \end{center}
\end{enumerate}
\end{exercice}

\begin{exercice}[Expression écrite : Mon budget idéal]
\textbf{Sujet :} « \emph{Présente ton budget mensuel idéal (revenus, dépenses fixes, loisirs, épargne) et explique tes choix.} »
\textbf{Consignes :}
\begin{itemize}
    \item Rédige un paragraphe cohérent en chinois (caractères simplifiés).
    \item Longueur : \textbf{80 à 100 caractères chinois} (ponctuation comprise).
    \item Utilise \textbf{au moins 6 mots} du glossaire de la leçon (ex. : \zh{工资、房租、水电费、伙食费、存、省、预算、余额、记账}).
    \item Utilise les structures : \zh{先...再...}, \zh{因为...所以...}, \zh{比较 / 很 + adjectif}.
    \item Soigne l'ordre des traits pour les caractères manuscrits (si copie papier) ou la saisie correcte (si numérique).
\end{itemize}
\begin{center}
\textit{Mots-clés pour t'aider :} \\
\zh{收入来源：兼职 / 零花钱 / 奖学金} \quad \zh{支出：房租、吃饭、交通、网费} \quad \zh{结余：存银行、买书 / 旅游}
\end{center}
\end{exercice}

\begin{methode}[Rédiger un paragraphe « budget » en 4 étapes]
\begin{enumerate}
    \item \textbf{Annoncer le revenu} : \zh{我每个月有……元。} (source : \zh{零花钱 / 兼职收入}).
    \item \textbf{Ordonner les dépenses} avec \zh{先...再...} : \zh{我先付……，再买……。}
    \item \textbf{Justifier un choix} avec \zh{因为...所以...} : \zh{因为……很贵，所以我要省钱。}
    \item \textbf{Conclure par l'épargne} : \zh{剩下的钱我存进银行。} + une phrase d'opinion : \zh{记账很重要。}
\end{enumerate}
\emph{Vérifie :} un verbe par phrase, l'ordre Sujet--Temps--Verbe--Complément, et les tons du pinyin si tu recopies la transcription.
\end{methode}

\begin{exemplebox}[Modèle de production (environ 90 caractères)]
\zh{我每个月有一百元零花钱。我先付交通费和网费，再买菜和文具。因为雅温得的生活比较贵，所以我不乱花钱。剩下的钱我存进银行，准备假期去旅游。我觉得记账很重要，钱要花在刀刃上。}\\
\emph{(Wǒ měige yuè yǒu yìbǎi yuán línghuāqián. Wǒ xiān fù jiāotōngfèi hé wǎngfèi, zài mǎi cài hé wénjù. Yīnwèi Yǎwēndé de shēnghuó bǐjiào guì, suǒyǐ wǒ bù luàn huā qián. Shèngxià de qián wǒ cún jìn yínháng, zhǔnbèi jiàqī qù lǚyóu. Wǒ juéde jìzhàng hěn zhòngyào, qián yào huā zài dāorèn shàng.)}\\
« Chaque mois, j'ai cent yuans d'argent de poche. Je paie d'abord le transport et Internet, ensuite j'achète nourriture et fournitures. Comme la vie à Yaoundé est assez chère, je ne dépense pas n'importe comment. L'argent restant, je le dépose à la banque pour voyager pendant les vacances. Je trouve que tenir ses comptes est très important : l'argent doit être dépensé utilement. »
\end{exemplebox}

\newpage

\coursec{Corrigés des exercices}

\courssub{Corrigés détaillés}

\begin{corrige}[Exercice 1 — QCM : Vocabulaire et structures du budget]
\begin{enumerate}
    \item \textbf{B.} \zh{工资} (gōngzī) = salaire. \zh{房租} = loyer ; \zh{水电费} = factures d'eau et d'électricité ; \zh{余额} = solde restant.
    \item \textbf{C.} \zh{存} (cún) = déposer, économiser ; la collocation est \zh{存钱} (cún qián, mettre de l'argent de côté). \zh{花} = dépenser ; \zh{付} = payer ; \zh{算} = calculer.
    \item \textbf{B.} Structure \zh{先 + verbe 1，再 + verbe 2} : \zh{先付房租，再买菜。} (Xiān fù fángzū, zài mǎi cài.) « D'abord payer le loyer, ensuite acheter à manger. » Les propositions C et D violent l'ordre des mots : \zh{先} et \zh{再} se placent \emph{avant} le verbe, jamais après.
    \item \textbf{A.} La clé de \zh{费} est \zh{贝} (bèi), le coquillage, ancienne monnaie en Chine. On la retrouve dans \zh{贵} (guì, cher), \zh{账} (zhàng, compte), \zh{购} (gòu, acheter), \zh{赚} (zhuàn, gagner).
\end{enumerate}
\emph{Point de vigilance :} ne pas confondre \zh{存} (épargner) et \zh{花} (dépenser) : même champ lexical, sens opposés. Retenir aussi que \zh{再} (zài) introduit l'action \emph{suivante} (« ensuite »), alors que \zh{又} (yòu) marque l'ajout ou la répétition d'une action déjà faite.
\end{corrige}

\begin{corrige}[Exercice 2 — Vrai ou Faux justifié]
\begin{enumerate}
    \item \textbf{V.} \zh{这个月工资到了。} (Zhège yuè gōngzī dào le.) « Le salaire de ce mois est arrivé. »
    \item \textbf{F.} Le texte dit \zh{我先付了房租和水电费，又买了一些菜。} (Wǒ xiān fù le fángzū hé shuǐdiànfèi, yòu mǎi le yìxiē cài.) : il a payé le loyer \emph{d'abord}, et il a acheté des \emph{provisions} (\zh{菜}), pas des vêtements (\zh{衣服}).
    \item \textbf{V.} \zh{剩下的钱我想存一点，买一张去喀麦隆的机票。} (Shèngxià de qián wǒ xiǎng cún yìdiǎn, mǎi yì zhāng qù Kāmàilóng de jīpiào.) « Avec l'argent restant, je veux en mettre un peu de côté pour acheter un billet d'avion pour le Cameroun. »
    \item \textbf{F.} Il lui \emph{reste} de l'argent (\zh{剩下的钱}), donc il n'a pas tout dépensé ; il a seulement acheté \zh{一些菜} (quelques provisions).
\end{enumerate}
\emph{Point de vigilance :} \zh{又} (yòu) marque ici l'enchaînement d'actions accomplies (« et puis »), à ne pas traduire par « encore » au sens de répétition future. Noter aussi le classificateur \zh{张} (zhāng) devant \zh{机票} : il sert pour les objets plats (billets, feuilles, tables).
\end{corrige}

\begin{corrige}[Exercice 3 — Vocabulaire : Catégorisation et collocations]
\textbf{Classement :}
\begin{itemize}
    \item \textbf{Revenu (\zh{收入}) :} \zh{工资}.
    \item \textbf{Dépense (\zh{支出}) :} \zh{房租、水电费、花钱、付账、买菜}.
    \item \textbf{Épargne (\zh{储蓄}) :} \zh{存钱、余额、理财}.
    \item \textbf{Action / gestion (\zh{动作}) :} \zh{预算} (on dit \zh{做预算}, faire un budget).
\end{itemize}
\textbf{Collocations attendues :}
\begin{itemize}
    \item \zh{付房租 / 付水电费 / 付账单} (payer le loyer / les factures / la note) ;
    \item \zh{存钱、把钱存进银行} (mettre de l'argent de côté, déposer à la banque) ;
    \item \zh{花钱、花时间} (dépenser de l'argent, passer du temps) ;
    \item \zh{做预算、理财} (faire un budget, gérer ses finances).
\end{itemize}
\emph{Remarque :} \zh{余额} (yú'é, solde) est le \emph{résultat} de l'épargne ; \zh{理财} (lǐcái) signifie « gérer ses finances ». La construction \zh{把 + objet + verbe + complément} (\zh{把钱存进银行}) met l'objet en avant : structure caractéristique du chinois, à réemployer dans les productions.
\end{corrige}

\begin{corrige}[Exercice 4 — Pinyin et Caractères : Reconnaissance visuelle]
\begin{center}
\begin{tabularx}{\linewidth}{|X|X|X|}
\hline
\rowcolor{popgreenL} \textbf{Pinyin} & \textbf{Caractère} & \textbf{Traduction} \\
\hline
gōngzī & \zh{工资} & salaire \\
\hline
fángzū & \zh{房租} & loyer \\
\hline
shuǐdiànfèi & \zh{水电费} & factures d'eau et d'électricité \\
\hline
cún & \zh{存} & déposer, économiser \\
\hline
yú'é & \zh{余额} & solde restant \\
\hline
\end{tabularx}
\end{center}
\emph{Point de vigilance :} dans \zh{余额} (yú'é), l'apostrophe du pinyin sépare les deux syllabes pour éviter la lecture *yu'é ; attention aussi au troisième ton de \zh{水} (shuǐ) dans \zh{水电费}. Rappel d'écriture : \zh{存} s'écrit de gauche à droite (la clé \zh{子}, enfant, à droite) ; \zh{费} s'écrit de haut en bas, avec la clé \zh{贝} (coquillage, monnaie) en bas.
\end{corrige}

\begin{corrige}[Exercice 5 — Traduction : Thème (Français $\rightarrow$ Chinois)]
\begin{enumerate}
    \item \zh{我爸爸先付房租，再买菜。} (Wǒ bàba xiān fù fángzū, zài mǎi cài.)\\
    \emph{Justification :} \zh{先} et \zh{再} précèdent immédiatement le verbe ; « la nourriture » se rend ici par \zh{菜} (provisions, plats).
    \item \zh{我不乱花钱，我想存钱。} (Wǒ bù luàn huā qián, wǒ xiǎng cún qián.)\\
    \emph{Justification :} \zh{乱} (luàn, désordonné) placé avant le verbe exprime « n'importe comment » ; \zh{想 + verbe} = vouloir faire.
    \item \zh{这个月水电费很贵。} (Zhège yuè shuǐdiànfèi hěn guì.)\\
    \emph{Justification :} le complément de temps \zh{这个月} se place après le sujet (ou en tête) ; pas de verbe « être » devant l'adjectif : \zh{很贵}.
    \item \zh{她收到工资以后做预算。} (Tā shōudào gōngzī yǐhòu zuò yùsuàn.)\\
    \emph{Justification :} « après avoir » = \zh{……以后} placé \emph{après} le groupe verbal, particularité de l'ordre chinois.
\end{enumerate}
\end{corrige}

\begin{corrige}[Exercice 6 — Transformation de phrases : Négation et Aspect]
\begin{enumerate}
    \item \zh{我没付房租。} (Wǒ méi fù fángzū.) — Avec \zh{没(有)}, la particule \zh{了} disparaît obligatoirement.
    \item \zh{他存了钱。} (Tā cún le qián.) — \zh{了} après le verbe marque l'action accomplie.
    \item \zh{我们做过预算。} (Wǒmen zuò guo yùsuàn.) — \zh{过} marque l'expérience vécue : « nous avons déjà fait un budget (au moins une fois) ».
    \item \zh{水电费不贵。} (Shuǐdiànfèi bù guì.) — La négation d'un adjectif se fait avec \zh{不}.
    \item \zh{他记着账。} (Tā jì zhe zhàng.) — \zh{着} (zhe) marque la durée ou l'état continu : « il tient ses comptes (en ce moment, régulièrement) ».
\end{enumerate}
\emph{Point de vigilance :} on ne dit jamais \zh{我没付了房租} : \zh{没} et \zh{了} sont incompatibles dans la même proposition. De même, la négation de \zh{过} se fait avec \zh{没(有)} : \zh{我们没做过预算。} (nous n'avons jamais fait de budget).
\end{corrige}

\begin{corrige}[Exercice 7 — Texte à trous]
\zh{我每个月有零花钱。我先买文具，再存一点钱。上个月我花了五十元，买了一本中文书。现在我的余额不多，我得省着花钱。}\\
\emph{(Wǒ měige yuè yǒu línghuāqián. Wǒ xiān mǎi wénjù, zài cún yìdiǎn qián. Shàngge yuè wǒ huā le wǔshí yuán, mǎi le yì běn Zhōngwén shū. Xiànzài wǒ de yú'é bù duō, wǒ děi shěng zhe huā qián.)}\\
« Chaque mois, j'ai de l'argent de poche. J'achète d'abord des fournitures, ensuite je mets un peu d'argent de côté. Le mois dernier, j'ai dépensé cinquante yuans et acheté un livre de chinois. Maintenant mon solde n'est pas élevé, je dois dépenser avec économie. »\\
\emph{Justifications :} \zh{先...再...} ordonne les deux actions ; \zh{花了} (dépenser, action accomplie au passé) ; \zh{买了} (acheter, idem) ; \zh{不多} (pas beaucoup, négation d'adjectif) ; \zh{省着花钱} = dépenser en économisant, \zh{着} marquant le mode continu de l'action. Noter le classificateur \zh{本} (běn) devant \zh{书} (livre).
\end{corrige}

\begin{corrige}[Exercice 8 — Appariement et Remise en ordre]
\begin{enumerate}
    \item \textbf{Revenu :} \zh{工资} (salaire), \zh{奖金} (prime), \zh{兼职收入} (revenu d'un emploi à temps partiel), \zh{零花钱} (argent de poche — c'est un revenu pour un élève).\\
    \textbf{Dépense :} \zh{房租} (loyer), \zh{水电费} (factures), \zh{买菜钱} (argent des courses).\\
    \textbf{Épargne :} \zh{存款} (dépôts, économies), \zh{理财产品} (produits de placement).
    \item \begin{itemize}
        \item \zh{我每月想存一点钱。} (Wǒ měiyuè xiǎng cún yìdiǎn qián.) « Je veux mettre un peu d'argent de côté chaque mois. » — Ordre : Sujet + temps + \zh{想} + verbe.
        \item \zh{今年房租很贵。} (Jīnnián fángzū hěn guì.) « Cette année, le loyer est très cher. » — Temps en tête, adjectif prédicat sans verbe « être ».
        \item \zh{我们应该做一个预算。} (Wǒmen yīnggāi zuò yí ge yùsuàn.) « Nous devrions faire un budget. » — \zh{应该} (devoir) avant le verbe ; classificateur \zh{个} devant \zh{预算}.
    \end{itemize}
\end{enumerate}
\end{corrige}

\begin{corrige}[Exercice 9 — Compréhension écrite type Baccalauréat]
\textbf{Barème indicatif (sur 10 points) :} Q1 : 1,5 pt (0,5 pour F + 1 pour la justification citée) ; Q2 : 2 pts (1 pt par élément) ; Q3 : 1,5 pt ; Q4 : 2,5 pts (chinois correct exigé) ; Q5 : 2,5 pts (structure \zh{先...再...} correcte).
\begin{enumerate}
    \item \textbf{F.} \zh{他住在学校宿舍，不用付房租。} (Tā zhù zài xuéxiào sùshè, bùyòng fù fángzū.) « Il loge en résidence universitaire, il n'a pas besoin de payer de loyer. »
    \item Il dépense environ \textbf{1 500 yuans} par mois pour la nourriture (\zh{一个月伙食费要一千五百元左右}). Cet argent vient des \textbf{2 000 yuans} que ses parents lui donnent chaque mois (\zh{父母给他两千元生活费}).
    \item Il dépose l'argent restant sur sa carte bancaire pour \textbf{voyager au Cameroun} pendant les vacances d'été (\zh{准备暑假去喀麦隆旅游}).
    \item Réponse modèle : \zh{因为钱要花在刀刃上，记账能帮他管好钱。} (Yīnwèi qián yào huā zài dāorèn shàng, jìzhàng néng bāng tā guǎn hǎo qián.) « Parce que l'argent doit être dépensé à bon escient ; tenir ses comptes l'aide à bien gérer son argent. »
    \item \zh{他先买球鞋，再存剩下的钱。} (Tā xiān mǎi qiúxié, zài cún shèngxià de qián.) « Il achète d'abord les chaussures de sport, ensuite il met de côté l'argent restant. »
\end{enumerate}
\emph{Points de vigilance :} pour Q4, exiger \zh{因为} en tête de réponse ; pour Q5, vérifier que \zh{先} et \zh{再} précèdent leurs verbes et que l'ordre logique du texte est respecté (achat puis épargne). \zh{左右} (zuǒyòu), placé \emph{après} la quantité, signifie « environ » : \zh{一千五百元左右} = environ 1 500 yuans.
\end{corrige}

\begin{corrige}[Exercice 10 — Traduction : Thème et Version]
\textbf{Barème indicatif (sur 10 points) :} Thème : 6 pts (3 pts par phrase : vocabulaire 1 pt, structure 1 pt, exactitude des caractères 1 pt) ; Version : 4 pts (fidélité 2 pts, qualité du français 2 pts).

\textbf{Thème :}
\begin{itemize}
    \item[a.] \zh{每个月我收到工资，先付房租和电费，再去买菜。} (Měige yuè wǒ shōudào gōngzī, xiān fù fángzū hé diànfèi, zài qù mǎi cài.)
    \item[b.] \zh{今年物价很高，所以我要省钱，存钱去喀麦隆旅游。} (Jīnnián wùjià hěn gāo, suǒyǐ wǒ yào shěng qián, cún qián qù Kāmàilóng lǚyóu.)
\end{itemize}
\emph{Justifications :} en (a), le temps \zh{每个月} ouvre la phrase et \zh{先...再...} enchaîne les deux dépenses ; en (b), « la vie est chère » se rend naturellement par \zh{物价很高} (les prix sont élevés) plutôt que par une traduction mot à mot ; « donc » = \zh{所以}, et « économiser » se rend par \zh{省钱} (dépenser moins au quotidien), tandis que \zh{存钱} exprime la mise de côté en vue du voyage.

\textbf{Version (proposition de traduction) :}\\
« Le salaire du vieux Zhang n'est pas élevé, mais il ne dépense pas n'importe comment. Les légumes au marché sont relativement bon marché, il s'y rend souvent pour faire ses achats. Chaque mois, il arrive à mettre de côté cinq cents yuans. Il dit : "Économiser, ce n'est pas être avare, c'est pour ne pas avoir de soucis d'argent à l'avenir." »

\emph{Points de vigilance :} \zh{存下} = réussir à mettre de côté (complément de résultat) ; \zh{小气} (xiǎoqì) = avare, pingre ; \zh{不愁钱} = ne pas se soucier de l'argent. Éviter la traduction mot à mot « économiser l'argent n'est pas petit esprit ».
\end{corrige}

\begin{corrige}[Exercice 11 — Expression écrite : Mon budget idéal]
\textbf{Barème indicatif (sur 10 points) :} respect de la longueur (80--100 caractères) : 1 pt ; au moins 6 mots du glossaire : 2 pts ; structures imposées (\zh{先...再...}, \zh{因为...所以...}, adjectif avec \zh{很/比较}) : 3 pts ; correction grammaticale et ordre des mots : 2 pts ; cohérence et richesse du propos : 2 pts.

\textbf{Proposition de rédaction modèle (environ 90 caractères) :}\\
\zh{我每个月有一百元零花钱。我先付交通费和网费，再买菜和文具。因为雅温得的生活比较贵，所以我不乱花钱。剩下的钱我存进银行，准备假期去旅游。我觉得记账很重要，钱要花在刀刃上。}\\
\emph{(Wǒ měige yuè yǒu yìbǎi yuán línghuāqián. Wǒ xiān fù jiāotōngfèi hé wǎngfèi, zài mǎi cài hé wénjù. Yīnwèi Yǎwēndé de shēnghuó bǐjiào guì, suǒyǐ wǒ bù luàn huā qián. Shèngxià de qián wǒ cún jìn yínháng, zhǔnbèi jiàqī qù lǚyóu. Wǒ juéde jìzhàng hěn zhòngyào, qián yào huā zài dāorèn shàng.)}\\
« Chaque mois, j'ai cent yuans d'argent de poche. Je paie d'abord le transport et Internet, ensuite j'achète nourriture et fournitures. Comme la vie à Yaoundé est assez chère, je ne dépense pas n'importe comment. L'argent restant, je le dépose à la banque pour voyager pendant les vacances. Je trouve que tenir ses comptes est très important : l'argent doit être dépensé utilement. »

\textbf{Points de vigilance :}
\begin{itemize}
    \item Un seul verbe principal par phrase ; ne pas juxtaposer deux verbes sans mot-outil.
    \item Respecter l'ordre Sujet + Temps + Verbe + Complément : \zh{我每个月有……} ; la variante \zh{每个月我有……} (temps en tête) est possible, mais il faut rester cohérent dans tout le texte.
    \item \zh{因为...所以...} : en chinois, les deux moitiés s'emploient ensemble, contrairement au français (« parce que... donc » est fautif en français, correct en chinois).
    \item Vérifier les tons en recopiant le pinyin : \zh{línghuāqián} (2-1-2), \zh{yùsuàn} (4-4), \zh{jìzhàng} (4-4).
    \item Compter les caractères : la ponctuation chinoise pleine chasse （，。） compte dans le total.
\end{itemize}
\end{corrige}

\newpage

\coursec{Fiche bilan}

\begin{popfigure}
\begin{tikzpicture}[
    mindmap,
    level 1 concept/.append style={level distance=3.5cm, sibling angle=45},
    level 2 concept/.append style={level distance=2.5cm},
    every node/.style={concept, execute at begin node=\hskip0pt},
    concept color=popblue,
    text=white,
    font=\small\bfseries,
    grow cyclic
]
\node[concept color=popdark, minimum size=2.2cm, align=center]{Gestion du budget\\预算管理 (Yùsuàn Guǎnlǐ)}
    child[concept color=popblue] { node[minimum size=1.8cm, align=center]{Vocabulaire\\核心词汇} 
        child { node[concept color=popblueL, font=\tiny] {收入 / 支出} }
        child { node[concept color=popblueL, font=\tiny] {储蓄 / 投资} }
        child { node[concept color=popblueL, font=\tiny] {贷款 / 利息} }
        child { node[concept color=popblueL, font=\tiny] {记账 / 平衡} }
    }
    child[concept color=poporange] { node[minimum size=1.8cm, align=center]{Grammaire\\重点语法} 
        child { node[concept color=poporangeL, font=\tiny] {比较句 (比)} }
        child { node[concept color=poporangeL, font=\tiny] {越来越 + Adj.} }
        child { node[concept color=poporangeL, font=\tiny] {只要...就...} }
        child { node[concept color=poporangeL, font=\tiny] {被字句 (Passif)} }
    }
    child[concept color=popgreen] { node[minimum size=1.8cm, align=center]{Écriture\\关键汉字} 
        child { node[concept color=popgreenL, font=\tiny] {Radical 贝 (argent)} }
        child { node[concept color=popgreenL, font=\tiny] {Ordre des traits} }
        child { node[concept color=popgreenL, font=\tiny] {结构 : 上下/左右} }
    }
    child[concept color=poppurple] { node[minimum size=1.8cm, align=center]{Compréhension\\阅读理解} 
        child { node[concept color=poppurpleL, font=\tiny] {Texte : Budget étudiant} }
        child { node[concept color=poppurpleL, font=\tiny] {Questions QCM / Ouvertes} }
        child { node[concept color=poppurpleL, font=\tiny] {Recherche d'infos} }
    }
    child[concept color=poppink] { node[minimum size=1.8cm, align=center]{Production\\表达应用} 
        child { node[concept color=poppinkL, font=\tiny] {Décrire son budget} }
        child { node[concept color=poppinkL, font=\tiny] {Conseils 理财} }
        child { node[concept color=poppinkL, font=\tiny] {Dialogue banque} }
    }
    child[concept color=popgold] { node[minimum size=1.8cm, align=center]{Culture\\中喀对比} 
        child { node[concept color=popgoldL, font=\tiny] {Mobile payment (CN)} }
        child { node[concept color=popgoldL, font=\tiny] {Mobile money (CM)} }
        child { node[concept color=popgoldL, font=\tiny] {Habitudes épargne} }
    };
\end{tikzpicture}
\legende{Carte mentale récapitulative : Leçon 25 — Gestion du budget}
\end{popfigure}

\begin{bilan}
\textbf{Lexique clé (HSK 3-4) — Thème : Gestion du budget / 预算管理}

\begin{center}
\begin{tabularx}{\linewidth}{|X|X|X|X|}
\hline
\rowcolor{popblueL} \textbf{Caractères} & \textbf{Pinyin} & \textbf{Nature} & \textbf{Traduction} \\ \hline
预算 & yùsuàn & n. & budget \\ \hline
收入 & shōurù & n./v. & revenus / percevoir des revenus \\ \hline
支出 & zhīchū & n./v. & dépenses / dépenser \\ \hline
结余 & jiéyú & n. & solde, reste (budget) \\ \hline
储蓄 & chǔxù & n./v. & épargne / épargner \\ \hline
投资 & tóuzī & n./v. & investissement / investir \\ \hline
贷款 & dàikuǎn & n. & prêt, crédit \\ \hline
利息 & lìxī & n. & intérêts (bancaires) \\ \hline
记账 & jìzhàng & v. & tenir les comptes, noter les dépenses \\ \hline
理财 & lǐcái & v./n. & gestion financière / finance personnelle \\ \hline
量入为出 & liàngrùwéichū & id. & adapter ses dépenses à ses revenus (vivre selon ses moyens) \\ \hline
月光族 & yuèguāngzú & n. & « clan du mois clair » : personnes qui dépensent tout leur salaire chaque mois \\ \hline
剁手 & duòshǒu & v. (argot) & « se couper la main » : arrêter d'acheter en ligne (acheteur compulsif) \\ \hline
\end{tabularx}
\end{center}

\textbf{Règles de grammaire essentielles (à connaître par cœur)}

\begin{center}
\begin{tabularx}{\linewidth}{|X|X|X|}
\hline
\rowcolor{poporangeL} \textbf{Structure} & \textbf{Exemple (Caractères + Pinyin)} & \textbf{Traduction} \\ \hline
\textbf{比较句 (Comparatif)} 
A + 比 + B + Adj. (pas de 很) & 我的收入比支出高。 
(Wǒ de shōurù bǐ zhīchū gāo.) & Mes revenus sont plus élevés que mes dépenses. \\ \hline
\textbf{越来越 + Adj./Verbe mental} 
(Évolution croissante) & 生活费越来越贵了。 
(Shēnghuófèi yuèláiyuè guì le.) & Le coût de la vie est de plus en plus cher. \\ \hline
\textbf{除了...以外/除了...还...} 
(Inclusion / Exclusion) & 除了房租，我还要交水电费。 
(Chúle fángzū, wǒ hái yào jiāo shuǐdiànfèi.) & En plus du loyer, je dois payer l'eau et l'électricité. \\ \hline
\textbf{只要...就...} 
(Condition suffisante) & 只要你记账，就能省钱。 
(Zhǐyào nǐ jìzhàng, jiù néng shěngqián.) & Tant que tu tiens tes comptes, tu pourras économiser. \\ \hline
\textbf{被字句 (Passif / Subi)} 
Sujet + 被 + Agent + Verbe + Complément & 我的钱包被偷了。 
(Wǒ de qiánbāo bèi tōu le.) & Mon portefeuille a été volé. \\ \hline
\textbf{把字句 (Disposition)} 
Sujet + 把 + Objet + Verbe + Résultat & 请把账单算清楚。 
(Qǐng bǎ zhàngdān suàn qīngchu.) & Calculez la facture clairement. \\ \hline
\end{tabularx}
\end{center}

\textbf{Méthodes express}
\begin{itemize}
    \item \textbf{Lecture rapide :} Repérer les mots-clés (chiffres, \cle{收入}, \cle{支出}, \cle{结余}) $\rightarrow$ Identifier le problème (déficit ? excédent ?) $\rightarrow$ Chercher la solution proposée (\cle{记账}, \cle{理财}).
    \item \textbf{Rédaction budget :} 1. Lister revenus (固定/变动). 2. Lister dépenses (必要/非必要). 3. Calculer \cle{结余} (收入 - 支出). 4. Proposer ajustement (\cle{量入为出}).
    \item \textbf{Caractères « Argent » :} Repérer le radical \cle{贝} (bèi) : 费 、 贷 、 账 、 赚 、 赔 、 货 。 Ordre : haut $\rightarrow$ bas, gauche $\rightarrow$ droite ; horizontal avant vertical.
\end{itemize}
\end{bilan}

\begin{aretenir}[Checklist avant le Bac — Leçon 25]
$\square$ Je sais écrire et dicter les 12 mots du tableau vocabulaire (caractères + pinyin + tons). \\
$\square$ Je reconnais le radical \textbf{贝} (bèi) et j'en déduis le sens « argent/valeur » pour 5 nouveaux caractères. \\
$\square$ Je maîtrise la structure comparative \textbf{A 比 B Adj.} (sans \textbf{很}) et sa forme négative \textbf{A 没有 B Adj.}. \\
$\square$ J'utilise correctement \textbf{越来越} pour décrire une tendance (prix, épargne, dettes). \\
$\square$ Je distingue \textbf{除了...以外} (exclusion) de \textbf{除了...还/也...} (inclusion/addition). \\
$\square$ Je sais transformer une phrase active en phrase passive avec \textbf{被} (ex: facture payée, carte volée). \\
$\square$ Je sais utiliser \textbf{把} pour indiquer la manipulation d'un objet précis (payer la facture, calculer le solde). \\
$\square$ Je comprends un texte authentique sur le budget d'un étudiant (Yaoundé / Pékin) et réponds aux questions (V/F, Ouvertes). \\
$\square$ Je produis un court paragraphe (80-100 caractères) pour présenter mon budget mensuel et un conseil \textbf{理财}. \\
$\square$ Je connais les différences culturelles : \textbf{移动支付} (Alipay/WeChat Pay) en Chine vs \textbf{Mobile Money} (Orange Money/MTN MoMo) au Cameroun. \\
$\square$ Je connais l'idiome \textbf{量入为出} et le terme argotique \textbf{月光族} / \textbf{剁手族}.
\end{aretenir}

\findecours

\end{document}
